\bookStart{Spae of the Wallow}[Vǫlu spǫ́]
\setBookCode{Voluspa}

\begin{flushright}%
\textbf{Dating} \parencite{Sapp2022}: C10th (0.865)

\textbf{Meter:} \Fornyrdislag%
\end{flushright}

\section{Introduction}

The \textbf{Spae of the Wallow} (ON \emph{Vǫlu spǫ́}, abbrev.~Voluspa) is undoubtedly the most important Norse mythological text surviving from Heathen times.  It is a \inx[C]{spae} (\emph{spǫ́} ‘prophecy’) told by a \inx[C]{wallow} (\emph{vǫlva} ‘seeress, sibyl, prophetess’) who has been summoned by \inx[P]{Weden}, the chief of the Gods, in order to relate the mythic history of the World.

\subsection{Origin and parallels}

As a repository of mythic lore \Voluspa\ resembles \textlink{Vafthrudnismal}, \textlink{Grimnismal}, \textlink{Baldrsdraumar}, and \textlink{Alvissmal}, but it differs from most of the aforementioned works in two key ways: \textlink{Voluspa} is a tightly structured spae rather than a motley collection of scattered mythological lore, and a monologue rather than a dialogue or riddle contest.  The unique importance of this poem lies in the fact that it offers a complete chronological overview of the Norse mythic timeline of the World from its creation to its coming destruction and rebirth.

The understanding of this timeline is, unfortunately for the study of Norse mythology, troubled by several factors.  The first is that the wallow continually speaks in the most obscure terms and relates the myths in an allusive fashion which presupposes that the audience is already familiar with them.  The second lies in the way we have to understand these allusions, for while some are entirely unknown to us, we can connect most of them to more complete narratives in Snorre’s Edda.  Snorre’s Edda is, however, not a perfect source and since Snorre is often relying on \Voluspa\ we risk running into circular reasoning.  What is worse, in at least one instance (st.~38/1) Snorre shows willingness to alter the wording of \Voluspa\ in citation in order to fit his own narratives, which calls his reliability into question.  The third factor is that the witness manuscripts disagree over the order and inclusion of numerous stanzas.

In order to properly elucidate the myths and their broader religious and cosmic significance, we therefore have to take into account all available resources: comparative mythology, folkloristics, linguistics, historical sources, and archæology, and nowhere is this more important than in \Voluspa.  Although the body of Norse pre-Christian poetry is rather small, the available material for a broader comparative study is truly vast, and since there is still a tremendous amount of work to be done there will undoubtedly creep in both omissions and errors.  Even so, it is hoped that the present commentary along with that of the other poems included in this edition will provide new and interesting perspectives to the general reader, student, and scholar, not just of Germanic mythology but of the history of religions in general.

\subsubsection{The wallow}

As noted above, \Voluspa\ is narrated by a wallow.  This figure represents the northern European belief that women, especially in elder age, were invested with a unique psychic or prophetic foresight, a power which became even greater when they died and might come into possession of truly hidden knowledge---knowledge perhaps unknown even to the chief god Weden Himself.  In \Voluspa\ this ability of sight is reflected in the continual repetition of the phrase \emph{sá hǫ̇n} ‘she saw’, or, in the later part of the poem, \emph{sér hǫ̇n} ‘she sees’.  In her trance the wallow \emph{sees} the past and the future as if they were before her very eyes.

Accounts of famous northern European seeresses are found from the time of our earliest Roman sources.  During the Batavian rebellion of the late 1st c.~\ce, the seeress Veleda (a Celtic word, cf.~OIr. \emph{filed} ‘seer’, Middle Welsh \emph{gwelet} ‘id.’) of the Bructeri made herself famous for her spaes.  Tacitus, \emph{Germania} 8, mentions her by name when talking about the Germanic respect for women: \emph{Inesse quin etiam sanctum aliquid et providum putant, nec aut consilia earum aspernantur aut responsa neglegunt.  Vidimus sub divo Vespasiano Veledam diu apud plerosque numinis loco habitam.} ‘Moreover, they think that there is in them something holy and prophetic, and they neither despise their counsels nor disregard their responses.  We saw, under the deified Vespasian, that Veleda was held in the place of a deity for a long time among the multitude.’

In Jordanes, \emph{Getica} 121, we find mention of the Gothic witches known as the \emph{haliurunnae}, thought to be the ancestors of the hostile Huns.  Their name is clearly cognate with OE \emph{hęlli-ru̇ne} ‘witch, sorceress’, formed from \emph{hęll} ‘hell, Underworld’ + \emph{ru̇ne} ‘whisperess, confidante’, thus signifying ‘those who convene with the dead’.

The Longbeards, too, were once led by a prophetic woman, Gambara, as told in their 7th c.~\emph{Origo} (see Appendix to \textlink{Grimnismal}).

\subsubsection{Weden}

It was to acquire their knowledge that Weden in his incessant quest for wisdom undertook frequent journeys in order to question wallows and hell-whisperers alike about events past and future, and \Voluspa\ is not the only Eddic poem reflecting that motif.  The closest parallel is \textlink{Baldrsdraumar}, wherein Weden summons a wallow from her grave in order to find out why the god \inx[P]{Balder} is having ominous nightmares.  Weden’s journeys to convene with the dead are also alluded to in \textlink{Havamal}[158] and \textlink{Harbardsljod}[44]--45.

For the god’s general quest for knowledge there are numerous parallels: Weden’s wisdom contest against the wise ettin \inx[P]{Webthrithner} (\textlink{Vafthrudnismal}), his theft of the mead of poetry (\textlink{Havamal}[103]--110), his self-hanging and acquisition of the runes (\textlink{Havamal}[138], 139), and his giving up his eye to Mimer in exchange for wisdom (\textlink{Voluspa}[28]).

\subsubsection{Cycles and eschatology}

As has been impressively shown by \textcite{EliadeReturn}, pre-Abrahamic religions are centered around cycles and the power of \emph{being} inherent in their return; they are based around cyclical mythic time rather than linear history.  The sun at dawn, the new moon, the daylight at the winter solstice, the warmth of spring---traditional man did not think the return of these things certain.  Rather, the sun, moon, year, and indeed the world itself had to be strengthened and re-established through rituals harkening back to mythic time and re-establishing Creation.  So it is that modern man thinks little of that which traditional man worried about with great fright and consequently heralded with great rejoicing.

In the traditional worldview the world consists of nested cycles, a series of ever-revolving epicycles of rises and declines: the day, the month, the year, the human lifetime, and the world itself.  Each must rise and fall, decline and be regenerated.  In this way history is not linear, but cyclical; everything that is done has already been done before and in fact it is just this precedent which gives the present occurrence its realness.

\Voluspa\ focuses on the greatest cycle, that of the World itself, from its beginning as formless matter (3) through its Creation by the Gods (4--6) and Golden Age (7--8) to its gradual decline (8--42), destruction (43--55) and rebirth (56--62).  Although the poem is not explicit about it, there is a hint that the seed of evil is not entirely defeated in the next world (cf.~62 n.), leaving the possibility open for yet another repetition of the cycle.

TODO: mention Indian yuga-s, Zoroastrian texts.

\subsubsection{The Wiking Age eschatological fixation}

One of the most important developments of the Germanic religion when compared with other Indo-European religions is its unique obsession with eschatology.  Although the relevant myths are preserved only in the Norse literature, likely traces of the same are found on continental West Germanic territory in the form of the OHG \textlink{Muspilli} and the OS \textlink{Heliand}, both of which use the likely pre-Christian inherited term \emph{Muspilli}, also found in ON as \emph{Muspell}; see st.~49/1 n.~below, \textlink{Muspilli} 56 n.

To be clear, what distinguishes the Germanic (and specifically Norse) religion from other Indo-European systems is not the mere presence of ideas of the destruction and rebirth of the World, which as discussed above are also found in Indo-Iranian sources, but specifically a fixation with the imminent Destruction of the World and most of the Gods at the hands of the evil cosmic forces during the so-called Rakes of the Reins (\emph{ragna rǫk}).

We find allusions to the Destruction in a large number of pre-Christian Norse texts: many of the mythological poems (\Voluspa\ 43--55, \textlink{Vafthrudnismal}[17]--18, 44--53, \textlink{Grimnismal}[4]/4, 17, 24, 40, \textlink{Baldrsdraumar}[14], \textlink{Lokasenna}[10]/1, 39, 41--42, \textlink{Hyndluljod}[41], 43), several Scaldic eulogies (\Hakonarmal\ 20--21, \Eiriksmal\ 7), and a few poetic pre-Christian runestones (Sö 154 Skarpåker, perhaps also Vg 119 Sparlösa).

The monsters of the Rakes of the Reins also figure iconographically, e.g.~on the 10th c.~Gosforth cross, showing Wider slaying the Middenyardswyrm.  Pre-Christian memorial runestones frequently depict a single ravenous wolf, e.g.~N 84 Vang and DR 285 and 286 Hunnestad; sometimes the wolf is accompanied by a ship (perhaps Nailfare, 48/4 n.), e.g.~DR 271 Tullstorp (Figure~\ref{fig:tullstorp}, p.~\pageref{fig:tullstorp}) and Ög 181 Ledsberg.  These pictorial funerary motifs should likely be considered visual representations of the same attitude as found in Sö 154 Skarpåker, \Hakonarmal\ 20--21, and \Eiriksmal\ 7.

In the ubiquity of these allusions, including on runestones, we see that the Rakes of the Reins were not a mere a literary construct.  The worry about the impending destruction of the world was felt acutely on a personal level by the practitioners of the pre-Christian religion.  It is clear that the Heathens felt that the destruction was imminent and could come, as it were, “like a thief in the night” (cf.~\textlink{Heliand}[52][4358]--61).  It is notable that the association is particularly strong in texts relating to death, whether they be runestones for dead family members raised by more or less every-day people, or funerary poems for beloved kings.  We do not know whether the practitioners of the religion found consolation in these eschatological events, but it is hard not to be moved by the feeling of the father who, having buried his own son, proclaims that “Earth and Up-heaven shall be rent” (Sö 154).

\begin{figure}[b]
\centering
\includegraphics[width=\textwidth]{Tullstorp}
\caption{The Tullstorp stone (DR 271), showing a wolf above a ship, perhaps the Fenrerswolf and Nailfare.  The inscription is a memorial and not related to the iconography.  Wiking Age, ca. 1000 \ce.  © Jorchr. \href{https://creativecommons.org/licenses/by-sa/4.0/deed.en}{CC BY-SA 4.0}.  \url{https://commons.wikimedia.org/wiki/File:Tullstorpstenen,_DR_271,Tullstorp_1-1,_Runristning.jpg}}
\label{fig:tullstorp}
\end{figure}

\subsection{The climate catastrophe of 536 \ce}

If we are to look for a historical explanation for the above-mentioned eschatological fixation, the most likely candidate is undoubtedly the climate catastrophe of 536--537 \ce.  TODO. \textcite{Gräslund2008}.

Although ideas of the Destruction probably existed before 536, the plague, massive depopulation and accompanying social upheaval, perhaps alongside Christian apocalypticism spreading from the converted Gots and Vandals, likely led to a new form of religious expression.

\subsection{Textual preservation and transmission}

\textlink{Voluspa} is attested in three independent recensions, \Regius, \Hauksbok\ and \GylfMS.

The most important recension, \Regius, is the second oldest (ca.~1275 \ce).  This is the complete version of \Voluspa\ preserved in \Regius\ on foll. 1r--3r, where \Voluspa\ is the first poem.

The other complete recension, \Hauksbok, is the youngest (ca.~1310 \ce).  This is the version of \Voluspa\ preserved in \Hauksbok\ (foll. 20r--21r), where the poem is found in the middle of an impressive collection of native Norse-Icelandic saws and learned Catholic works.

The third recension, \GylfMS, is the oldest (ca.~1225 \ce), and represents the stanzas of \Voluspa\ (3, 5, 9, 10--13, 15, 18, 24--25, 28, 37--40, 44--46, 48--54, 61) cited in \Gylfaginning, the first part of Snorre’s Edda.  This recension is fragmentary, though many now-lost stanzas belonging to it may be gleaned through Snorre’s paraphrases of them.  At times the paraphrases can provide important critical readings and serve as a tiebreaker for variants that differ between \Regius\ and \Hauksbok, e.g. in st.~19/2, where \emph{sal} ‘hall’ in the \Gylfaginning\ paraphrase agrees with \Hauksbok\ against \Regius\ \emph{sę́} ‘lake’.  For the four mss. of \Gylfaginning\ (\RegiusProse, \Trajectinus, \Wormianus~TODO, and \Upsaliensis) see the Introduction to Mythic Poetry above.

\subsubsection{Table of stanzas}

The order of stanzas varies greatly across the three recensions and is best illustrated by a table.

Stanzas cited in \Gylfaginning\ have been handled in the following way: those cited independently are marked with a plus sign, while those belonging to a sequence of several stanzas are marked with an incrementing letter based on the sequence and a number for its position, so that \emph{A1} is the first stanza in the first sequence, \emph{B1} the first stanza in the second sequence, and so forth.

When a stanza found in a manuscript is strongly divergent (e.g.~st.~10, where \Gylfaginning~14:15 omits the first two half-lines), its mark is followed by a star.

The stanzas beginning with the formula \emph{Þȧ gingu ręgin ǫll} ‘Then went all the Reins’ et c. are represented by the immediately following half-line.

\begin{small}\begin{longtabu} to \textwidth {|c c c c c c|}
	\hline
	\multicolumn{2}{|c}{\emph{pres. ed.}} & \Regius & \Hauksbok & \RegiusProse\Trajectinus\Wormianus & \Upsaliensis \\ [0.5ex]
	\hline\hline\endhead
	\hline\endfoot
	1 & Hljóðs bið’k allar & 1 & 1 & − & − \\
	2 & Ek man jǫtna & 2 & 2 & − & − \\
	3 & Ár vas alda & 3 & 3 & + & + \\
	4 & áðr Burs synir & 4 & 4 & − & − \\
	5 & Sól varp sunnan & 5 & 5 & +* & +* \\
	6 & \dots\ nǫ́tt ok niðjum & 6 & 6 & − & − \\
	7 & Hittu⸗sk ę̇sir & 7 & 7 & − & − \\
	8 & Tęflðu ï tu̇ni & 8 & 8 & − & − \\
	9 & \dots\ hvęrr skyldi dverga & 9 & 9 & B1 & B1 \\
	10 & Þar vas Móð·sognir & 10 & 10 & B2* & B2* \\
	11--15 & \emph{Dwarf-tallies} & 11--15 & 11--16 & + & + \\
	16 & Und’s þrír kvǫ̇mu & 16 & 17 & − & − \\
	17 & Ǫnd þau né ǫ́ttu & 17 & 18 & − & − \\
	18 & Ask vęit’k standa & 18 & 19 & + & + \\
	19 & Þaðan koma męyjar & 19--20 & 20--21 & − & − \\
	20 & Þat man hǫ̇n folk-víg & 21--22 & 27 & − & − \\
	21 & Hęiði hana hétu & 23 & 28 & − & − \\
	22 & \dots\ hvárt skyldu ę̇sir & 24 & 29 & − & − \\
	23 & Flęygði Óðinn & 25 & 30 & − & − \\
	24 & \dots\ hvęrr hęfði lopt alt & 26 & 22 & C1 & C1 \\
	25 & Þȯrr ęinn þar vá & 27 & 23 & C2* & C2* \\
	26 & Vęit hǫ̇n Hęim·dalar & 28 & 24 & − & − \\
	27 & Ęin sat hǫ̇n úti & 29 & − & − & − \\
	28 & Alt vęit’k, Óðinn & 29 & − & + & + \\
	29 & Valði hęnni Hęr-fǫðr & 30 & − & − & − \\
	30 & Sá hǫ̇n val-kyrjur & 31 & − & − & − \\
	31 & Ek sá Baldri & 32 & − & − & − \\
	32 & Varð af męiði & 33 & − & − & − \\
	33 & Þó ę́va hęndr & 34 & − & − & − \\
	34 & þar sitr Sigyn & 35 & 32* & − & − \\
	35 & Ǫ́ fęllr austan & 36 & − & − & − \\
	36 & Stóð fyr norðan & 36 & − & − & − \\
	37 & Sal sá hǫ̇n standa & 37 & 36 & E1 & E1 \\
	38 & Sá hǫ̇n þar vaða & 38 & 37 & E2* & E2* \\
	39 & Austr býr hin aldna & 39 & 25 & A1 & A1 \\
	40 & Fylli⸗sk fjǫrvi & 40 & 26 & A2 & A2 \\
	41 & Sat þar ȧ haugi & 41 & 34 & − & − \\
	42 & Gól of ǫ̇sum & 42 & 35 & − & − \\
	43, 47, 55 & Gęyr (nú) Garmr mjǫk & 43, 46, 55 & 33, 38, 43, 48, 51 & − & − \\
	44 & Brǿðr munu bęrja⸗sk & 44 & 39 & − & − \\
	45 & Lęika Mïms synir & 45 & 40 & D1* & D1* \\
	46 & Hvat ’s með ǫ̇sum? & 49 & 42 & D2 & D2* \\
	48 & Hrymr ękr austan & 47 & 44 & D3 & − \\
	49 & Kjóll fęrr austan & 48 & 45 & D4 & − \\
	50 & Surtr fęrr sunnan & 50 & 46 & +, D5 (n.) & + \\
	51 & Þȧ kømr Hlïnar & 51 & 47 & D6 & − \\
	52 & Þȧ kømr hinn mikli & 52 & − & D7 & − \\
	53 & Þȧ kømr hinn mę́ri & 53* & 49* & D8 & − \\
	54 & Sól tér sortna & 54 & 50 & D9 & − \\
	56 & Sér hǫ̇n upp koma & 56 & 52 & − & − \\
	57 & Finna⸗sk ę̇sir & 57* & 53 & − & − \\
	58 & Þar munu ęptir & 58 & 54 & − & − \\
	59 & Munu ȯ·sánir & 59 & 55 & − & − \\
	60 & Þȧ kná Hø̇nir & 60 & 56 & − & − \\
	61 & Sal sér hǫ̇n standa & 61 & 57 & + & + \\
	62 & Þar kømr hinn dimmi & 62 & 59 & − & − \\
	H31 & Þȧ kná Váli & − & 31 & − & − \\
	H41 & Hrę́ða⸗sk allir & − & 41 & − & − \\
	H48 & Gïnn lopt yfir & − & 48 & − & − \\
	H58 & Þȧ kømr hinn ríki & − & 58 & − & − \\ [1ex]

	\hline
\end{longtabu}\end{small}

\subsection{Summary}

It seems that immediately prior to the poem’s beginning, Weden has awakened the wallow out of her grave.  He commands her to speak.

On the basis of its refrains and content, the poem may be divided into the following rough sections: introduction (1--2), Creation and the early age of the World (3--17), cosmological and mythological exposition and decline (18--42), Destruction (43--55), and Rebirth (56--62).

The wallow’s \emph{spae} begins with a bid for silence (1), followed by her earliest memories of being raised amongst ettins (2).

She begins by recounting the Creation (3--6) and the following Golden Age (7--8), which was brought to an end by the intrusion of three unidentified ettin-maidens (8).  After this, she describes the making of the dwarfs (9--10), whose names are listed in several separate \emph{dwarf-tallies} which are without doubt later inserts (11--15).  She mentions how the Gods gave life to the first human beings (16--17), and then describes the Ugdrassle’s Ash (18) and the three \inx[P]{norns} living under it (19).

After st.~19, the order of the stanzas in the two full recensions---\Regius\ and \Hauksbok---diverges significantly, and does not come to agree again until st.~43.  With a single exception (st.~46), the order adopted for the present edition is that of \Regius, but since it is hard to decide which one is more original both are summarised here.

In \Regius, the wallow recalls how a woman named Goldwey was sacrificed and reborn three times (20), and how she, under the name Hode, practiced sorcery and witchcraft (21).  She then recalls the first war in the world, between the Eese and Wanes (22--23), and the killing of the smith who was promised \inx[P]{Frow} and the sun and moon in exchange for building the wall of Osyard (24--25).  This is followed by a cryptic verse describing Mimer’s Well (26).  She then seems to tell of her own origin; how she was approached by Weden, who gave her rings and jewels in exchange for her wisdom (27--29).  She then tells of the walkirries (30) and relates the death of Balder (31--33) and the subsequent binding of Lock (34).  After this, she describes a sword-river (35) and two beer-halls (36), then the torture of criminals in the afterlife by Nithehewer (37--38).  She sees Ironwood in the East, where the wolves who will swallow the Sun and Moon are being reared (39--40).  She foresees the crowing of three cocks: one in Ironwood (41), the other among the Gods, and the third in \inx[P]{Hell} or the Underworld (42).  After this, we get the first occurrence of the refrain-stanza of the destruction; it describes how the wolf breaks free, and is repeated several times throughout the part of the spae that deals with the destruction of the world (43).

In \Hauksbok\ the structure is quite different.  The Eese immediately go to decide what action to take regarding the Master Builder (24--25), and Mimer’s well is described (26).  Then follow the two stanzas about the wolves which will swallow the sun and moon (40--41), and the story of Goldwey and the war between the Eese and Wanes (20--23).  The stanzas about Goldwey’s origin (27--29), the walkirries (30), and the death of Balder (31--33) are left out of \Hauksbok, along with the sword-river and beer-halls (35--36), but it does describe the binding of Lock (34).  Next we get the first occurrence of the refrain-stanza (43), then the crowing of the three cocks (41--42), and lastly the punishment of criminals (38--39).  \Hauksbok\ then repeats the refrain-stanza a second time, and \Regius\ and \Hauksbok\ are now almost completely aligned for the rest of the poem.

The Destruction of the world begins with numerous catastrophes and a decay in human morality: kinslaying, incest, adultery, and great wars become widespread (44).  Homedall blows the great Yellhorn, the Ash of Ugdrassle is set ablaze, and a trapped ettin comes loose (45).  The Eese, Elves, and Dwarfs are all distressed (46).  Here the refrain-stanza is repeated again (47).  Now the enemies of the Gods launch their attack: the ettin Rim drives from the east, the wyrm Ermingand writhes in the sea and propels huge waves, and the ship Nailfare comes loose (48).  The ship drives from the east and is led by Lock and the Wol, both of whom have at last come free (49).  From the south Surt brings a flaming sword to set fire to the world; there is a crack in Heaven and Hell is joined by many men (50).  Weden goes out to fight the Wolf and falls, as does Free against Surt (51).  Wider avenges Weden by killing the Wolf (52), while Thunder and the Wyrm kill each other; the great battle leaves the world uninhabited (53).  Instead of 52, \Hauksbok\ has a stanza describing how the Wyrm gapes over the entire world, and how Thunder will ‘meet its venom’ (H48).  The Destruction culminates in the sun turning black, the stars disappearing, and the earth sinking into the sea.  The whole world burns (54).  The refrain-stanza repeats a final time (55).

The world is reborn.  The earth rises from the sea and soon turns green again; the cycle of nature is renewed (56).  The surviving Eese find each other on the Idewold, where they had first settled, and reminisce (57).  A new Golden Age begins (58).  The fields grow high without toil and Hath and Balder reconcile (59).  Heener becomes the head priest of the Gods and the sons of Hath and Balder settle the sky (60).  The wallow sees a fair, golden hall in which virtuous generations reside (61), but her vision concludes with the sight of a flying, dusky dragon---it is Nithehewer, carrying corpses in its mouth, a sign that there is still a seed of evil left in the new world.  With this portent she sinks back down into her grave (62).

\clearpage

\section{Text}

\bvg\bva\mssnote{\Regius~1r/2, \Hauksbok~20r/1}%
\Ballnote{The wallow has just been awoken from her grave (cf.~st.~62/4b below) and commences her speech by asking for the silence of gods and men, a meristic expression \parencite[99--100]{West2007}.  The whole introductory formula has Indo-European parallels; see \textcite[63,92--93,312]{West2007}.}%
\llap{„}\edtrans{\edtext{\alst{H}ljóðs}{\Afootnote{The \emph{H} is an ornamented initial 5 lines in height, noticeably taller than any other initial save the one introducing \textlink{HelgakvidaOne}, in \Regius \emph{;} introduced by a small initial 2 lines in height in \Hauksbok.}} bið’k}{For hearing I ask}{\Bfootnote{This is the only call for hearing found in an Eddic poem, but many parallels are found in Scaldic praise poetry, most often involving the related words \emph{hljóð} ‘hearing’ and \emph{hlýða} ‘hear, listen’ \parencite[424, 250]{Wood1960}.
Examples abound (see \cite{Wood1960}), but most notably among them, the same introductory expression is found in Eyel’s Head-ransom (Egill \emph{Hfl} 2 in \Skp\ 5): \emph{hljóðs biðjum hann} ‘for hearing we [I] ask him’.}} allar \hld\ \edtrans{\alst{h}ęlgar}{holy}{\Afootnote{so \Hauksbok \emph{;} om. \Regius}\Bfootnote{That the omission of this word in \Regius\ is nothing more than a scribal error is clearly shown by the meter;
the a-verse in \emph{Hljóðs bið \alst{e}k \hld\ \alst{a}llar kindir} is only three syllables long and has highly unnatural alliteration on the unstressed pronoun \emph{ek} rather than the first nominal \emph{hljóðs}.}} kindir, &
\edtrans{\alst{m}ęiri ok \alst{m}inni}{greater and lesser}{\Bfootnote{It is ambiguous to which noun phrase these adjectives belong.  It can either be (a) “all greater and lesser holy races”, which might be equivalent to the phrase \inx[F]{Eese and Elves} (earthly and heavenly supernatural beings; see Index for occurrences), or (b) “the greater and lesser sons of Homedall”.  (b) is preferred as the simpler reading since it avoids enjambment.  The sense of “greater or lesser” is then more likely to be literal and physical (i.e., the younger and older members of the audience) rather than figurative (i.e., the varying social classes).}} \hld\ \edtrans{\alst{m}ǫgu Hęim·dalar}{sons of Homedall \ken{men}}{\Bfootnote{Usually interpreted as a kenning for \textsc{men} on the basis of \textlink{Rigsthula}, according to which Homedall sired the three social classes of mankind.  This convention is followed here as it is felt to give the better sense and conforms to the Indo-European poetic phrase “gods and men” (see note to ALL), but such a kenning is otherwise unparalleled.
A more typologically likely interpretation, following \textcite[23, 252]{Meissner1921}, would be to see \emph{męgir} in the sense ‘lads, comrades’.  This makes the kenning equivalent to \emph{Hropts męgir} ‘Roft’s lads’ (\textlink{Lokasenna}[45]/3), \emph{allar áttir Ingvi-Fręys} ‘all the kin of Ing-Free \ken{gods}’ (\Haustlong\ 10), \emph{Njarðar niðjar} ‘Nearth’s kinsmen’ (Hfr Lv 10 in \Skp\ 5).}}; &
\alst{v}ilt at, \edtrans{\alst{V}al-fǫðr}{Walfather}{\Bfootnote{That is, the ‘Father of the Slain’, a name for Weden.  The wallow probably uses this name since Weden, acting in his role as necromancer and psychopomp, has just awoken her out of her grave (perhaps with a \emph{val-galdr} ‘slain-galder’ as in \textlink{Baldrsdraumar}[4]/3 and \textlink{Havamal}[158]).}}, \hld\ \alst{v}ęl fram tęlja’k &
\alst{f}orn spjǫll \alst{f}ira, \hld\ \edtrans{þau’s \alst{f}ręmst of man}{which I earliest recall}{\Bfootnote{Cf.~\textlink{Vafthrudnismal}[34]--35 with similar phrasing.}}.\eva

\bvb “{\huge F}\textsc{or hearing I ask} all holy races \ken{gods}, \\
the greater and lesser sons of \inx[P]{Homedall} \ken{men}! \\
Thou wouldst, O \inx[P]{Walfather}, that I well tell forth \\
the ancient sayings of men, the earliest I recall.\evb\evg


\bvg\bva\mssnote{\Regius~1r/4, \Hauksbok~20r/2}%
\alst{E}k man \alst{jǫ}tna \hld\ \alst{á}r of borna, &
þȧ’s \alst{f}orðum mik \hld\ \alst{f}ǿdda hǫfðu; &
\edtext{\alst{n}íu man’k hęima}{\lemma{níu \dots\ hęima ‘nine Homes’}\Bfootnote{Or ‘nine realms’.  The identity of these is by no means clear nor is there any canonical list of which realms are supposed to be included.}}, \hld\ \alst{n}íu \edtext{ï·viðjur}{\Afootnote{so \Regius\Hauksbok; previously read \emph{‘iviði’} in \Regius, but this was made obsolete by an x-ray scan undertaken by \textcite{StefanKarlsson1979} which revealed a previously overlooked abbrev.~mark standing for \emph{-ur}.}\lemma{ï·viðjur ‘Inwithies’}\Bfootnote{Evil-working women or ogresses; their identity is by no means clear.
The f. \emph{jōn}-stem \emph{ï·viðja} also appears in a list of poetic synonyms for trollwomen (Þul \emph{Trollkvenna} 3 in \Skp\ 3).  A commonly suggested etymology is \emph{í} ‘in’ + \emph{viðr} ‘tree, forest’ (i.e., ‘tree, forest-dwellers’), but this would be an unusual formation and leaves the \emph{-j-} unexplained.  A more plausible etymology is a derivative of \emph{*ï·við} ‘malice’, attested in the cpd. \emph{ï·við-gjarn} ‘malicious’ (\textlink{Volundarkvida}[28]).
This derivation can also explain the \emph{-j-}, since the WGmc. cognates OE \emph{in·wid}, OS \emph{in·wid}, and OHG \emph{in·wit} show it to be a neutr. \emph{ja}-stem.}}, &
\edtrans{\alst{m}jǫt-við \alst{m}ę́ran \hld\ fyr \alst{m}old neðan.}{the famed Measure-Tree beneath the soil.}{\Bfootnote{Probably \inx[P]{Ugdrassle’s Ash}, still but a seed.  The sense of \emph{mjǫt-viðr} ‘Measure-Tree’ is not clear.  It is probably not the \emph{mjǫtuðr} in 45/1 below.}}\eva

\bvb I recall \inx[P]{Ettins} born of yore, \\
those who formerly had reared me. \\
I recall nine \inx[C]{Home}[Homes], nine \inx[P]{Inwithies}, \\
the famed Measure-Tree beneath the soil.\evb\evg


\bvg\bva\mssnote{\Regius~1r/6, \Hauksbok~20r/4, \\
\RegiusProse~2r/25, \Trajectinus~3r/2, \\
\Wormianus~TODO, \Upsaliensis~TODO}%
\Ballnote{The wallow begins her retelling of the mythic history of the World.  Told in sts.~3--4 is the creation of the Earth and the Heavens.  Other Norse poetic sources for the creation are \textlink{Grimnismal}[41]--42 and \textlink{Vafthrudnismal}[21].  A Christian creation hymn apparently drawing on older pagan sources is the OHG \textlink{Wessobrunn} edited below under Poetry on Christian Subjects; cf.~also the more thoroughly Christian \textlink{Cadman} in the same section.
The beginning of the Norse creation narrative is treated at length in \Gylfaginning~4--5, according to which the world first consisted of two extremities: the frozen Nivelham in the north and the burning Muspellsham in the south.  From Nivelham the icy venomous rivers called the \inx[P]{Ilewaves} rushed forth until they lost their momentum and froze to ice, while from Muspellsham flowed burning lava.  The ice and the lava met in the Gap of Ginnings (\emph{Ginnunga-gap}, which \emph{var svá hlę́tt sem lopt vind-laust} ‘was as calm as windless air’, \Gylfaginning~5:7) and there combined to form the first being, \inx[P]{Yimer}.}%
\edtrans{\alst{Á}r vas \edtext{\alst{a}lda}{\Afootnote{\emph{halda} \RegiusProse}}}{It was early of ages}{\Bfootnote{By these words the wallow begins her narrative in mythic time, in the earliest age (\emph{ǫld}, cf.~st.~44/4--5 below), \emph{in illo tempore}.
It is formulaic for a mythological or heroic-legendary Eddic poem to begin with \emph{ár} ‘of yore, early’; cf.~\textlink{Hymiskvida}[1]/1, \textlink{HelgakvidaOne}[1]/1 (which likewise has \emph{Ár vas alda}, prob.~borrowed from the pres.~st.), \textlink{GudrunOne}[1]/1,
\textlink{Sigurdskamma}[1]/1.}} \hld\ \edtext{þar’s}{\lemma{þar’s Ymir byggði ‘where Yimer dwelled’}\Afootnote{\emph{þat’s ękki vas} ‘when nothing was’ \RegiusProse\Trajectinus\Wormianus \emph{;} \emph{þar’s ękki vas} ‘where nothing was’ \Upsaliensis}} \edtrans{\alst{Y}mir}{Yimer}{\Bfootnote{The first ettin, from whom the rest were begotten, at first a- or ambisexually.  Yimer was the Heaven and Earth conjoined and as such hermaphroditic, containing both a male (heavenly) and female (earthly) aspect.
He is described in \Gylfaginning~5:8: \emph{Ok þá er mǿttisk hrímin ok blę́r hita’ns svá at bráðnaði ok draup, ok af þeim kviku-dropum kviknaði með krapti þess er til sendi hita’nn, ok varð manns líkandi, ok var sá nefndr Ymir.} ‘And when the rime-frost and hot air from the heat met so that it thawed and dripped, then from those flowing droplets there arose a life with the power which brought the heat and became the likeness of a man, and he was named Yimer.’

The comparative mythological idea of Yimer as the reflex of a PIE cosmic twin must be disregarded, for his name derives from PIE \emph{*ym̥H-yós}, which, although ultimately from the root \emph{*yemH-} ‘twin’ (whence OI \emph{Yamá} ‘twin’), is most closely related to Latvian \emph{jumis} ‘paired couple, double fruit’ \parencite[274]{Kroonen2013} and thus seems to suggest hermaphroditism or the union of the copulating Earth-Heaven.
Yimer is probably (so \Gylfaginning~5:9) to be identified with the ambisexually reproducing first ettin Earyelmer in \textlink{Vafthrudnismal}[30]--33.
Comparatively he corresponds to the sexually joined \textgreek{Οὐρανός} (‘Heaven’) and \textgreek{Γαῖα} (‘Earth’) prior to the castration of the former at the hands of their son \textgreek{Κρόνος} (‘the Cutter’; \Theogony\ 159--172), who was only later identified trough folk etymology with \textgreek{Χρόνος} (‘Time’).  This violent separation ended the primæval hermaphroditic begetting and ushered in the patriarchal rule of the male Sons of Byre; it is of note that Yimer, like \textgreek{Οὐρανός}, is described as evil (\Gylfaginning~5:15).

Yimer further resembles the Vedic \textoi{Púruṣa} and the Chinese \emph{Pangu} in his role as the primæval being from whose body parts the World is shaped \parencite{Witzel2017}.  Cf.~\textlink{Vafthrudnismal}[21], \textlink{Grimnismal}[43] n.}} byggði, &
vas⸗a \alst{s}andr né \edtext{\alst{s}ę́r}{\Afootnote{\emph{sjár} \Wormianus \emph{;} \emph{sjór} \Hauksbok\Upsaliensis}}, \hld\ né \edtext{\alst{s}valar unnir}{\Afootnote{\emph{†sva†}, with a large space following \Trajectinus}}; &
\edtext{\alst{jǫ}rð fann⸗sk \edtrans{\alst{ę́}va}{never}{\Afootnote{\emph{ęigi} ‘not’ \RegiusProse\Trajectinus\Upsaliensis}} \hld\ né \alst{u}pp-himinn}{\lemma{jǫrð \dots\ né upp-himinn ‘Earth \dots\ nor Up-heaven’}\Bfootnote{A well-attested formulaic cosmological word-pair found in all four Old Germanic languages with alliterative poetic traditions (ON, OE, OS, and OHG), especially in the context of the creation and destruction of the world.  See Index: \inx[F]{Earth and Up-heaven}.}}; &
\edtrans{\alst{g}ap vas \alst{g}innunga}{there was the gap of hawks \ken{air/midspace}}{\Bfootnote{\emph{gap ginnunga} ‘gap of hawks’ has usually been interpreted on the basis of the cosmogony in \Gylfaginning~4--5, where Snorre presents the \emph{Ginnunga-gap} ‘Gap of Ginnings’ as a primordial physical location where cold and hot elements combine to form the first differentiated cosmic being, Yimer, in the place where the Earth will later be situated.
For this reason the phrase is typically rendered in translation simply as “Ginnungagap” and explained as a sort of supernatural void.  The present stanza, however, is the only conjunction of the words \emph{gap} and \emph{ginnunga} outside of Snorre’s Edda, and his cosmogony is very peculiar.  There is therefore good reason to reexamine these two words in the present context in order to determine the exact nature of the ‘gap’.%
% TODO: Read and cite deVries1930

To begin with, we must reject the traditional renditions “yawning gap” or “supernatural void” as unlikely.  The former may be eliminated immediately since the oft-connected ON verb \emph{gína} ‘yawn, gape’ has a different root (\emph{gín-}) than \emph{ginnunga} (\emph{ginn-}) and since there is no clear derivation from \emph{gína} to \emph{ginnunga}, which in any case must be a noun, not an adjective and certainly not a participle.
\emph{ginnunga} has been explained variously.  \textcite{FinnurEdda} sees it as gen.~sg.~of \emph{*ginnungi} (from the lost equivalent to OE \emph{ginne} ‘wide, spacious’ + the nominal suffix \emph{-ungi}).  This is, however, a strange and decidedly non-Norse construction, especially since the suffix \emph{-ungi} forms masculine nominals with animate senses, not abstract nouns.
A more common explanation is to see \emph{ginnunga} as gen.~of a pl.~\emph{*ginnungar} ‘supernatural beings’, derived from \emph{ginn-} ‘yin-’ in expressions like \emph{ginn-ru̇nar} ‘yin-Runes’, \emph{ginn-hęilǫg goð} ‘yin-holy Gods’, and \emph{ginn-ręgin} ‘yin-Reins’.  It is not at all clear that \emph{ginn-} means ‘supernatural’; it is more likely to be cognate with OE \emph{ginne} ‘wide, spacious’, and thereby serves to strengthen the nouns it follows (e.g.~‘vast runes’, etc.).  This explanation also requires us to assume the existence of nebulous ‘supernatural beings’ called ginnings in the primæval nothingness, an idea for which there is no support.

I find a much more natural reading in \citeauthor{Meissner1921}, who sees \emph{gap ginnunga} as a kenning ‘gap of hawks \ken{air/midspace}’, where \emph{ginnunga} is gen.~pl.~of the well-attested poetic word \emph{ginnungr} ‘hawk’.
The proposed kenning is conventional and belongs to the type “land, path of the bird \ken{air}” \parencite[108]{Meissner1921}.  Indeed, the determinant \emph{ginnungr} is also found in a \textsc{sky}-kenning in \Haustlong\ 15: \emph{ǫll endi-lǫ́g ginnunga vé} ‘all the low mansions of hawks \ken{skies} from end to end’.
This interpretation is confirmed by an underappreciated passage in \Skaldskaparmal\ 74 which lists \emph{ginnunga-gap} among poetic synonyms for the air: \emph{Lopt heitir ginnunga-gap ok meðal-heimr, fogl-heimr, veðr-heimr.} ‘Air is called \emph{ginnunga-gap} (‘gap of hawks’) and middle-home, bird-home, weather-home.’  The sense of \Voluspa ’s \emph{gap ginnunga} ‘gap of hawks’ is thus the same as \Haustlong’s \emph{vé ginnunga} ‘mansions of hawks’, viz.~‘air, atmospheric sky’.

Having determined the likely sense of the phrase, we should consider its significance in the present context.  This requires a cosmological excursion.
The Old Germanic pagan cosmology, like that of many other ancient religions (e.g. the Vedic with its two \textoi{rṓdasī} ‘World-Halves’, \Rigveda\ 1.10.8, 5.85.3, 7.86.1, et c.) sees the cosmos as consisting of two world-halves: Earth and Up-heaven, the land of Men and the land of Gods.  The Earth is held to be a flat disc, while the Heaven (or Firmament) is a round dome above it, held up by four dwarfs (ON \emph{dvergar} ‘dwarfs’ can also mean ‘structural pillars’) lest it fall down (\Gylfaginning~8:10; cf.~Arn \emph{Þorfdr} 24 in \Skp\ 2).
Both Earth and Heaven are believed to be physically solid realms and made out of different parts of Yimer’s body: the Earth from his torso, the Heaven from his skull (\textlink{Grimnismal}[41], \textlink{Vafthrudnismal}[21]).  Like the Biblical Firmament, the Heaven is made out of stone (cf.~\textlink{Harbardsljod}[15]), which also explains the falling of meteors.

Between the two solid World-halves are the atmospheric skies, the Air (ON \emph{lopt}) or the Midspace (ON \emph{meðal-hęimr} ‘middle-home’) equivalent to the \textoi{antárikṣam} of the Vedic hymns (e.g. \Rigveda\ 4.40.5, 5.85.3), wherein the clouds---Yimer’s brains---float (\textlink{Grimnismal}[42]).  The important distinction between the empty skies and the solid Up-heaven or Firmament is seen in the aforementioned \Haustlong, which distinguishes between the ‘low \textsc{skies}’ (\Haustlong\ 15) and \emph{upp-himinn} ‘Up-heaven’ (\Haustlong\ 16).

To conclude, the primordial state expressed in the present stanza is one where there is only emptiness or Chaos (cf.~\textlink{Vafthrudnismal}[25]/2 n.), while Earth and Stoney Heaven are still monstrously conjoined in the body of Yimer.}}, \hld\ ęn \alst{g}ras \edtrans{hvęr⸗gi}{nowhere}{\Afootnote{\emph{ękki} ‘not’ \Hauksbok\RegiusProse\Trajectinus\Wormianus}},\eva%TODO: bibliography, Meissner is 20. Luft, b

\bvb It was early of ages where \inx[P]{Yimer} dwelled; \\
there was not sand nor sea nor cool waves. \\
\inx[P]{Earth} was never found nor \inx[P]{Up-heaven}; \\
there was the \inx[P]{gap of hawks} \ken{air/midspace} but grass nowhere,\evb\evg


\bvg\bva\mssnote{\Regius~1r/8, \Hauksbok~20r/5}%
\Ballnote{The Sons of Byre, the first Gods, came to life subsequent to Yimer.  They sacrificed the primordial ettin through decapitation and created the World out of his body, as told in \textlink{Grimnismal}[41]--42, \textlink{Vafthrudnismal}[21], and \Gylfaginning~6--8.}%
áðr \edtrans{\alst{B}urs synir}{the Sons of Byre}{\Bfootnote{In \Gylfaginning~6:8--10 identified as the three brothers Weden (ON \emph{Óðinn}), Will (\emph{Vili}), and Wigh (\emph{Véi}), the first Gods.  For their identities cf.~st.~16/2: \emph{ę̇sir} n.~below.}} \hld\ \edtrans{\alst{b}jǫðum of ypðu}{brought up the flatlands}{\Bfootnote{The verb \emph{yppa} is a causative derived from ON \emph{upp} ‘up’, \emph{uppi} ‘above, up, on high, remembered’ (cf.~below st.~15/6a).  In poetry it typically means ‘extol, lift up, announce’.  Its OE cognate \emph{yppan} means ‘bring up, bring forth, reveal’.

There is an apparent contradiction here, since the verb implies that the earth was covered by the sea; something stated clearly in st.~56 below.  This is problematic since st.~3/2 states plainly that there was neither earth nor sea.  The problem can be resolved in two ways.
(1) According to \Gylfaginning~7:3--5, the killing of Yimer flooded the world with so much blood that it drowned all the ettins except for Bearyelmer and his family, who survived by taking refuge on his canoe (\emph{lúðr}, see \textlink{Vafthrudnismal}[35] for discussion on this word and a citation of the relevant passage).  Thus, the present stanza would likewise allude to this primæval post-sacrificial flood and the subsequent raising of the earth, Yimer’s blood-covered trunk.
(2) If the \GylfMS\ version of l.~3/1b (\emph{þat’s ękki vas} ‘when nothing was’) is authorial, the wallow may not have intended to allude to the myth of Yimer at all, but may instead have been thinking of something like Genesis 1:2, 9--10, where Elohim makes dry earth rise up from the deep waters.  Something similar is found in \Iliad\ 14.201, 246, which state that \textgreek{Ὠκεανος}, Ocean, is the origin of all Gods and, indeed, of all things.  The problem with explanation (2) is that st.~3/2 above states plainly that there was not yet any sea before the sons of Byre.  The Greek Ocean is equivalent to the Ilewaves, the primæval venomous waters whence himself Yimer originated (\textlink{Vafthrudnismal}[31]).  As such, (1) should be preferred.}}, &
þęir es \edtrans{\alst{M}ið·garð}{Middenyard}{\Bfootnote{The Middle Enclosure, which was created as the home of Men (\textlink{Grimnismal}[42]).}} \hld\ \alst{m}ę́ran skópu; &
\alst{s}ól skęin \alst{s}unnan \hld\ ȧ \alst{s}alar stęina; &
þȧ vas \alst{g}rund \alst{g}róin \hld\ \edtrans{\alst{g}rø̇num lauki}{green leek}{\Bfootnote{A sign of the golden age, for the leek (in the Old Germanic languages referring more broadly to any plant of the genus \emph{Allium}) was in ancient times held to be the noblest plant.  See Index.}}.\eva

\bvb before the \inx[P]{Sons of Byre} brought up the flatlands, \\
they who created famed \inx[P]{Middenyard}. \\
The sun shone from the south on the stones of the hall; \\
then was the ground grown with green leek.\evb\evg


\bvg\bva\mssnote{\Regius~1r/11, \Hauksbok~20r/7, \\
\RegiusProse~3r/21, \Trajectinus~4r/2, \\
\Wormianus~TODO, \Upsaliensis~TODO}%
\Ballnote{The Sun lifts herself up over the horizon and rises for the first time, but she does not know her proper course.  The heavenly bodies are in a state of confusion and disorder.}%
\edtext{\alst{S}ól}{\linenum{|1|||2}\lemma{Sól \dots\ himin-jǫður ‘The Sun \dots\ heaven’s rim’}\Afootnote{om. \GylfMS.}} varp \alst{s}unnan, \hld\ \edtrans{\alst{s}inni Mȧna}{Moon’s companion}{\Bfootnote{At times mistranslated as ‘her moon’, understanding \emph{sinni} as dat. sg. f. of \emph{sïnn} ‘its (reflexive)’.  This cannot be correct since ON possessives are inflected based on the gender of the noun possessed, not the gender of the possessor.  \emph{mȧni} ‘moon’ is masculine, and so ‘her moon’ would be \emph{*sïnum Mȧna}.}}, &
\alst{h}ęndi hinni \alst{h}ǿgri \hld\ \edtrans{of \alst{h}imin-jǫður}{over heaven’s rim}{\Afootnote{composite; \emph{of himin-†iodyr†} \Regius\emph{;} \emph{of jǫður} \Hauksbok.}\Bfootnote{A problematic line.
(1) The reading of \Regius\ may be normalised as \emph{of himin-jó-dýr} ‘over the heaven-horse-beast(s)’ (so some recent eds., e.g.~\cite{PettitEdda}), but such a compound is nonsensical and gives the whole verse the otherwise unattested stress pattern x\#pxSS.  On the other hand, the reading of \Hauksbok\ (\emph{‘of ioður’} ‘over the rim, edge’) is clearly deficient, since it lacks the necessary alliteration on \emph{h-}.
(2) If we instead see \Regius’s \emph{‘iodyr’} as corrupted from an archetypal \emph{*iodur}, we can restore to \emph{of *himin-jǫður}, as is done here (so also e.g.~\cite{FinnurEdda}, \cite[292]{EddukvaediOne}).
Although the better reading, (2) also has problems, for in \Voluspa\ a type C verse with resolution on the first lift cannot have a short second lift, which \emph{of *himin-jǫður} (x\#pxsx) has \parencite[161]{Males2024Emend}.}}; &
\alst{S}ól þat né vissi, \hld\ hvar hǫ̇n \alst{s}ali átti; &
\edtext{\edtrans{\alst{M}ȧni þat né vissi, \hld\ hvat hann \alst{m}ęgins átti}{the Moon knew not what sort of might he had}{\Bfootnote{For a concrete example of this “might” cf.~\textlink{Havamal}[137]/7, where the moon is to be invoked in a blood feud.  A belief in the power of the moon is found in most ancient religions.  Its influence on the human mind was remarked already by \textcite{PlinyNatHist}, 2.102, 18.75 and is witnessed in folklore by the influence of the full moon on the werewolf; cf.~in English the word \emph{lunacy} (from Latin \emph{lūna} ‘moon’).  It is known that women are particularly affected by the moon through the synchronisation of the monthly and menstrual cycles (the two in fact being doublets, “menstrual” deriving from Latin \emph{mēnsis} ‘month’, originally ‘moon’) likely caused by the moon’s gravitational pull on the body, as proven by modern science \parencite{Forster2025}.}}; &
\alst{st}jǫrnur þat né vissu, \hld\ hvar þę́r \alst{st}aði ǫ́ttu}{\lemma{Mȧni \dots\ ǫ́ttu ‘the Moon \dots\ they had.’}\Afootnote{These lines are reversed (i.e., Sun, Stars, Moon) in \Regius\Hauksbok.}}.\eva

\bvb The Sun cast from the south---the Moon’s companion--- \\
her right hand over heaven’s rim. \\
The Sun knew not where halls she had; \\
the Moon knew not what sort of might he had; \\
the Stars knew not where seats they had.\evb\evg


\bvg\bva\mssnote{\Regius~1r/13, \Hauksbok~20r/9}%
\Ballnote{The Gods assemble at their Thing and solve the problem of the disorder of the heavenly bodies, implicitly by causing them to turn around in the Heavens (cf.~\textlink{Vafthrudnismal}[22]--25, \textlink{Grimnismal}[38]).  Thereafter they establish the reckoning of time by giving names to the natural cycles.  This name-giving, the establishment of norms, is seen as a part of the Godly ordering of the universe done to the weal of mankind \parencite[737]{Nordberg2020Time}; in fact it is the very act of naming and thereby categorising these phenomena which sets them apart from the disordered Chaos.  Thus the categories and norms, rather than being arbitrarily invented by Man, are seen as stemming directly from the Gods \parencite[21--22]{EliadeReturn}.
This view is not exactly identical to the modern, positively-oriented Christian scientific attitude towards Nature, which is based on the idea of understanding the order in the perfect design of an intelligent Creator.  In the Indo-European worldview, the raw matter already exists before the Gods in the disordered form of Yimer and other Ettins (cf.~\textlink{Vafthrudnismal}[31]).  The Creation consists in \emph{shaping} it according to the Godly Will, that is, the heroic masculine creative intellect.}%
\edtext{Þȧ gingu \edtrans{\alst{r}ęgin ǫll}{all the Reins}{\Bfootnote{All the Gods (\emph{ręgin} ‘Reins’, i.e., ‘Powers, Counsels’, a pantheistic word), here acting as one body.  According to \Skaldskaparmal\ 1 twelve Gods were present at the Thing, apart from Weden; these were Thunder, Nearth, Free, Tew, Homedall, Bray, Wider, Wonnel, Wolder, Heener, Foresitter, and Lock.  However, Wonnel and Lock can hardly have been seated at the Thing at the same time, since Wonnel was born to avenge Balder and Lock was bound soon after Balder’s death (cf.~st.~34 below, \textlink{Lokasenna}[P8]).  After Balder’s death, the Gods at the Thing thus numbered twelve, including Weden, as attested in the late \GautreksSaga\ 7.
The number twelve clearly belongs to historical Germanic Thing-proceedings and is the origin of the English twelve-man jury; cf.~\textlink{Baldrsdraumar}[1]/1--3 n.}} \hld\ ȧ \edtrans{\alst{r}ǫk-stóla}{council-seats}{\Bfootnote{This word is only found in the present refrain but corresponds to the \emph{sę́ti} ‘seats’ (\Gylfaginning~14:2, 4, 11) and \emph{dóm-stólar} ‘judgment-seats’ (\Gylfaginning~42:18, \Skaldskaparmal\ 43) in Snorre’s prose.}}, &
\edtrans{\alst{g}inn-hęilǫg \alst{g}oð}{yin-holy Gods}{\Bfootnote{A reverent epithet, ‘the vast holy (i.e., inviolable) Gods’.  The adj.~\emph{ginn-hęilagr} ‘yin-holy’ is composed of the augmentative prefix \emph{ginn-} ‘yin-’ (cf.~OE \emph{ġinn} ‘wide, vast, ample’) + \emph{hęilagr} ‘holy, inviolable’.
Outside of this refrain, \emph{ginn-hęilagr} only has one other occurrence in the Norse corpus, \textlink{Lokasenna}[11]/2: \emph{ǫll ginn-hęilǫg goð} ‘all yin-holy Gods’.}}, \hld\ ok umb þat \alst{g}ę̇ttu⸗sk}{\lemma{Þȧ \dots\ gę̇ttu⸗sk ‘Then \dots\ of that.’}\Bfootnote{The first refrain (\emph{stęf}) of the poem, describing the convening of the \inx[P]{Thing of the Gods} repeated below in sts.~9/1--2, 22/1--2, and 24/1--2.
The establishment of the Thing of the Gods is further described in \Gylfaginning~14:2--7 as cited in st.~7 below; the Thing is also mentioned in st.~46/2 below, \textlink{Hymiskvida}[39]/1, \textlink{Grimnismal}[29]--30, \textlink{Baldrsdraumar}[1]/1--3, \textlink{Thrymskvida}[14]/1--3, and \Haustlong\ 10/3.

A similar group of lines describing the convening of the Thing is found in \textlink{Baldrsdraumar}[1]/1--3 and \textlink{Thrymskvida}[14]/1--3, which follow the same structure as the present refrain.
In every occurrence of either refrain outside of the present stanza, the demonstrative pronoun \emph{þat} ‘this, that’ clearly refers to an immediately following question introduced by a \emph{hv}-word, which describes a problem the Gods have to deal with (e.g. \textlink{Thrymskvida}[14]/3--4: ‘and of this counseled the mighty Tews: how they Loride’s hammer might find.’)
Following this pattern, we would expect to find such a question following \emph{umb þat gę̇ttu⸗sk} ‘took counsel of that’ in the present stanza, and so it seems reasonably plausible (though not certain) that one has been lost in transmission.}}. &
\edtrans{\alst{N}ǫ́tt ok \alst{n}iðjum}{night and the moon-phases}{\Bfootnote{Also found in \textlink{Vafthrudnismal}[24]/4.}} \hld\ \alst{n}ǫfn of gǫ́fu, &
\alst{m}orgin hétu \hld\ ok \alst{m}iðjan dag, &
\alst{u}ndurn ok \alst{a}ptan, \hld\ \edtrans{\alst{ǫ́}rum at tęlja}{the years for to tally}{\Bfootnote{Cf.~\textlink{Vafthrudnismal}[22]--25, where it is said that the Gods made the sun and moon turn round in heaven and created the moon-phases \emph{ǫldum at ár-tali} ‘for mankind’s tally of years’.}}.\eva

\bvb Then went all the Reins onto their council-seats: \\
the yin-holy Gods, and from each other took counsel of that. \\
To night and the moon-phases names they gave; \\
morning they named, and middle day, \\
afternoon and evening the years for to tally.\evb\evg


\bvg\bva\mssnote{\Regius~1r/16, \Hauksbok~20r/10}%
\Ballnote{The Gods build their settlement. ---

The whole st.~is clearly a source for \Gylfaginning~14:2--7: \emph{Hárr mę́lti: „Í upp·hafi setti hann stjórnar-menn í sę́ti ok beiddi þá at dǿma með sér ør·lǫg manna ok ráða um skipun borgar’innar.  Þat var þar, sem heitir Iða-vǫllr í miðri borg’inni.  Var þat hit fyrsta þeira verk at gera hof þat, er sę́ti þeira tólf standa í ǫnnur en há-sę́tit, þat er Al·fǫðr á.  Þat hús er betst gert á jǫrðu ok mest.  Allt er þat útan ok innan svá sem gull eitt.  Í þeim stað kalla menn Glaðs-heim.  Annan sal gerðu þeir; þat var hǫrgr, er gyðjur’nar áttu, ok var hann all-fagr.  Þat hús kalla menn Vin-gólf.  Þar nę́st gerðu þeir hús, er þeir lǫgðu afla í, ok þar til gerðu þeir hamar ok tǫng ok steðja ok þaðan af ǫll tól ǫnnur.}

‘High spoke: “In the beginning He [= Weden, the Allfather] set men of authority in seats and invited them to judge alongside Him over the orlays of men and rule the arrangements of the stronghold.  It [i.e., the judgment-place] was at the place called the Idewold in the middle of the stronghold.  It was the first of their works to make the hove where their twelve seats stand, apart from the high seat which Allfather owns.  That house is the best one made on the earth and the greatest.  It is all without and within as if it were of gold alone.  That place men call Gladsham.  Another hall they made; it was a harrow, which the giddens owned, and it was all-fair.  That house men call Wingolf.  Thereafter they made houses wherein they laid hearths, and thereto they made a hammer and tongs and an anvil and therewith all other tools.’}%
Hittu⸗sk \alst{ę̇}sir \hld\ ȧ \edtrans{\alst{I}ða-vęlli}{the Idewold}{\Bfootnote{Dat.~sg.~of \emph{Iða-vǫllr} ‘the Plain of Industry.’}}, &
\edtrans{þęir’s \alst{h}ǫrg ok \alst{h}of \hld\ \alst{h}ǫ́-timbruðu}{they who timbered on high harrow and hove}{\Afootnote{\emph{afls kostuðu \hld\ alls freistuðu} ‘they exerted strength, they tried everything’ \Hauksbok}\Bfootnote{Two formulæ. --- \emph{hǫrgr ok hof} ‘harrow and hove’ is a merism, i.e., ritual structures made of stone and wood; cf.~\textlink{Vafthrudnismal}[38] and \textlink{HelgakvidaHjorvardssonar}[4], as well as the Saw of Anlaf Trueson in \FlatMS, which says that King Anlaf \emph{lét [...] brenna hof en brjóta hǫrga ok reisa í stað’inn kirkjur} ‘had hoves burnt and harrows broken and churches raised in their stead’. ---%TODO: FlatMS
\emph{hǫ́-timbra} ‘timber on high’ is a rare compound.  Its only other occurrence in the ON corpus is in \textlink{Grimnismal}[16] where it describes a harrow ruled by Nearth. --- This passage has often been wondered at.  Why would the Gods themselves make cultic buildings?  Whom do \emph{they} worship?
Let it be added to the present instance that they are also said to partake in the ritual sacrifice of animals, divining with their blood and feasting on their flesh (e.g. \textlink{Voluspa}[61], \textlink{Hymiskvida}[1], 39, \textlink{Lokasenna}, \Haustlong\ 2).
This question is really not as difficult as it is made out to be, however, since the answer must surely be that the behaviour of the Gods serves as the model for virtuous human behaviour and justifies the ideology of settlement.  Colonising new lands, tilling the earth, building houses and enclosures and ritual structures, sacrificing animals---if these are virtuous for humanity, they must find their model among the Gods, for ”Man constructs according to an archetype” and all settlements reenact Creation \parencite[10, 76--78]{EliadeReturn}.
This does not mean that the Gods worship themselves (or anything else); the Gods operate on the level of the macrocosm, and it is from them that the Godly Will flows.  Humans connect with this Godly Will through imitating the Gods on the level of the microcosm by sacrifice and settlement; so the sacrifice of Yimer and the enclosing of “the famous Middenyard” are reenacted through the sacrifice of horses and oxen within ritual and settlement enclosures \parencite[\textit{passim}]{Lincoln1986}.  See further \textlink{Grimnismal}[41]--43 and notes.}}; &
\edtrans{\alst{a}fla}{Hearths}{\Bfootnote{More specifically forges, by which we should understand that the Gods established fire in their settlement.

The civilising role of fire is found in Norse sources in the description of Icelandic settlers carrying fire around their land-claims, so e.g.~\Landnamabok\ 3.6: \emph{En meðan Sę́·mundr fór eldi um land-nám sitt} ‘But while Seamund went with fire around his land-claim’, 5.3: \emph{Þat land fór Jǫrundr eldi ok lagði til hofs} ‘Earwind went with fire throughout that land and laid it under the hove’, \EyrbyggjaSaga\ 4: \emph{fór Þór·ólfr eldi um land-nám sitt} ‘Thurolf went with fire around his land-claim’.
In \Landnamabok\ 5.1 the practice is framed as a rule put in place by King Harold Fairhair to prevent Icelandic settlers claiming excessive tracts of land: \emph{engi skyldi víðara nema en hann mę́tti eldi yfir fara á degi með skipverjum sínum. Menn skyldu eld gera, þá er sól vę́ri í austri. Þar skyldi gera aðra reyki, svá at hvára sę́i frá ǫðrum, en þeir eldar, er gervir váru, þá er sól var í austri, skyldi brenna til nę́tr. Síðan skyldu þeir ganga til þess, er sól vę́ri í vestri, ok gera þar aðra elda.}
‘No one was to claim more [land] than that over which he might bring fire in a day with his shipmates.  Men were to kindle fire when the sun was in the east.  There they were to make other smoke-fires so that each fire could be seen from the other, but the fires which were made when the sun was in the east were to burn till night.  Thereafter they were to walk until the sun was in the west and there make other fires.’

It seems likely that this explanation of the fire-ritual as a mere rule set down to prevent the excessive claiming of land is a later historicised explanation developed after the original reason behind the use of fire had been forgotten.  The ancient roots of the practice is demonstrated by the account of the legendary Gotlandish civilising founder, Thelver, in \Gutasaga\ 1: \emph{Gut·land hitti fyrsti maðr þann, sum Þielvarr hít.  Þá var Gut·land so eliust, et þet dagum sank ok náttum var uppi.  En þann maðr kvam fyrsti eldi á land, ok síðan sank þet aldri.} ‘Gotland was first discovered by the man who was named Thelver.  At that time Gotland was so uncanny [?] that it sank during the days and was up during the nights.  But that man was the first who brought fire to the land, and thereafter it never sank.’
Here we see that the bringing of fire fills a role in exorcising the land and making it civilised, that is, making it conform with the natural order.
This use of fire in settlements has Vedic parallels \parencite[11]{EliadeReturn} and fire also features prominently in Norse ritual; cf.~\textlink{Grimnismal}[43]/2.}} lǫgðu, \hld\ \alst{au}ð smíðuðu, &
\alst{t}angir skópu \hld\ ok \alst{t}ól gęrðu.\eva

\bvb The Eese met each other on the \inx[P]{Idewold}, \\
they who timbered on high \inx[C]{harrow} and \inx[C]{hove}. \\
Hearths they laid, wealth they smithed, \\
tongs they shaped and tools they made.\evb\evg


\bvg\bva\mssnote{\Regius~1r/18, \Hauksbok~20r/12}%
\Ballnote{After having settled down, the Gods are merry and content without troubles until their joy is spoiled by the coming of Ettin-women. ---
The whole stanza is paraphrased in \Gylfaginning~14:8--10, which explicitly identifies this period as the Golden Age (cf.~l.~2 n.): \emph{Ok því nę́st smíðuðu þeir málm ok stein ok tré ok svá gnóg⸗liga þann málm, er gull heitir, at ǫll bús-gǫgn ok ǫll reiði-gǫgn hǫfðu þeir af gulli, ok er sú ǫld kǫlluð gull-aldr, áðr en spillti⸗sk af til-kvámu kvinna’nna; þę́r kómu ór Jǫtun-heimum.} ‘And thereafter they forged metal and stone and wood, and so abundantly [did they forge] that metal which is called gold that all their household tools and riding tools were golden.  And that age is called the golden age, before it was spoiled by the arrival of the women; they came out of the Ettinhomes.’}%
\edtrans{\alst{T}ęflðu}{played Tables}{\Bfootnote{3rd pl. pret. ind. of \emph{tęfla}, a causative verb derived from \emph{tafl} ‘Tables, board game’, an old borrowing from Latin \emph{tabula}.  ‘Tables’ is used as a cognate translation; the exact type of board game referred to is unimportant, although it is certainly a game of strategy (e.g. \emph{hnefa-tafl} or nine men’s morris) rather than one of chance.  For \emph{hnefa-tafl} (which held the role as the chief board game in Scandinavia before the introduction of chess in the C12th) and its relatives see \textcite[55--64]{Murray1952}.}} ï \alst{t}u̇ni, \hld\ \alst{t}ęitir vǫ́ru, &
\edtrans{\alst{v}as þęim \edtrans{\alst{v}éttir⸗gis}{nothing}{\Bfootnote{An archaic gen.~of \emph{vétt⸗ki} ‘nothing’, the \emph{-ir-} representing a fossilized \emph{i}-stem gen.~sg., as ON \emph{véttr} ‘thing’ comes from PN \emph{*wihtiʀ}.  The only other occurrence of this form is in the notably archaic old Icelandic Homily Book (ms.~Holm perg 15 4°, fol. 36v/30).}} \hld\ \alst{v}ant ór golli}{for them was nothing golden wanting}{\Bfootnote{Indeed, even the bricks they played with were of gold (cf.~st.~58).  The Golden Age at the initial fulfillment of Creation is found in many mythologies \parencite[9, 112, 122, 127, 132, 135--136, 149, 154]{EliadeReturn}.}}, &
und’s \edtext{\alst{þ}ríar kvǫ̇mu \hld\ \alst{þ}ursa męyjar}{\lemma{þríar \dots\ þursa męyjar ‘three maidens of Thurses’}\Bfootnote{These three maidens are never mentioned again (unless they are taken to be the three norns in st.~19, but they would then be introduced twice). It is possible that an additional stanza giving further information about them has been lost. If it originally existed, it was already absent in the version used for \Gylfaginning, since no additional information is found there.}}, &
\edtrans{\alst{ȧ}m-át⸗kar}{uncanny}{\Bfootnote{The word \emph{ȧm-átt⸗igr} has a clear association with supernatural ettins.  It occurs in four other places in \Regius: in \textlink{Grimnismal}[11], \textlink{Skirnismal}[10], and \textlink{HelgakvidaHjorvardssonar}[17] it is paired with \emph{jǫtunn} ‘ettin’, while in \textlink{HelgakvidaHjorvardssonar}[14] it is used by the daughter of an ettin to refer to a human hero, a clear inversion of its usual use.}} mjǫk, \hld\ ór \alst{Jǫ}tun-hęimum.\eva

\bvb They played \inx[C]{Tables} in the yard, merry were they, \\
for them was nothing golden wanting \\
until three maidens of \inx[P]{Thurses} came, \\
most uncanny, out of the \inx[P]{Ettinham}[Ettinhomes].\evb\evg


\bvg\bva\mssnote{\Regius~1r/20, \Hauksbok~20r/14, \\
\RegiusProse~4v/3, \Trajectinus~TODO, \\
\Wormianus~TODO, \Upsaliensis~TODO}%
\Ballnote{After the Golden Age is spoiled, the Gods must get their gold in some other way.  For this they need the Dwarfs, who are associated with precious metals and ore-veins, and make many great treasures for the Gods (e.g.~Glapner, \textlink{EddicFragments}[F8][VIII], and Shidebladner, \textlink{Grimnismal}[44]).  The conception of spontaneous generation is clearly also at play; according to the prevalent view in antiquity maggots and worms were thought to spontaneously generate in rotting flesh and earth, and the dwarfs were apparently likewise thought to have spontaneously generated in stone.  As much is stated in \Gylfaginning~14:11--14, which continues with its paraphrase:

\emph{Þar nę́st settu⸗sk goð’in upp í sę́ti sín ok réttu dóma sína ok minntu⸗sk, hvaðan dvergar hǫfðu kviknat í mold’inni ok niðri í jǫrðu’nni, svá sem maðkar í holdi.  Dvergar’nir hǫfðu skipat⸗sk fyrst ok tekit kviknun í holdi Ymis ok váru þá maðkar, en af at·kvę́ðum goða’nna urðu þeir vitandi mann-vits ok hǫfðu manns líki ok búa þó í jǫrðu ok í steinum.  Móð·sognir var ǿðstr ok annarr Durinn.}

‘Thereafter the Gods set themselves high in their seats and made their judgments and remembered whence the Dwarfs had come to life in the ground and down in the earth like maggots in flesh.  The dwarfs had first taken shape and come to life in Yimer’s flesh and were then maggots, but by the decrees of the Gods they became knowing of manwit and had the forms of a man, and yet they live in the earth and in stones.  Moodsowner was the first in rank, and second was Dorn.’

Immediately after this passage, \Gylfaginning~14:15 quotes the present st.~and st.~10/3--4 in sequence, introduced by the words \emph{svá segir í Vǫlu-spá:} ‘so it says in the Spae of the Wallow:’.}%
Þȧ gingu \alst{r}ęgin ǫll \hld\ ȧ \alst{r}ǫk-stóla, &
\alst{g}inn-hęilǫg \alst{g}oð, \hld\ ok umb þat \alst{g}ę̇ttu⸗sk: &
\edtrans{Hvęrr skyldi \alst{d}verga}{Who should \dots\ of dwarfs}{\Afootnote{so \Regius\Wormianus\Upsaliensis \emph{;} \emph{at skyldi dverga} ‘That they should \dots\ of dwarfs’ \RegiusProse\Trajectinus \emph{;} \emph{hverir skyldu dvergar} ‘Which Dwarfs should [shape the troops out of \dots]’ \Hauksbok}} \hld\ \edtrans{\alst{d}rótt}{the troop}%
{\Afootnote{so \GylfMS \emph{;} \emph{drotin} ‘the lord’ \Regius \emph{;} \emph{dróttir} ‘the troops’ \Hauksbok}} \edtext{of}%
{\Afootnote{so \GylfMS \emph{;} om. \Regius\Hauksbok}} \edtrans{skępja}{shape}{\Afootnote{\emph{spękja} (norm.) ‘soothe’ \Upsaliensis}} &
\edtrans{ór \edtrans{\alst{B}rimis \alst{b}lóði}{Brimmer’s blood}{\Afootnote{so \Regius\Trajectinus \emph{;} \emph{brimi blóðgu} ‘the bloody surf’ \Hauksbok\RegiusProse\Wormianus\Upsaliensis.}\Bfootnote{Brimmer is perhaps a name for an ettin in st.~36/4 below, in which case it should be taken as a generic name here, i.e., ‘out of the ettin’s blood’.}} %
\hld\ ok ór \edtrans{\alst{B}láins}{Blown’s}{\Afootnote{so \Hauksbok\Wormianus \emph{;} \emph{†bláms†} (corrupted from \emph{*Bláins}) \RegiusProse\Trajectinus\Upsaliensis \emph{;} \emph{blǫ́um} \emph{(‘blám’)} ‘blue-black’ \Regius.}\Bfootnote{\emph{Bláinn} ‘Blown’ (from \emph{blár} ‘blue, black’) has usually been read as a poetic name for the ettin Yimer, and this is assumed here, although it is not attested as a name for him anywhere else.
It is, however, attested as a dwarf-name in Þul \emph{Dverga} 1 (\Skp\ 3).  Following Gurevich (\Skp\ 3, p. 693), we may therefore take the b-verse as a kenning ‘the legs of Blown \name{dwarf} \ken{stones}’, but this suffers from the fact that the dwarfs have not yet been created at this point in the narrative.}} lęggjum}{out of Brimmer’s blood and out of Blown’s legs}{\Bfootnote{The adopted reading for this line is a composite, but is based on the distributions of variants across the mss.~and the principle that the more difficult reading is the likelier.  Both \emph{brimi blóðgu} ‘the bloody surf’ and \emph{blǫ́um} \emph{(‘blám’)} ‘blue-black’ are thus to be taken as instances of trivialisation due to the confusion of rarely attested proper names \parencites[196--197]{Males2023Textual}[16--17]{MalesMyrvollForthcoming}.

The blood and legs are taken to be those of Yimer.  From his bones were made the rocks and from his blood the sea (\textlink{Grimnismal}[41], \textlink{Vafthrudnismal}[21]).
Explaining the creation of dwarfs from Yimer’s legs is easy, for dwarfs dwell in rocks (cf.~e.g.~\Ynglingatal\ 2, where the Swedish king Swayther (ON \emph{Svęigðir}) runs into a rock in pursuit of a dwarf).
More difficult to explain is how the dwarfs would be made out of the sea.  One suggestion is that this may be referring to the formation of salt-stones by means of evaporating salty sea-water.  Another interpretation is that we are dealing with the motif of underwater stones which may have been thought to have spontaneously generated in the sea; cf.~Scaldic \textsc{rock}-kennings like \Haleygjatal\ 2: \emph{sę́var bęin} ‘bone of the sea’ and \Thorsdrapa\ 16: \emph{fjarð-ępli} ‘firth-apple’ \parencite[90]{Meissner1921}.}}?\eva

\bvb Then went all the Reins onto their council-seats: \\
the yin-holy Gods, and from each other took counsel of this: \\
Who should shape the troop of \inx[P]{Dwarfs} \\
out of Brimmer’s blood and out of Blown’s legs?\evb\evg


\bvg\bva\mssnote{\Regius~1r/21, \Hauksbok~20r/15, \\
\RegiusProse~4v/5, \Trajectinus~TODO, \\
\Wormianus~TODO, \Upsaliensis~TODO}%
\Ballnote{Ll. 3--4 are cited in \Gylfaginning~14:15, immediately following st.~9 (see there).}%
\edtext{\edtext{Þar vas \alst{M}óð·sognir}{\Afootnote{so \Hauksbok \emph{,} (in the paraphrase) \GylfMS\ \emph{;} \emph{Þar Mót·sognir vitnir} ‘there Mootsowner wolf(?)’ \Regius.}} \hld\ \alst{m}ę́tstr of orðinn &
\alst{d}verga allra, \hld\ ęn \alst{D}urinn annarr;}{\lemma{Þar \dots\ annarr ‘There \dots\ second’}\Afootnote{om. \GylfMS}{\Bfootnote{These lines are not found in \Gylfaginning, but the author must have had the full stanza since \Gylfaginning~14:14 (“Moodsowner was the first in rank, and second was Dorn”, see st.~9) clearly derives from ll.~1--2.  They were likely omitted by the author of \Gylfaginning\ as superfluous.}}} &
þęir \edtrans{\alst{m}an-líkun}{man-likenesses}{\Bfootnote{The Dwarfs were in the shape of men, but not quite human.  In the creation of the Dwarfs before men, we may find a reflex of a motif common throughout world mythology: the creation of several failed or incomplete generations of sentient beings by the God(s) prior to the current race of humanity.  These failed or ancient generations are typically created out of other materials than those used for present humanity, display non-human features and uncanny abilities, and are either destroyed or turned into sprites or animals.
Mythological parallels include the five ages of man (viz. the Golden Age, Silver Age, Bronze Age, Heroic Age, and Iron Age) in Hesiod’s \emph{Works and Days} 109--201; the Jinn in Islam, created out of smokeless fire prior to the creation of Man out of clay (\Quran\ 15:26--27); and the several attempts at creating humans in the pre-Columbian Mayan \emph{Popol Vuh}, where the first generation was made out of earth and mud, but destroyed for its weakness (ll. 434--517), the second out of effigies of carved wood, becoming the spider monkeys (518--837), while the third and successful creation of humanity was out of maize (4822--4939) \parencite[64--75,178--181]{Christensen2007}.
The Norse Dwarfs fit surprisingly well into this typology.  They are created out of stone (whereas the first humans are made out of wood), are described in sources as unusually lustful and greedy, but also as masterful craftsmen, and are a form of sprite or earthly dæmon.}} \edtext{\hld\ \edtext{\alst{m}ǫrg of gørðu}{\lemma{þęir man-líkun \hld\ mǫrg of gørðu}\Afootnote{so \Regius\Hauksbok\Upsaliensis \emph{;}
\emph{þar man-líkun \hld\ mǫrg of gørðu⸗sk} \RegiusProse\Trajectinus\Wormianus}}, &
\alst{d}vergar \edtrans{ï}{in}{\Afootnote{so \Hauksbok\GylfMS \emph{;} \emph{ór} ‘out of’ \Regius}} jǫrðu, \hld\ \edtrans{sęm \alst{D}urinn sagði}{as Dorn said}%
{\Afootnote{so \Regius\Hauksbok\RegiusProse\Wormianus \emph{;} \emph{sem †dur menn† sagði} ‘as door-men(?) said’ \Trajectinus \emph{;} \emph{sem †þeim dyrinn kendi†} ‘as the beasts(?) taught them’ \Upsaliensis}}.}{\lemma{þęir \dots\ sagði ‘Those \dots\ said.’}%
\Bfootnote{The mss.~readings offer two conflicting versions of the creation of the dwarfs.
(1) \RegiusProse\Trajectinus\Wormianus: The dwarfs arise on their own through spontaneous generation; this is supported by the prose of \Gylfaginning~14:11--14 (see note to previous st.) and by the form of the stanza quoted there in the non-\Upsaliensis\ mss., but it may have been changed to correspond to the author’s vision (cf.~st.~38/3 n.~below).
(2) \Regius\Hauksbok\Upsaliensis: the dwarfs Moodsowner and Dorn shape “man-likenesses”, most likely lesser Dwarfs, out of the soil.  For stemmatic reasons, the present ed.~adopts (2).}}\eva

\bvb There was Moodsowner made the chiefest \\
of all Dwarfs, but Dorn the second. \\
Those Dwarfs made many man-likenesses \\
in the earth, as Dorn said.\evb\evg

\sectionline\setcounter{stanza}{15}

\begin{center}(For sts.~\textbf{11--15}---the \emph{Dwarf-tallies}---see Appendix to the poem.)\end{center}

\sectionline

\bvg\bva\mssnote{\Regius~1v/1, \Hauksbok~20r/26}%
\Ballnote{Sts.~16--17 are paraphrased in \Gylfaginning~9:3--7:

\emph{Þá er þeir gengu með sę́var-strǫndu Burs synir, fundu þeir tré tvau ok tóku upp tré’n ok skǫpuðu af menn.  Gaf inn fyrsti ǫnd ok líf, annarr vit ok hrǿring, þriði á·sjónu, mál ok heyrn ok sjón; gǫ́fu þeim klę́ði ok nǫfn.  Hét karl-maðr’inn Askr, en kona’n Embla, ok ólsk þaðan af mann-kind’in, sú er byggð’in var gefinn undir Mið·garði.}

‘When the sons of Byre (cf.~n.~to l.~2: \emph{ę̇sir}) were walking along the sea-shore they found two pieces of wood (\emph{tré}) and they took up the pieces of wood and shaped men out of them.  The first [god] gave breath and life; the second [gave] intelligence and movement; the third [gave] outward appearance, speech and hearing and sight.  They gave them clothes and names.  The man was called Ash and the woman Emble, and from them was begotten mankind, to whom the settlement within Middenyard was given.’

Based on the just-cited passage, the Gods are generally understood to have created mankind out of pieces of driftwood, and that is most likely what is being referred to by Snorre.  Yet no such detail is found in the present stanza; it may, therefore, instead be a later Icelandic or Snorronean reinterpretation.  As pointed out by \textcite{Hultgård2006}, the evidence from comparative mythology in fact suggests that the first humans were originally thought of as being made out of living, growing trees.

The descent of humans from trees is probably the reason behind the extensive use of words for trees in kennings for men and women in Norse poetry (see \Skp\ I, p. lxxv ff., \cite[245,266--272, 410]{Meissner1921}).  This kenning-type is most common in Scaldic poetry, but at times also appears in Eddic poetry (e.g. \textlink{Sigrdrifumal}[5]: \emph{bryn-þings apaldr} ‘apple-tree of the byrnie-\inx[C]{Thing} \ken{battle > warrior}’).}%
\edtrans{Und’s}{Until}{\Bfootnote{The first line, beginning with the conjunction \emph{und’s} ‘until’, is clearly supposed to connect syntactically to an immediately preceding stanza describing some ongoing action, as in st.~4/1 above: \emph{áðr} ‘before’.  For other instances of sts.~beginning with \emph{und’s} cf.~\textlink{Hymiskvida}[30], \textlink{Sigurdskamma}[3], 51, \textlink{GudrunTwo}[3]; in \Skp: \Haustlong\ 11, \Thorsdrapa\ 10, Anon \emph{Pl} 31.

If the dwarf-tallies (sts.~11--15) are later inserts, the preceding stanza would be st.~10, in which case the creation of men is probably to be understood as following and surpassing the creation of the dwarfs; cf.~note to st.~10/3: \emph{man-líkun}.}} %
\edtrans{\edtext{\alst{þ}r\emph{í}r}{\Afootnote{emend.; \emph{þrjár} \Regius\Hauksbok}} kvǫ̇mu \hld\ \edtext{ór \alst{þ}ví liði}{\Afootnote{so \Regius \emph{;} \emph{þussa brúðir} \Hauksbok.}}}%
{three came out of that host}{\Bfootnote{Both mss.~show corruption from st.~8 in their use of the f.~\emph{þrjár} ‘three’ for the m.~\emph{*þrír}.  \Hauksbok\ goes further in replacing the b-verse \emph{ór því liði} ‘out of that host’ with \emph{þussa brúðir} ‘brides of thurses’.  That this is nothing other than corruption is clearly shown by the m.~pl.~\emph{ǫfl⸗gir ok ǫ̇st⸗kir ę̇sir} ‘strong and loving Eese’ in l.~2.}} &
\edtrans{\alst{ǫ}fl⸗gir ok \edtrans{\alst{ǫ̇}st⸗kir}{loving}{\Bfootnote{The Wallow depicts the creation of Mankind as an act of love.  For our benefit, the Gods also created Middenyard (\textlink{Voluspa}[4]/2, \textlink{Grimnismal}[42]), the sun, and the moon with its phases (\textlink{Voluspa}[6], \textlink{Vafthrudnismal}[22]--25).}}}{strong and loving}{\Afootnote{\emph{ǫ̇st⸗kir ok ǫfl⸗gir} ‘loving and strong’ \Hauksbok}} \hld\ %
\edtrans{\alst{ę̇}sir}{Eese}{\Bfootnote{Weden, Heener, and Lother.  This stands in contradiction to \Gylfaginning~9:3--4 which tells us that the gods who made the first men were the sons of Byre, in \Gylfaginning~6:8--10 (cf.~st.4) identified with Weden, Will, and Wigh.  If we are to reconcile this difficulty we would have to see Heener and Lother as the equivalents of the latter two, but that identification is rather unlikely.

To begin with, not even the notoriously systematising \Gylfaginning\ equates Heener with Will or Wigh, nor does it mention Lother at all.  In Scaldic poetry the two gods Heener and Lother are found in a few kennings for Weden, but unlike kennings involving Will, these are always of the type ‘friend of Lother/Heener’ \ken*{= Weden} and never ‘brother of Lother/Heener’ \ken*{= Weden}.  Heener and Lother instead seem to be distinct gods.  See further the notes to their names in st.~17.}} \edtrans{at húsi}{along the settlement}{\Bfootnote{Lit.~‘along the house.’  An adverbial; the Gods were walking on the outskirts of their settlement.}}; &
fundu ȧ \alst{l}andi \hld\ \edtext{\alst{l}ítt megandi &
\edtrans{\alst{A}sk ok \alst{Ę}mblu}{Ash and Emble}{\Bfootnote{The ancestors of all human beings; cf.~Life and Lafthrasher who repopulate the world after its rebirth (\textlink{Vafthrudnismal}[44]--45).  Ash (nom. \emph{Askr}) is easily identified with the same-named wood species (\emph{Fraxinus excelsior}), but the form Emble (nom. \emph{Ęmbla}) is more difficult.  Her name is often translated as “Elm”, but the ON word for that tree (\emph{Ulmus glabra}) is the m.~\emph{almr}.
A possible etymology is metathesis from earlier \emph{*Ęlmla} (< PN \emph{*almilō} ‘young elm tree’), a f. \emph{n}-stem derivative of the same type as \emph{þęlla} ‘young fir tree’ (< PN \emph{*þallilō}) \char`~\ \emph{þǫll} ‘fir tree’ (< PN \emph{*þallu}), although the \emph{-b-} remains unexplained.}} \hld\ \alst{ø}r·lǫg-lausa}{\lemma{lítt megandi \dots\ ør·lǫg-lausa ‘little availing \dots\ orlay-less’}\Bfootnote{Prior to the gifts from the Gods the human beings were immobile and without \inx[C]{orlay} (fate, destiny), which is thus seen to be one of the distinguishing features of mankind.}}.\eva

\bvb Until three came out of that host, \\
strong and loving Eese along the settlement. \\
They found on the land the little availing \\
Ash and Emble, \inx[C]{orlay}-less.\evb\evg


\bvg\bva\mssnote{\Regius~1v/3, \Hauksbok~20r/27}%
\edtrans{\alst{Ǫ}nd}{Breath}{\Bfootnote{The breath of life or life force which sets living beings apart from dead matter.  The word \emph{ǫnd} originally means simply ‘breath’ (cf.~Gothic \emph{uz·anan} ‘to breathe out’), but the semantic development \textsc{breath} > \textsc{vitality} is by no means unparalleled, as seen by the Latin root-cognate \emph{animus} ‘life force, vitality; (later) soul’ vs. the Greek \textgreek{ᾰ̓́νεμος} ‘wind, breeze’.  All these words are ultimately from PIE \emph{*h₂enh₁-} ‘breathe’ \parencite[26--27]{Kroonen2013}.

After the conversion to Christianity, ON \emph{ǫnd} takes on the derived sense ‘spirit, soul’, as seen on Scandinavian memorial runestones where it is used interchangeably with the Anglo-Saxon borrowing \emph{sǫ́l} ‘soul’ (e.g. Sö 10 \emph{Guð hjalpi ǫnd hans} ‘God help his breath’ vs. Sö 8 \emph{Guð hjalpi sǫ́lu hans} ‘God help his soul’) and paired with it in at least 14 separate inscriptions (e.g. U 358 \emph{Guð hjalpi hans ǫnd ok sǫ́lu} ‘God help his breath and soul’).
This new semantic, probably calqued from the Latin \emph{spiritus} ‘breath; spirit’, is also charted in the Scaldic corpus, where \emph{ǫnd} (regardless of sense) makes it first securely dateable appearance in Hallfrith’s memorial poem to Anlaf Trueson (Hfr \emph{ErfÓl} 27 in \Skp\ 1) dated to ca. 1001 \ce.

The lack of pre-Christian Scaldic attestations of \emph{ǫnd} may lead one to question the authenticity of the sense ‘life (force)’ in the pagan period, but I think that it must have existed to begin with in order for the new sense of ‘spirit’ to have been acceptable to the Norse converts.  If \emph{ǫnd} had meant merely ‘respiration’ it would probably have been displaced entirely by the Christian loan-word \emph{sǫ́l} ‘soul’ rather than have seen acceptance even amongst former pagan poets like Hallfrith.
In any case the sense ‘life’ is attested in the most likely pre-Christian Eddic poem \textlink{Sigurdskamma}[33], 53 (\emph{ǫndu látit} ‘lost his life’), 60 (\emph{ǫndu tẏna} ‘id.’).
Also in later poetry a person who dies is often said to \emph{láta} or \emph{tẏna ǫndu} (cf.~\textlink{HelgakvidaHjorvardssonar}[37]; in \Skp\ Ótt \emph{Knútdr} 9, Þorm Lv 13, ÞSjár \emph{Þórdr} 3, Ív \emph{Sig} 30, ESk \emph{Geisl} 60, Bjbp \emph{Jóms} 41, Þórðh Lv 9, ǪrvOdd \emph{Ævdr} 12, 39, Anon \emph{Brúðv} 27), whereas it is not said once of the \emph{sǫ́l}; a Christian could hardly say that a person who dies “loses his soul”.

The use of \emph{ǫnd} in \Gylfaginning~3:10---printed at the beginning of this book for its simple beauty---is distinctly in the Christian sense ‘spirit’: \emph{Hitt er þó mest, er hann gerði mann’inn ok gaf honum ǫnd þá, er lifa skal ok aldri týna⸗sk, þó’tt líkam’inn fúni at moldu eða brenni at ǫsku} ‘Yet the greatest thing is when He [= Weden, the Allfather] made the man and gave him the “breath” which shall live and never perish although the body may moulder to dust or burn to ashes.’}} þau né \alst{ǫ́}ttu, \hld\ \edtrans{\alst{ó}ð}{wode}{\Bfootnote{Outside of the present st., ON \emph{óðr} ‘wode’ (gen.~sg.~\emph{óðar}) is used only in Scaldic poetry or in prose dealing with Scaldic poetry and then always with the sense ‘poem, poetry’.  This is most likely also its sense in the cpd.~\emph{Óð·rǿrir} ‘Woderearer’, one of the vats in which the Mead of Poetry was kept (\textlink{Havamal}[107]).

The present stanza thus gives us the word’s only Eddic attestations, but its exact sense in the present st.~is unlikely to be merely ‘poetry’ (which was after all given by Weden, not Heener, and that at a later time; see \textlink{Havamal}[103]--110).  To determine the proper sense of ‘wode’ we ought to consider its etymology.

As a m.~\emph{u}-stem the noun \emph{óðr} must derive from PGmc.~\emph{*wóþuz}.  PGmc.~words from the same root include the nouns \emph{*wódaz} ‘madness, fury’ (> OE \emph{wód} ‘madness’, OHG \emph{wuot} ‘fury’) and \emph{*wóþó} (> OE \emph{wóþ} ‘song’) and the adj.~\emph{*wódaz} ‘mad, furious’ (> ON \emph{óðr} ‘id.’, OE \emph{wód} ‘mad, rabid’, Gothic \emph{wóds} ‘possessed with the devil’).

Two Germanic theonyms are clearly derived from the root: Weden (ON \emph{Óðinn} < PGmc. \emph{*Wódanaz}) and Wode (ON \emph{Óðr} < PGmc. \emph{*Wódaz}).  Weden literally means ‘lord, possessor of wode’ (suffixed with PIE \emph{*-nós} ‘lord, master of’, cf.~his name \emph{Hęrjann} < PGmc. \emph{*harjanaz} < PIE \emph{kor-yo-nos} ‘lord of hosts, chieftain’ > Greek \textgreek{κοίρᾰνος} ‘king, ruler’), while Wode (unlike the noun \emph{óðr} an \emph{a}-stem as seen by \textlink{Voluspa}[24]) is probably a nominalisation of the adjective and thus means ‘mad, frenzied, inspired’.

The PIE root, which is found only in Germanic, Italic, and Celtic, is \emph{*wéh₂t-} and the alternation between the senses ‘poetry’ \char`~\ ‘fury’ is also seen in the other branches.  PGmc.~\emph{*wóþuz} derives from PIE \emph{*wéh₂t-u-s} (> OIr. \emph{fáth} ‘cause, motive, subject, prophecy, composition’, Old Welsh \emph{gwawt} ‘song of praise’).  Also related are Latin \emph{vātēs} ‘prophet, seer; poet’ and Gaulish \textgreek{ὀυάτεις} ‘a type of priest’ (pl., found in Strabo’s \emph{Geography} 4.4.4) \parencite[592]{Kroonen2013}.

The IE root sense, clearly retained in Celtic and Germanic, thus appears to be a state of frenzied poetic inspiration, \emph{furor poeticus}, but the sense in the present stanza may also cleave closer to ‘intelligent thought, mind’, or the power to express one’s thoughts with words, which is unique to man.}} þau né hǫfðu, &
\edtrans{\alst{l}ǫ́}{guile}{\Bfootnote{The exact sense of this word is uncertain.  It is otherwise found referring to sea-water near the shore, poetically to the ‘sea’, for which reason its sense here has sometimes been interpreted as ‘blood’.  \Skaldskaparmal\ 85 gives \emph{lǫ́} as a poetic synonym for ‘hair’, but it seems unlikely that hair would be mentioned next to the breath of life and intellect.
I tentatively take it to be an otherwise unattested variant of the poetic word \emph{lę́} ‘guile, deceit’, here in the sense ‘skill, craft, cunning’.
For the vowel alternation cf.~\emph{hrę́} ‘corpse’ (declined identically and rhyming with \emph{lę́}) < PGmc. \emph{*hraiwą}, irregularly metathesised from PIE \emph{*krow(h₂)-yóm} \parencite[242]{Kroonen2013} vs.~\emph{hrár} ‘raw, bloody’ < \emph{*hrawaz} < PIE \emph{*krow(h₂)-os} \parencite[244]{Kroonen2013}.}} né \alst{l}ę́ti \hld\ né \edtrans{\alst{l}itu góða}{goodly hues}{\Bfootnote{The ruddy complexion of the living in contrast to the pallor of the dead. ---
The reading \emph{litu goða} ‘hues of the Gods’ must be rejected for metrical reasons since the lift in \emph{goða} cannot be short \parencite[161]{Males2024Emend}; cf.~st.~5/2 n.~above.}}; &
\alst{ǫ}nd gaf \alst{Ó}ðinn, \hld\ \alst{ó}ð gaf \edtrans{Hø̇nir}{Heener}{\Bfootnote{An obscure god; he is one of the twelve Eese present at the \inx[C]{Thing of the Gods} (\Skaldskaparmal\ 1) and takes over the role as \emph{sacerdos civitatis} of the Gods after Weden dies, for which cf.~\textlink{Voluspa}[60] below.  His name is found thrice in Scaldic poetry, all in kennings for Lock of the type ‘\textsc{friend/companion} of Heener \ken*{= Lock}’ (\Haustlong~3/4b, 7/4a, 12/2a).

Heener plays a minor role in the mythology.  In \YnglingaSaga\ 4 he is one of the gods sent as a hostage from the Eese to the Wanes alongside Mimer, the wisest of the Eese; the Wanes in turn give over Quasher, their wisest man.
Heener is soon made a chieftain among the Wanes but is not very clever, and when present at the Thing without Mimer, the only thing he can say is \emph{ráði aðrir} ‘let others decide’.

Heener further figures in two adventures alongside Weden and Lock, viz.~in the encounter with Thedse (\Haustlong\ 1--13, \Skaldskaparmal\ 2--3) and in the origin of the Nivling hoard (\textlink{Reginsmal}[P1]--9, \Skaldskaparmal\ 46), but he plays no active role in either narrative.

\textcite[552]{Rydberg1886} convincingly argues that Heener is connected with the stork, deriving his name from the same root as Greek \textgreek{κύκνος} ‘swan’ and Sanskrit \emph{şakuná} ‘bird of omen’ (cf.~\cite[240]{Kroonen2013}), and noting that his epithets \emph{langi fótr} ‘long foot’ and \emph{aur-konungr} ‘mud-king’ (both found in \Skaldskaparmal\ 22) accurately describe that bird.
According to Rydberg, Heener’s giving of \inx[C]{wode} would then be associated with the later folk-myth of the stork that brings children to their parents, just as it is Heener who gives humans their mental faculties and thereby their unique human aspect.}}, &
\alst{l}ǫ́ gaf \edtrans{\alst{L}óðurr}{Lother}{\Bfootnote{An even more obscure god than Heener; his name is not found in the Poetic Edda and only twice in Scaldic poetry (\Haleygjatal\ 8/4a (\Skp\ 1) and HaukrV \emph{Ísldr} 1/1b (\Skp\ 4)), in both cases in the kenning \emph{vinr Lóðurs} ‘friend of Lother \ken*{= Weden}’.
Lother is possibly to be identified with the obscure figure \textbf{logaþore} (‘flame-darer(?)’) who appears alongside the two gods \textbf{wodan} (\emph{Wódan}, ‘Weden’) and \textbf{wigiþonar} (‘battle(?)-Thunder’) on the 6th c.~German Nordendorf fibula I (KJ 151) from near Augsburg, Bavaria, although this is uncertain \parencite[143]{TurvillePetre1975}.
Since the triad Weden, Heener, and Lock figures in other sources, Lother may also be identified with Lock, whose name may well be a nickname formed from the older Lother \parencite[143--144]{TurvillePetre1975}.  This identification is strongly supported if the above conjecture is right that Lother’s gift \emph{lǫ́} is an alternative form of \emph{lę́} ‘guile, deceit’, a word closely associated with Lock.}} \hld\ ok \alst{l}itu góða.\eva

\bvb Breath they had not; \inx[C]{wode} they had not, \\
not guile nor sound nor goodly hues. \\
Weden gave breath; Heener gave wode; \\
Lother gave guile and goodly hues.\evb\evg


\bvg\bva\mssnote{\Regius~1v/5, \Hauksbok~20r/29, \\
\RegiusProse~5r/30, \Trajectinus~TODO, \\
\Wormianus~TODO, \Upsaliensis~TODO}%
\Ballnote{This st.~is cited in \Gylfaginning~16:13, introduced by the words \emph{Svá sem hér segir} ‘Just as it says here’.  \Gylfaginning~16:10--12 provides context:

\emph{Enn er þat sagt, at nornir þę́r, er byggja við Urðar brunn, taka hvern dag vatn í brunni’num ok með aur’inn þann, er liggr um brunn’inn, ok ausa upp yfir ask’inn, til þess at eigi skuli limar hans tréna eða fúna. En þat vatn er svá heilagt, at allir hlutir, þeir er þar koma í brunn’inn, verða svá hvítir sem hinna sú, er skjall heitir, er innan liggr við egg-skurn.}

‘It is Furthermore said that the Norns who dwell by the Well of Weird fetch water in the Well every day, and along with it the mud which lies around the well, and they sprinkle it up over the Ash so that its branches should not harden or rot.  But that water is so holy that all things which come into the well turn as white as the membrane called the shell-membrane which lies inside the eggshell.’

Following the st.~\Gylfaginning~16:14--15 continues: \emph{Sú dǫgg, er þaðan af fellr á jǫrð’ina, þat kalla menn hunang-fall, ok þar af fǿðask bý-flugur. Fuglar tveir fǿðask í Urðar brunni. Þeir heita svanir, ok af þeim fuglum hefir komit þat fugla-kyn, er svá heitir.}
‘That dew which falls thereof onto the earth do men call ‘honey-fall’, and thereof bees feed themselves.  Two fowls feed themselves in the Well of Weird.  They are called ‘swans’, and from those fowls has come the race of fowls so called.’}%
\edtext{\alst{A}sk vęit’k \edtrans{standa}{standing}{\Afootnote{so \Regius\Hauksbok\Upsaliensis \emph{;} \emph{ausinn} ‘sprinkled’ \RegiusProse\Trajectinus\Wormianus}}, \hld\ hęitir \edtext{\alst{Y}gg·drasils}{\Afootnote{so \RegiusProse \emph{;} \emph{Ygg·drasill} all others}}}{\lemma{Ask \dots\ Ygg·drasils ‘Ash \dots\ Ugdrassle’s’}\Bfootnote{The Ash of Ugdrassle, the World Tree.
That the Ash stands in the very center of the world is shown by \textlink{Grimnismal}[26]/4, where we learn that all waters run their course from it, and \textlink{Grimnismal}[31], which tells us that its roots reach and wrap around the world of Men, the Underworld, and the Ettinhomes.}}, &
\alst{h}ǫ́r \edtrans{baðmr}{tree}{\Afootnote{\emph{borinn} ‘born’ (wo. doubt corrupt) \Upsaliensis.}}, \edtrans{\edtrans{ausinn}{sprinkled}{\Afootnote{\emph{hęilagr} ‘holy’ \GylfMS}} \hld\ \alst{h}víta auri}{sprinkled with white mud}{\Bfootnote{Cf.~st.~26/3--4 below.}}; &
þaðan koma \alst{d}ǫggvar \hld\ \edtext{þę́r’s}{\Afootnote{\emph{es} \RegiusProse\Trajectinus}} ï \edtext{\alst{d}ala}{\Afootnote{\emph{dali} \RegiusProse\Trajectinus\Upsaliensis}} falla; &
stęndr \edtext{\alst{ę́}}{\Afootnote{om. \Upsaliensis}} \alst{y}fir \edtext{grø̇nn}{\Afootnote{\emph{†grvnn†} \RegiusProse \emph{;} \emph{†grein†} \Upsaliensis}} \hld\ \alst{U}rðar brunni.\eva

\bvb An ash I know standing, it is called \inx[P]{Ugdrassle}’s: \\
a high tree sprinkled with white mud. \\
Thence come the dew-drops which fall in the dales; \\
it stands ever green over the \inx[P]{Well of Weird}.\evb\evg


\bvg\bva\mssnote{\Regius~1v/8, \Hauksbok~20r/31}%
\Ballnote{This st.~is paraphrased in \Gylfaginning~15:26--29: \emph{Þar stendr salr einn fagr undir aski’num við brunn’in, ok ór þeim sal koma þrjár meyjar, þę́r er svá heita: Urðr, Verðandi, Skuld.  Þessar meyjar skapa mǫnnum aldr; þę́r kǫllum vér nornir.} ‘There stands a single fair hall beneath the ash-tree by the well, and out of that hall come three maidens who are called thus: Weird, Werthing, Shild.  These maidens shape the ages of men; we call them Norns.’

For the three Norns, we may compare a strikingly similar description in the late ancient Gallo-Roman source \emph{De medicamentis} (TODO: cite), which relates a traditional Celtic myth of a tree standing in the middle of the sea, kept by three virgins.  A similar motif is found on numerous syncretic Gallo-Roman altars from western Germany, which depict the three \emph{Matronae Andrusteihae} ‘Tree-Standing Mothers’ next to a great tree (TODO: cite).
The three Norns also correspond closely to the Greek \textgreek{Μοῖραι}, the three weavers of fate (cf.~\textlink{HelgakvidaOne}[3]/1 n.).}%
Þaðan koma \alst{m}ęyjar \hld\ \alst{m}args vitandi &
\alst{þ}ríar ór þęim \edtrans{sal}{hall}{\Afootnote{so \Hauksbok \emph{,} (in the paraphrase) \GylfMS\ \emph{;} \emph{sę́} ‘lake’ \Regius}} \hld\ es \edtrans{und}{under}{\Afootnote{\emph{ȧ} ‘on’ \Hauksbok}} \edtrans{\alst{þ}olli}{tree}{\Bfootnote{Literally ‘fir’, but the word is only used for the alliteration and the species is generic.  The same may perhaps apply to \emph{askr} ‘ash’ above.}} stęndr; &
\edtrans{\alst{U}rð}{Weird}{\Bfootnote{Acc.~of \emph{Urðr} ‘fated course of events,’ lit.~‘that which has become’.  The cognates of this word feature prominently in a fatalistic context in West Germanic poetry (TODO).}} hétu \alst{ęi}na, \hld\ \alst{a}ðra \edtrans{Verðandi}{Werthing}{\Bfootnote{‘That which is becoming’, the pres.~participle of \emph{verða} ‘to become, happen’.  The name is only found here and in \Gylfaginning~15:27.}}, &
---\edtrans{\alst{sk}ǫ́ru ȧ \alst{sk}íði}{they scored billets}{\Bfootnote{\emph{skíð} refers to a thin piece of split wood, esp.~as used for firewood, but can also mean ‘plank’ or ‘ski’.  The sense here is most likely that the norns cut notches for the number of years allotted to each human being.
This notion is unparalleled (TODO: is it?) in mediæval literary sources, but the method of time keeping may represent a predecessor of the mediæval calendars used by the Northern European peasantry to track the lunar cycle and the Catholic feast-days.  These consist of carved wooden sticks where each day is represented either by a notch (in Norway and England) or by a rune of the futhark (in Sweden and Estonia) and go under a variety of names; \emph{primstav} in Norway, \emph{runstav} in Sweden, \emph{sirvilaud} in Estonia, and clog almanac in England.
Although these later wooden calendars are undoubtedly influenced by the timekeeping of the Church, the present line may just about hint that there existed a native, pre-Christian Germanic calendrical tradition.}}--- \hld\ \edtrans{\alst{Sk}uld}{Shild}{\Bfootnote{Acc.~of \emph{Skuld(r)} ‘that which shall be’.  The nom.~should probably have an \emph{-r}, though that is not attested in \Gylfaginning~15:27.  The name is only found here and in \Gylfaginning~15:27.}} ina þriðju &
þę́r \edtrans{\alst{l}ǫg}{laws}{\Bfootnote{With this word we should understand ‘decisions’, for pre-Christian Scandinavian Law was based on custom and precedent.  Similar legalistic vocabulary is found in the \emph{dómar} ‘judgments’ and \emph{kviðir} ‘decrees, rulings’ of the Norns (\textlink{Fafnismal}[11]/1, \textlink{Hamdismal}[30]/4).}} \alst{l}ǫgðu, \hld\ þę́r \alst{l}íf køru, &
\alst{a}lda bǫrnum, \hld\ \alst{ø}r·lǫg \edtrans{sęggja}{of youths}{\Afootnote{\emph{at sęgja} ‘to say’ \Hauksbok}}.\eva

\bvb Thence come maidens much-knowing: \\
three out of the hall which stands beneath the tree. \\
One they called Weird, the other Werthing \\
---they scored billets---Shild the third. \\
They laid down laws; they chose lives \\
for the children of men, the \inx[C]{orlay} of youths.\evb\evg


\bvg\bva\mssnote{\Regius~1v/11, \Hauksbok~20v/5}%
\Ballnote{A very cryptic stanza.  The name Goldwey is not mentioned in any other source, nor is the underlying story, but she was apparently slain, burned, and reborn three times.  Goldwey is perhaps to be identified with the wallow of the poem herself, which would explain the cryptic “yet she lives still”.}%
Þat man hǫ̇n \edtrans{\alst{f}olk-víg}{troop-conflict}{\Bfootnote{\emph{folk} here carries its older meaning ‘troop, band’, as seen e.g.~in the Slavic borrowing exemplified by Russian \textrussian{полк} ‘regiment, host, army’.}} \hld\ \alst{f}yrst ï hęimi, &
es \alst{G}oll·vęigu \hld\ \alst{g}ęirum studdu &
ok ï \alst{h}ǫll \edtrans{\edtext{\alst{H}á\emph{a}rs}{\Afootnote{\emph{Hárs} \Regius\Hauksbok}}}{Hougher}{\Bfootnote{Hiatus must be restored for metrical reasons; cf.~\textcite[265--266]{Males2023Sievers}.  \emph{Háarr}, usually normalised as \emph{Hǫ́arr}, is a name for Weden also found in Scaldic poetry, perhaps < \emph{*haiha-harjaʀ} ‘one-eyed warrior’.
It should not be confused with Weden’s names \emph{Hárr} ‘Hoary’ and \emph{Hǫ́r} ‘High’.}} \hld\ \edtext{\alst{h}ȧna}{\Bfootnote{The meter requires the pronoun to have a long syllable.}} bręnndu, &
\edtext{\alst{þ}rysvar bręnndu}{\Afootnote{written twice \Hauksbok}} \hld\ \alst{þ}rysvar borna, &
\alst{o}pt, \alst{ȯ}·sjaldan, \hld\ þó hǫ̇n \alst{ę}nn lifir.\eva

\bvb That troop-conflict she recalls first in the \inx[C]{Home} \\
when Goldwey they impaled with spears \\
and in the hall of \inx[P]{Hougher} \name{= Weden} \ken*{= Walhall} burned her. \\
Thrice they burned the one thrice born, \\
often, unseldom---yet she lives still.\evb\evg


\bvg\bva\mssnote{\Regius~1v/13, \Hauksbok~20v/7}%
\edtrans{\alst{H}ęiði}{Hode}{\Bfootnote{Acc.~sg.~of \emph{*Hęiðr} ‘bright, splendid’.}} \alst{h}étu, \hld\ hvar’s til \alst{h}úsa \edtext{kom,}{\Afootnote{add.~\emph{ok} \Hauksbok}} &
\alst{v}ǫlu \alst{v}ęl-spáa, \hld\ \alst{v}itti \edtrans{ganda}{gands}{\Bfootnote{Familiar spirits.  TODO: elaborate.}}; &
\alst{s}ęið hǫ̇n \edtext{hvar’s hǫ̇n}{\Afootnote{so \Hauksbok \emph{;} om.~\Regius}} kunni, \hld\ \alst{s}ęið hǫ̇n \edtext{hug}{\Afootnote{so \Hauksbok \emph{;} om.~\Regius}} lęikinn; &
\alst{ę́} vas hǫ̇n \alst{a}ngan \hld\ \alst{i}llrar brúðar.\eva

\bvb Hode they called her where to houses she came, \\
the well-spaeing \inx[C]{wallow}; she bewitched \inx[C]{gand}[gands]. \\
She ensorcelled where she could; she ensorcelled deluded minds; \\
she was ever the delight of an evil bride.\evb\evg


\bvg\bva\mssnote{\Regius~1v/16, \Hauksbok~20v/9}%
\Ballnote{Sts. 22--23 deal with the obscure war between the Eese and the Wanes, the two tribes of the Gods.  It is most fully described in \YnglingaSaga\ 4, which also features a long description of the hostages: \emph{Óðinn fór með her á hendr Vǫnum, en þeir urðu vel við ok vǫrðu land sitt, ok hǫfðu ymsir sigr; herjuðu hvárir á land annarra ok gerðu skaða.  En er þat leiddi⸗sk hvárum-tveggjum, lǫgðu þeir milli sín séttar-stefnu, ok gerðu frið ok seldu⸗sk gíslar.} ‘Weden went with a host against the Wanes, but they responded well and defended their land, and both sides had alternating victories; either side raided the land of the other and did harm.  But when that seemed loathsome to both sides, they held a conciliatory settlement between them and made peace and traded hostages.’

The war is also alluded to in \Gylfaginning~23:6 (edited in the pres.~ed.~under \textlink{EddicFragments}[F2][II]), which mentions how Nearth was one of the hostages sent by the Wanes, and \Skaldskaparmal\ 5, which explains the origin of the mead of poetry (for which cf.~\textlink{Havamal}[103]--110): \emph{Þat váru upp-hǫf til þess, at goðin hǫfðu ó·sę́tt við þat folk, er Vanir heita. En þeir lǫgðu með sér frið-stefnu ok settu grið á þá lund, at þeir gengu hvárir-tveggju til eins kers ok spýttu í hráka sínum} ‘This was the origin of it, that the Gods had hostile relations with the folk called the Wanes.  But they held a peace-settlement amongst each other and made a truce by the following means, that each of them went to a vat and spit their spittle into it.’

As the relevant sources (the Prose Edda and \YnglingaSaga) are closely related and probably both written by the same author, Snorre, the Eese-Wane war and indeed the very existence of the Eese and the Wanes as distinct tribes (rather than synonyms for the Gods as a collective) has been quite powerfully argued by \textcite{Simek2010} to be a Snorronean learned concoction as part of his learned pre-history of Scandinavia.  This is supported by data from \textcite{FrogRoper2011}, showing that the ON term \emph{vanir} ‘Wanes’ is used only in Eddic poetry and exclusively in alliterative positions.
Even so, a straightforward reading of \Voluspa~33/3--4 does see the ‘Wanes’ breaking down the fortification of the ‘Eese’, and when compared with \textlink{Vafthrudnismal}[38]--39 which explicitly contrasts the ‘Wanes’ and the ‘Eese’ they do in fact seem to be distinct groups also in the Poetic Edda.}%
Þȧ gingu \alst{r}ęgin ǫll \hld\ ȧ \alst{r}ǫk-stóla, &
\alst{g}inn-hęilǫg goð, \hld\ ok umb þat \alst{g}ę̇ttu⸗sk: &
Hvárt skyldu \alst{ę̇}sir \hld\ \alst{a}f·ráð gjalda, &
eða skyldu \edtrans{\alst{g}oð’in ǫll}{all the Gods}{\Bfootnote{The clitic definite \emph{-in} is very rare in older ON poetry.  This is its only occurrence in \textlink{Voluspa}, but it should be understood in the context of preceding the adjective \emph{ǫll}, for which cf.~\textlink{Lokasenna}[52]/4. --- Here “all the Gods” (viz., the Eese \emph{and} the Wanes) seem to be contrasted with the \inx[P]{Eese}, a subset.}} \hld\ \alst{g}ildi ęiga?\eva

\bvb Then went all the Reins onto their council-seats: \\
the yin-holy Gods, and from each other took counsel of this: \\
Whether the Eese should yield tribute \\
or all the Gods should hold a banquet?\evb\evg


\bvg\bva\mssnote{\Regius~1v/17, \Hauksbok~20v/11}%
\edtrans{\alst{F}lęygði Óðinn \hld\ ok ï \alst{f}olk of skaut}{Weden hurled and shot into the troop}{\Bfootnote{The object, a spear, is understood. --- This first spear-throw was ritually reenacted during battle, so \EyrbyggjaSaga\ 44: \emph{þá skaut Stein·þórr spjóti at fornum sið til heilla sér yfir flokk Snorra} ‘then Stanthur shot a spear according to ancient custom for his fortune over Snorre’s troop’ and \StyrbjarnarThattr\ 2, where the Swedish king Eric devotes himself as a sacrifice within a period of ten years in exchange for victory.  In return, Weden gives him a reed and tells him to hurl it into the opposing host with the phrase \emph{Óðinn á yðr alla} ‘Weden owns you all!’  Eric does so, and the reed turns into a spear in the air (for this motif cf.~\GautreksSaga\ 7 as summarized in \textlink{Havamal}[138] n.).  As it flies over his foes they are stricken by terror and blindness before being buried in a landslide.

It was customary among the Germanic pagans that the battle-slain (\emph{valr}) were devoted to Weden; cf.~Þhorn \emph{Harkv} 12 (\Skp\ 1): \emph{Valr lá þar ȧ sandi \hld\ vitinn inum ęin-ęygja // Friggjar faðm-byggvi} ‘The slain lay there on the sand set apart for the one-eyed bosom-dweller \ken{husband} of Frie \ken*{= Weden}’.
The consecration could also be done as a thank-offering after a victory, as in \HervararSaga\ 7: \emph{Heið·rekr kveð⸗sk nú gjalda fyrir son sinn þetta lið allt, er drepit var, ok gaf hann nú þennan val Óðni} ‘Heathric now declares that he, in exchange for his son[’s life], pays up the whole part of the army which was killed, and he now gave these battle-slain to Weden.’
In the afterlife these men would join the god as \inx[P]{Oneharriers} in \inx[P]{Walhall} (cf.~\textlink{Grimnismal}[14], \textlink{Harbardsljod}[24], runic inscription N B380 edited below under Galders).  The sacrifice also included captives, which were at times executed by hanging (\textlink{Havamal}[138] n.) and which may be the reason why Weden is also called \emph{Hapta·guð} ‘God of Captives’ (\Gylfaginning~20:10).  As seen below, even the weapons of the vanquished were part of the sacrifice.

This type of devotion of an entire army was not uncommon in Germanic antiquity and is securely attested in both history and archeology.  Among historical sources I particularly wish to mention Orosius, \emph{History} 5.16.5--6, describing the battle of Arausio fought in 105 \bce\ between the Romans and the Germanic tribes of the Cimbri and Teutons: “After capturing the two Roman camps and a vast amount of booty, the enemy destroyed everything that they had laid their hands on in some new, unexpected form of curse.  6 Clothing was ripped up and discarded, gold and silver thrown into the river, the men’s armour was torn apart, the horses’ harness scattered and the horses themselves drowned in the river, while the men had nooses tied round their necks and were hanged from trees. In this way the victor knew no booty nor the vanquished any mercy.” (\cite[235]{FearOrosius}).

The most relevant archeological finds are the various Northern European bogs and wetlands where horses and arms, presumably belonging to defeated foes, have been mutilated and defaced (e.g. by whipping the horses, bending the swords) before their submersion.  This defacing is to be explained by the fact that the objects, having been devoted and made holy unto the god, had to be made materially worthless and unusable to men, thereby releasing their inherent worldly power.  This is confirmed by Tacitus, \emph{Annals} 13.57 who describes how the Chatti, \emph{victores diversam aciem Marti ac Mercurio sacravere, quo voto equi, viri, cuncta occidioni dantur} ‘had devoted, in the event of victory, the enemy’s army to Mars and Mercury, a vow which consigns horses, men, everything indeed (on the vanquished side) to destruction.’
For Germanic devotional sacrifice cf.~further \textcite[617]{Schjodt2020VariousWays}, \textcite[1166--1168]{Schjodt2020Odin}, for bog burials cf.~\textlink{Voluspa}[38] n. below, for human sacrifice to Weden generally cf.~\textlink{Havamal}[138] n.

Outside of the Germanic sphere similar customs of holy war were practiced among the Romans (Latin \emph{devotio}), as told in cf.~Livy 8.9--10 where the Roman consul of 340 \bce, Publius Decius Mus, (like king Eric in \StyrbjarnarThattr) devotes not just the opposing legions but also himself, although Livy comments that this self-devotion was not strictly necessary for the ritual.
Similarly, the Jews, as part of their martial rites (Hebrew \texthebrew{חֵרֶם} \emph{ḥērem}; \cite{Lohfink1986}), would devote entire cities or nations to YHWH (Deuteronomy 20:16--17, Joshua 6:17--21, Joshua 8:1--29; note in Joshua 8:18 the detail of the god-given spear that leads to victory), after which, like in the Germanic bog deposits, also the arms and property of the vanquished were to be destroyed or set aside for the god (Leviticus 27:28--29).}}; &
þat vas ęnn \alst{f}olk-víg \hld\ \edtrans{\alst{f}yrr}{earlier}{\Afootnote{so \Hauksbok \emph{;} \emph{fyrst} ‘first’ \Regius}\Bfootnote{The \Regius\ reading cannot be correct as this st.~is describing a different war than the first; it has probably arisen due to the similarity with st.~20/1.}} ï hęimi; &
\alst{b}rotinn vas \edtrans{\alst{b}orð-vęggr}{plank-wall}{\Bfootnote{A dis legomenon in the Old Norse corpus; its only other occurrence is in Þloft \emph{Glækv} 6 (\Skp\ 1) from ca.~1032, where it describes some kind of wooden structure in which church-bells ring, probably a part of a stave-church.  I take it to refer to the wooden fortification of the stronghold of the Eese, i.e., Osyard.  Cf.~note to st.~24.}} \hld\ \alst{b}orgar ȧsa, &
knǫ́ttu \alst{v}anir \edtrans{\alst{v}íg-spǫ́}{war-spae}{\Bfootnote{The Wanes used a magic prophecy (\emph{spǫ́} ‘spae’) to win the battle and sack \inx[P]{Osyard}, the stronghold of the Eese.}} \hld\ \alst{v}ǫllu sporna.\eva

\bvb Weden hurled and shot into the troop; \\
that was yet a troop-conflict earlier in the \inx[P]{Home}. \\
Broken was the plank-wall of the stronghold of the Eese; \\
the Wanes by a war-\inx[C]{spae} did tread the fields.\evb\evg


\bvg\bva\mssnote{\Regius~1v/19, \Hauksbok~20r/34, \\
\RegiusProse~10v/25, \Trajectinus~TODO, \\
\Wormianus~TODO, \Upsaliensis~TODO}%
\Ballnote{After the wooden stronghold of the Eese, protected only by a plank-wall (\emph{borð-vęggr}), has been sacked by the Wanes, they decide to build stronger fortifications.  This leads into the story of the Master Builder as told at length in the following way in \Gylfaginning~42, which concludes by citing \Voluspa\ 24--25:
An \inx[P]{Ettins}[ettin] craftsman arrives at Osyard and offers to build the Eese a great wall in exchange for \inx[P]{Frow}’s hand and the \inx[P]{Sun} and \inx[P]{Moon}.  He further promises that he will finish the entire wall alone over a single winter, or else the contract will be void.  This seems reasonable to the Eese, who assume that he cannot possibly finish it in time, but things go awry when the ettin asks for permission to have the help of his workhorse \inx[P]{Swaddlefare}.  Lock grants him this without first consulting the other Gods, after which the agreement is sealed with strong oaths.  The horse turns out to be unexpectedly strong, and when only three days remain before summer, the wall is almost finished.  The panicked Eese blame Lock and threaten him with violence if he does not get rid of the horse.  His solution is to turn into a mare and distract it.  The trick works and the two stay out all night; Lock is made pregnant, later giving birth to \inx[P]{Slopner}.  When the ettin realises that he has been tricked he comes into his greatest ettin wrath, at which point the Eese call on Thunder who shortly appears and kills him.}%
Þȧ gingu \alst{r}ęgin ǫll \hld\ ȧ \alst{r}ǫk-stóla, &
\alst{g}inn-hęilǫg \alst{g}oð, \hld\ ok umb þat \alst{g}ę̇ttu⸗sk: &
Hvęrr hęfði \alst{l}opt alt \hld\ \alst{l}ę́vi blandit &
eða \alst{ę́}tt \alst{jǫ}tuns \hld\ \edtrans{\alst{Ó}ðs męy}{Wode’s maiden \ken*{= Frow}}{\Bfootnote{According to \Gylfaginning~35:7 and \Skaldskaparmal\ 28, 44 \emph{Óðr} ‘Wode’ (not to be confused with \emph{Óðinn} ‘Weden’) was Frow’s husband. ---
The ettins seem to have a particular obsession with abducting Frow, as seen also in \textlink{Thrymskvida}[8].  This may be explained if Frow is the Indo-European dawn-goddess, something suggested not only by her association with the Sun and Moon in the Master Builder myth, but also by her weeping of golden tears (\Gylfaginning~35:9, \Skaldskaparmal\ 40, 44; possibly a poetic description of the dawn), her epithet Meredall (\emph{Mar·dǫll} ‘Sea-brilliant’, another likely poetic description; for the second element see \textlink{Vafthrudnismal}[25]), and her possession of the fire-symbolic Blaze-Necklace (\emph{brísinga-męn}, cf.~\textlink{Thrymskvida}[13]/3 n.).

For the attempted theft of the sun and moon we may compare \emph{Kalevala} (songs 47--49), where the evil witch \emph{Louhi} of the Underworld steals the sun, moon, and fire and thereby deprives the world of its light; for further Finnic parallels cf.~the Introduction to \textlink{Thrymskvida}.
The attempted abduction of the Dawn has a more distant parallel in the Vedic \textoi{Valá}-myth (part of which is told in \Rigveda\ 10.108; \cite[70, 71]{JamBrer2020}).  In that myth the enemies of the Gods, the \textoi{Paṇi}-s, hide the Dawns in the form of kine in the \textoi{Valá}-cave and thereby prevent the natural cycle between Night and Day.  With his great song the Storm-god \textoi{Índra} breaks open the cave, killing \textoi{Valá} and releasing the dawns.
The common thread in these narratives is the fear that evil forces will steal away the light and warmth of the world, which in the Norse mythology reflects itself as the Fimblewinter, for which cf.~\textlink{Vafthrudnismal}[44]/3--4 n.}} gefna?\eva

\bvb Then went all the Reins onto their council-seats: \\
the yin-holy Gods, and from each other took counsel of this: \\
Who had blended all the air with deceit, \\
or to the ettin’s lineage given \inx[P]{Wode}’s maiden \ken*{= Frow}?\evb\evg


\bvg\bva\mssnote{\Regius~1v/20, \Hauksbok~20r/36, \\
\RegiusProse~10v/27, \Trajectinus~TODO, \\
\Wormianus~TODO, \Upsaliensis~TODO}%
\Aallnote{The order of the lines follows \Regius\Hauksbok \emph{;} the two helmings are reversed (i.e., 3--4, 1--2) in \GylfMS.}%
\alst{Þ}ȯrr ęinn \edtrans{\alst{þ}ar vá}{fought there}{\Afootnote{so \Hauksbok\Trajectinus\Upsaliensis \emph{;} \emph{þar var} ‘was there’ \Regius \emph{;} \emph{þat vann} ‘accomplished it’ \RegiusProse; \emph{þat vá} ‘fought it out’ \Wormianus}} \hld\ \alst{þ}runginn móði, &
\edtrans{hann \alst{s}jaldan \alst{s}itr \hld\ es \alst{s}·líkt of fregn;}{he seldom sits when of such he learns}{\Bfootnote{When he learns of an ettin encroaching on the Gods (see note to st.~24).  Thunder, as the defender of the Gods (\textlink{Thrymskvida}[18], Þdís \emph{Þórr} in \Skp\ 3), is willing to break even the strongest oaths if the safety of the Gods require it (cf.~\textlink{Lokasenna}[57]--64).}} &
\edtext{\alst{ȧ} gingu⸗sk \alst{ęi}ðar, \hld\ \alst{o}rð ok sǿri, &
\alst{m}ǫ́l ǫll \alst{m}ęgin⸗lig, \hld\ es ȧ \alst{m}eðal \edtext{fóru}{\lemma{fóru ‘had gone’}\Afootnote{\emph{vǫ́ru} ‘had been’ \Hauksbok\Trajectinus}}.}{\lemma{ȧ \dots\ fóru.}\Afootnote{om. \Wormianus}}\eva

\bvb Thunder alone fought there, pressed by wrath; \\
he seldom sits when of such he learns. \\
Trampled were oaths, speeches, and vows, \\
all the mighty treaties which between them had gone.\evb\evg


\bvg\bva\mssnote{\Regius~1v/23, \Hauksbok~20v/1}%
\Ballnote{The purpose of the stanza is quite obscure, but it appears to cryptically describe Mimer’s Well, into which Weden deposited his eye.

Cf.~\Gylfaginning~15:10--13: \emph{En undir þeiri rót, er til hrím-þursa horfir, þar er Mímis-brunnr, er spekð ok mann-vit er í folgit, ok heitir sá Mímir, er á brunn’inn.  Hann er fullr af vísendum, fyrir því at hann drekkr ór brunni’num af horni’nu Gjallar-horni. Þar kom Al·fǫðr ok beiddi⸗sk eins drykkjar af brunni’num, en hann fekk eigi, fyrr en hann lagði auga sitt at veði.} ‘But enwrapped by the root [of Ugdrassle] which turns towards the Rime-Thurses there is Mimer’s Well.  Wisdom and manwit are hidden within it and he who owns the well is called Mimer.  He is full of wise sayings, for he drinks from the Well out of the horn Yellhorn.  Thither Allfather came and asked for a single drink from the well, but He did not get it before He laid down his eye as a pledge.’

Immediately after this passage \Gylfaginning~15:14 quotes st.~28; cf.~there.}%
Vęit hǫ̇n \edtrans{\alst{H}ęim·dalar \hld\ \alst{h}ljóð}{Homedall’s sound}{\Bfootnote{Interpreted by Snorre (see note to ALL) as the Yellhorn, for which cf.~st.~45 below.}} of folgit &
und \edtrans{\alst{h}ęið-vǫnum}{shady}{\Bfootnote{Literally ‘light-less’, \emph{hęiðr} referring especially to the light of a clear sky.}} \hld\ \alst{h}ęlgum baðmi; &
\edtext{\alst{ȧ} sér hǫ̇n \alst{au}sa⸗sk \hld\ \edtrans{\alst{au}rgum}{muddy}{\Bfootnote{Most likely the same mud (\emph{aurr}) as in st.~19, there said of the Well of Weird.}} forsi &
af \edtrans{\alst{v}ęði \alst{V}al-fǫðrs}{Walfather’s pledge}{\Bfootnote{Weden’s eye, which he gave up in exchange for a sip from the Well.}}. \hld\ \edtrans{\alst{V}ituð ér ęnn eða hvat?}{Know ye yet, or what?}{\Bfootnote{I.e., ‘Do you, Weden, still know what I am saying, or what more do you want?’  The second refrain (\emph{stęf}) of the poem, repeated in 28, 33, 34, 38, 40, 46, 59, and 60.
Similar refrains are found in \textlink{Baldrsdraumar}[7]/4 et c., \textlink{Hyndluljod}[17]/4 et c.}}}{\lemma{ȧ \dots\ Val-fǫðrs ‘She sees \dots\ pledge’}\Bfootnote{The interpretation of these lines is ambiguous, as \emph{‘á’} \Regius, \emph{‘aa’} \Hauksbok\ can be read either (1) as the preposition \emph{ȧ} ‘on’ or (2) as the noun \emph{ǫ́} ‘river’.  In the latter case, the lines would mean “She sees a river being fed by a muddy torrent from Walfather’s pledge”.
Note that this ambiguity is not intentional, since the distinction between nasal and oral vowels and between long \emph{á} and \emph{ǫ́} was still upheld in the 12th century \parenciteshorthand{FGTHaugen}.  (1) is adopted here since it is clearly the interpretation preferred by Snorre, who, in his prose, has the Norns \emph{ausa} ‘sprinkle, pour’ \emph{aur’inn} ‘the mud’ over the ash (\Gylfaginning~16:10--11, cited in st.~19 n.~above).
Since Snorre’s version of the poem may theoretically have preserved the distinction, and since he was a native speaker of the language, his opinion is deferred to.}}\eva

\bvb She knows Homedall’s sound \ken*{= Yellhorn} hidden \\
beneath the shady, holy tree \ken*{= Ugdrassle’s Ash}. \\
She sees it being sprinkled with a muddy torrent \\
from Walfather’s pledge \ken*{= Mimer’s Well}. \emph{Know ye yet, or what?}\evb\evg


\bvg\bva\mssnote{\Regius~1v/25}%
\Ballnote{Stanzas 27--29 appear to relate how the wallow was summoned by Weden.}%
\edtrans{\alst{Ęi}n sat hǫ̇n \alst{ú}ti}{Alone she sat outside}{\Bfootnote{ON \emph{sitja úti} ‘sit outside’ has a cultural connotation of meditating in order to communicate with the otherworld; cf.~the noun \emph{úti-seta}. --- This line is directly repeated in \textlink{Sigurdskamma}[6]/1a.}}, \hld\ þȧ’s hinn \alst{a}ldni kom &
\alst{y}ggjungr \alst{ȧ}sa \hld\ ok ï \alst{au}gu lęit: &
\speakernote{[Vǫlva:]}%
‚Hvęrs \alst{f}regnið mik? \hld\ hví \edtrans{\alst{f}ręistið}{tempt}{\Bfootnote{\emph{fręista} ‘tempt’ has a sense of testing someone, especially intellectually.  Cf.~\textlink{Havamal}[26], \textlink{Vafthrudnismal}[3], 5.}} mïn?\eva

\bvb Alone she sat outside when the old one came, \\
the Frightener of the Eese \ken*{= Weden}, and looked into her eyes. \\
\speakernoteb{[The wallow:]}%
‘Of what ask ye me? Why tempt ye me?\evb\evg


\bvg\bva\mssnote{\Regius~1v/26, \RegiusProse~4v/27, \\
\Trajectinus~TODO, \Wormianus~TODO, \\
\Upsaliensis~TODO}%
\Ballnote{This st.~is quoted in \Gylfaginning~15:14, introduced by the words \emph{svá segir í Vǫlu-spá:} ‘so it says in the Spae of the Wallow’.  For the immediately preceding prose description of Mimer’s Well (\Gylfaginning~15:10--13) see note to st.~26 above.}%
\alst{A}lt vęit’k, \alst{Ó}ðinn, \hld\ hvar \alst{au}ga \edtext{falt}{\Afootnote{add., but marked with erasure dots \emph{þitt} \Regius}} &
\edtrans{ï inum \alst{m}ę́ra}{in the famed}{\Afootnote{so \Regius\Wormianus \emph{;} \emph{í þeim enom meira} ‘in the greater’ \Trajectinus \emph{;} \emph{i þeim enum mę́ra} ‘in the famed’ \Upsaliensis \emph{;} \emph{úr þeim enum mę́ra} ‘out of the famed’ \RegiusProse}} \hld\ \alst{M}ímis brunni; &
drekkr \alst{m}jǫð \alst{M}ímir \hld\ \alst{m}orgin hvęrjan &
af \edtrans{\alst{v}ęði}{pledge}{\Afootnote{\emph{†veiði†} \RegiusProse}} \alst{V}al-fǫðrs.‘ \hld\ \alst{V}ituð ér ęnn eða hvat?\eva

\bvb I know it all, Weden: where thine eye thou hidst \\
in the famed \inx[P]{Mimer’s Well}. \\
Mimer drinks mead every morning \\
from Walfather’s pledge.’ \emph{Know ye yet, or what?}\evb\evg


\bvg\bva\mssnote{\Regius~1v/29}%
Valði hęnni \alst{H}ęr-fǫðr \hld\ \alst{h}ringa ok męn, &
\edtrans{fe\emph{kk} \alst{sp}jǫll \alst{sp}ak⸗lig}{got foreseen tidings}{\Afootnote{emend.; \emph{‘fe spioll spaclig’} \Regius}\Bfootnote{The reading of \Regius\ may be interpreted either as (1) \emph{fé-spjǫll spak⸗lig} ‘foresighted wealth-tidings’ or (2) \emph{fé, spjǫll spak⸗lig} ‘wealth, foresighted tidings’.  Both are metrically deficient; (1) since a second element in a cpd.~like \emph{fé-spjǫll} cannot carry alliteration, and (2) since it has three strongly stressed nominals.  In both cases the first nominal \emph{fé} would be metrically expected to carry the alliteration.  The word \emph{fé} ‘wealth, cattle’ also makes little sense in context, since there is no transaction; Weden would be giving her expensive golden jewellery, of which he himself has no lack, in exchange for more gold.
The emendation replaces \emph{fé} with the verb \emph{fekk} ‘got, received’.  Verbs carry less stress than nominals, and the line thus becomes metrically equivalent to 28/3b \emph{drekkr mjǫð Mímir}.  Semantically we get a parallel to st.~1, where the wallow likewise says that she will relate \emph{spjǫll} ‘tidings, sayings’ (cf.~English \emph{gospel}, lit. ‘good news’ which originally translates the Greek \textgreek{εὐαγγέλιον}).
For further discussion on the metrics see \textcite[51--53]{Haukur2020}, \textcite[16]{Males2023Textual}.}} \hld\ ok \edtrans{\alst{sp}á-ganda}{spae-gands}{\Bfootnote{Spirits (familiars, gands; cf.~st.~21/2 above) sent out in order to gather hidden wisdom.}}; &
sá \alst{v}ítt ok of \alst{v}ítt \hld\ umb \alst{v}er-ǫld hvęrja.\eva

\bvb Host-father \name{= Weden} chose for her rings and necklaces; \\
he got foresighted tidings and \inx[C]{spae}-\inx[C]{gands}--- \\
she saw far and yet further through every world.\evb\evg


\bvg\bva\mssnote{\Regius~1v/30}%
Sá hǫ̇n \alst{v}al-kyrjur \hld\ \alst{v}ítt of komnar, &
\alst{g}ǫrvar at ríða \hld\ til \edtrans{\alst{g}oð-þjóðar}{folk of men}{\Bfootnote{The sense is uncertain, for in \Regius\ the difficult cluster \emph{tþ} in \emph{got-þjóð} (attested with \emph{-t-} in \HervararSaga\ st.~94/8b, \Skp\ 8) has been changed to \emph{ðþ}.  This means that the word \emph{goð-þjóð} may contain either (1) \emph{goti} ‘Got, man’ (thus ‘folk of the Gots’ or ‘folk of men’) or (2) \emph{goð} (‘folk of the Gods’).
Of these two alternatives, the element \emph{goti} in the sense ‘man’ is preferable, since the walkirries are presumably thought to be gathered in Walhall, ready to fly out into the world of men in order to attend battles and claim the best warriors for Weden (cf.~\Hakonarmal\ 1).  The present reading thus sees the cpd.~as semantically identical to \emph{ver-þjóð} ‘folk of men’, \textlink{Lokasenna}[24]/3.  The Walkirries do also have a particular association with the Gots, who were remembered as having fought the greatest battles of the Migration Period (cf.~note to \textlink{Volundarkvida}[1]/1).  For the formally identical \emph{Gót-þjóð} ‘folk, land of the Gots’ see \textlink{Helreid}[8]/1 n.}}: &
\edtext{\alst{Sk}uld hélt \alst{sk}ildi, \hld\ ęn \alst{Sk}ǫgul ǫnnur, &
\alst{G}unnr, Hildr, \alst{G}ǫndul \hld\ ok \alst{G}ęir-skǫgul; &
\alst{n}ú eru talðar \hld\ \edtrans{\alst{N}ǫnnur Hęrjans}{Nans of Harn \name{= Weden}}{\Bfootnote{\emph{Nanna} ‘\inx[P]{Nan}’ (the name itself is a nursery word) was the wife of \inx[P]{Balder}, but the word is here certainly being used to refer generically to ‘maidens, women’.  Cf.~Þul \emph{Ásynja} (\Skp\ 3), where the Walkirries are kenned \emph{Óðins męyjar} ‘Weden’s maidens’.}}, &
\alst{g}ǫrvar at ríða \hld\ \alst{g}rund, val-kyrjur.}{\lemma{Skuld \dots\ val-kyrjur. ‘Shild \dots\ earth.’}\Bfootnote{Judging especially by the phrase \emph{nú eru talðar} ‘now are tallied’, these four lines seem to be a later insert from a \inx[C]{thule} counting the Walkirries.}}\eva

\bvb She saw \inx[P]{Walkirries} come from afar, \\
ready to ride to the folk of men. \\
Shild held a shield and Shagle another, \\
Guth, Hild, Gandle and Goreshagle--- \\
now are tallied the Nans of Harn \name{= Weden}, \\
Walkirries ready to ride o’er the earth.\evb\evg


\bvg\bva\mssnote{\Regius~2r/2}%
\Ballnote{Told allusively in sts.~31--33 is the myth about Balder’s death.  The outline of this myth is as follows: Balder, the son of Weden and Frie, was killed by a mistletoe shot by his blind half-brother Hath, whose hand was guided by Lock.  As Balder’s father, it was Weden’s duty to avenge him, yet he could not kill Hath who was also his son.
Weden thus devised a plan and through sorcery seduced the woman Wrind (so KormǪ \emph{Sigdr} 3, \Skp\ 3: \emph{sęið Yggr til Rindar} ‘Ug won Wrind through sorcery’).
Wrind gave birth to Wonnel, who grew up fast, and after just one day was big enough to kill Hath and avenge Balder.

Other sources for Balder’s death include the poems \textlink{Baldrsdraumar}[8]--11, \textlink{Lokasenna}[28], and \textlink{Hyndluljod}[28] and the prose texts \Gylfaginning~49 and \textcite{Saxo} 3.4.1--8.
For the punishment of Lock for his role in the killing see sts.~34 and H31 below along with \textlink{Lokasenna}[P8].

The most similar text to \Voluspa\ 31--33 is \textlink{Baldrsdraumar}[8]--11, with which it is clearly textually related (cf.~st.~32/4 n.).  The chief narrative difference between the two passages is that \textlink{Baldrsdraumar}[11]/1 mentions Wrind, who is not found in \Voluspa.

Balder’s death is treated at greater length in \Gylfaginning~49:3--25, which may be briefly summarised as follows: Balder has terrible nightmares about his own death, and so his mother Frie makes all sorts of things (fire, water, venom, metals, stones, trees, diseases, animals) swear oaths not to harm him.  After this, the Eese make sport of shooting and cutting him since he cannot be harmed.  Lock is annoyed by the spectacle and tricks Frie into telling him that there was one thing that did not swear the oath---the mistletoe, which was thought too young.
Lock makes a mistletoe-arrow and gives it to the blind god Hath, guiding his hand towards Balder.  Hath shoots and Balder dies.
After this, \Gylfaginning~49 describes Balder’s funeral, treated poetically in \Husdrapa, at some length, alongside how the Eese failed to “weep Balder out of hell” (see \textlink{EddicFragments}[F5][V]).  The binding of Lock is described in \Gylfaginning~50 (cf.~st.~34 below).
Most notably, \Gylfaginning~49 neglects to mention both Wonnel and Wrind.  That part of the myth may have been left out of the Prose Edda for moral reasons due to its motifs of rape and kinslaying, but that it was known to the author is proven by two passages: \Gylfaginning~30:1--2: \emph{Áli eða Váli heitir einn, sonr Óðins ok Rindar.  Hann er djarfr í orrostum ok mjǫk happ-skęytr.} ‘Onnel or Wonnel is the name of one, the son of Weden and Wrind.  He is brave in battles and a very lucky shot.’  \Skaldskaparmal\ 19: \emph{Hvernig skal kenna Vála? Svá, at kalla hann son Óðins ok Rindar, [\dots] hefni-ǫ́s Baldrs, dólg Haðar ok bana hans [\dots]} ‘How shall one ken Wonnel?  Namely, by calling him the son of Weden and Wrind, [\dots] avenging \inx[C]{os} of Balder, the foe of Hath and his bane [\dots]’

The other lengthy prose retelling of the myth is \textcite{Saxo} 3.4, which relates the revenge narrative in typical euhemerized form, turning the gods Hath and Balder into flawed human rulers.  The outline of that version is as follows: After Balder (Latin \emph{Balderus}) has succumbed to a mortal wound dealt to him by Hath (\emph{Høtherus}) with a treacherous sword, his father Weden (\emph{Othinus}) takes counsel from a group of wizards.  One of them, Horsethief the Finn (\emph{Rostiophus Phinnicus}), foresees that Wrind (\emph{Rinda}), daughter of the Russian king, will bear him a son to avenge Balder.

Weden soon enlists in the Russian king’s army and leads it to great victories, but he is spurned by his daughter.  He tries various disguises but is continually refused, until at last he disguises himself as an old witch and volunteers himself as her private physician.  When Wrind becomes ill, he binds her, ostensibly in order to administer a certain foul potion, but instead rapes her, implicitly with her father’s consent.  Their son Bo grows up to become a fierce raider.  One day Weden summons him and reminds him of his duty to avenge his half-brother Balder; Bo slays Hath in a fierce duel but soon himself perishes from his wounds.
It is worth noting that the seduction of Wrind as told by \citeauthor{Saxo} is unlikely to be a fantasy of that author, but on the contrary seems to derive from a now-lost West Norse source.  The narrative bears a particular structural resemblance to Shirner’s coercion of Gird in \textlink{Skirnismal}[19]--36; see further \textlink{Skirnismal}[19] n., \textcite{Brate1913}.}%
Ek sá \alst{B}aldri, \hld\ \alst{b}lóðgum \edtrans{tífur}{victim’s}{\Bfootnote{This word is rather difficult and possibly corrupt.  (1) It has been connected with \emph{týr} ‘tew, god’, but the dat. sg. of \emph{týr} is \emph{tívi} and the intrusive \emph{r} is unexplained.
(2) A better explanation is given by \CV, which connects it with OE \emph{tiber, tifer} ‘victim, hostage’, but this also has problems.  \emph{blóðgum} ‘bloody’ is m.~dat.~sg., but OE \emph{tiber} is neuter.  If we are dealing with a m.~noun \emph{*tífurr} with the same declension as \emph{jǫfurr}, we would expect dat.~sg.~\emph{*tífri}--- not \emph{tífur} (which would, however, be the expected acc.~sg.).}}, &
\alst{Ó}ðins barni, \hld\ \alst{ø}r·lǫg \edtrans{folgin}{hid}{\Bfootnote{The sense is either that it is Balder’s orlay to be hid, that is, in the Underworld, whither he goes after his death (see \textlink{EddicFragments}[F5][V]), or that his orlay is metaphorically sealed and unalterable.
The verb \emph{fela} ‘hide, conceal’, of which \emph{folgin} is the n.~acc.~pl. past participle, is used in poetry to describe burial in mounds, as in \Ynglingatal\ 24 (“[...] And afterwards the victory-raisers hid (\emph{fǫ́lu}) the ruler on Borrey”) and the 10th c.~Karlevi runestone (Öl 1, ed.~below under Runic Poetry from Sweden and Gotland).}}; &
stóð of \alst{v}axinn \hld\ \edtext{\alst{v}ǫllu\emph{m}}{\Afootnote{\emph{vǫllu} (with omitted abbreviation mark for \emph{-m}) \Regius}} hę́ri &
\alst{m}jór ok \alst{m}jǫk fagr \hld\ \alst{m}istil-tęinn.\eva

\bvb I saw Balder’s, the bloody victim’s, \\
Weden’s child’s \inx[C]{orlay} hid: \\
there stood grown higher than the plains \\
a slender and most fair mistletoe.\evb\evg


\bvg\bva\mssnote{\Regius~2r/4}%
Varð af \alst{m}ęiði, \hld\ þęim’s \alst{m}jór sẏndi⸗sk, &
\alst{h}arm-flaug \alst{h}ę̇ttlig, \hld\ \alst{H}ǫðr nam skjóta. &
\alst{B}aldrs \alst{b}róðir vas \hld\ of \alst{b}orinn snimma, &
\edtext{sá nam, \alst{Ó}ðins sonr, \hld\ \alst{ęi}n-nę́ttr vega}{\Bfootnote{The three lines making up \Voluspa\ 32/4--33/2 are found near-identically (but with different tense) as \textlink{Baldrsdraumar}[11]/2--4.  The direction of influence is unclear, since both poems appear fairly old.  It is also plausible that both are borrowing from a now-lost third source dealing with Balder’s death; that such a text must have existed is proven by the stanza in \textlink{EddicFragments}[F5][V].}}.\eva

\bvb Of that tree which slender seemed \\
came a baneful harm-flier---Hath took to shoot. \\
Balder’s brother \ken*{= Wonnel} was born early; \\
he took, Weden’s son, one night old, to fight.\evb\evg


\bvg\bva\mssnote{\Regius~2r/6}%
\edtext{\edtext{Þó}{\lemma{Þó \dots\ kęmbði ‘washed \dots\ combed’}\Bfootnote{A collocation; see note to \textlink{Havamal}[61] for discussion and other examples.  Wonnel, being bound by the sworn mission to avenge his brother, could not engage in such acts of personal vanity.  For the motif of a warrior bound not to engage in acts of personal hygiene until having fulfilled a vow we may compare the famous vow of Harold Fairhair, \emph{at aldri [...] skera hár mitt né kemba, fyrr en ek hefi eignat⸗sk allan Nóreg með skǫttum ok skyldum ok for·ráði, en deyja at ǫðrum kosti.} ‘never to cut my hair nor comb it before I have come to own all of Norway with taxation and duties and stewardship, but to die otherwise.’ (\emph{Haralds saga ins hár-fagra}, 4).}} ę́va \alst{h}ęndr \hld\ né \alst{h}ǫfuð kęmbði, &
áðr ȧ \alst{b}ál of \alst{b}ar \hld\ \alst{B}aldrs and·skota}{\lemma{Þó \dots\ and·skota ‘He washed \dots\ opponent’}\Bfootnote{Found near-identically in \textlink{Baldrsdraumar}[11]/3--4; see note to st.~32/4 above.}}; &
ęn \alst{F}rigg of grét \hld\ ï \edtrans{\alst{F}ęn-sǫlum}{the Fenhalls}{\Bfootnote{Never mentioned elsewhere in poetry, but \Gylfaginning~35:1--2 explains: \emph{Þá mę́lti Gangleri: „Hverjar eru ǫ́synjurnar?“  Hárr segir: „Frigg er ǿðst.  Hon á þann bǿ, er Fen-salir heita, ok er hann all-veg⸗ligr.} ‘Then spoke Gangler: “Which are the Ossens?”  High says: “Frie is the noblest.  She owns the dwelling which is called the Fenhalls, and it is all glorious. [...]”’}} &
\edtrans{\alst{v}ǫ́ \alst{V}al-hallar}{the woe of Walhall}{\Bfootnote{The deaths of two sons: Balder and Hath.}}. \hld\ \alst{V}ituð ér ęnn eða hvat?\eva

\bvb He ne’er washed his hands nor combed his head \\
before onto the pyre he bore Balder’s opponent \ken*{= Hath}, \\
but Frie lamented in the Fenhalls \\
the woe of \inx[P]{Walhall}. \emph{Know ye yet, or what?}\evb\evg


\bvg\bva\mssnote{\Regius~2r/8, \Hauksbok~20v/13}%
\Ballnote{After Balder was avenged, the Eese went to catch his true killer, Lock, as told in \textlink{Lokasenna}[P8] and \Gylfaginning~50:3--25.  They bound him with his own son’s intestines and a snake was placed above his face to drool venom onto it.  His wife, Syein, sat over him and collected the venom in a small basin, but when she had to empty it, he shook so furiously from the pain of the acidic venom that the earth shook.}%
\edtext{\alst{H}apt sá hǫ̇n liggja \hld\ und \alst{H}vera-lundi &
\edtrans{\edtext{\alst{l}ę́-gjarn\emph{s}}{\Afootnote{gramm.~emend.; \emph{lę́-gjarn} (acc.~sg.) \Regius}}}{guile-eager}{\Bfootnote{An epithet of Lock, who is strongly associated with \emph{lę́} ‘guile, deceit’.
His name more often occurs paired with the adjective \emph{lę́-víss} ‘guile-wise’, so \textlink{Hymiskvida}[37]/4, \textlink{Lokasenna}[54]/5, and \Sorlathattr\ 2.}} \alst{l}íki \hld\ \alst{L}oka ȧ·þękkjan;}{\lemma{Hapt \dots\ ȧ·þękkjan ‘A captive \dots\ Lock’}\Afootnote{om.~and replaced by H31 \Hauksbok.}} &
\alst{þ}ar sitr Sigyn \hld\ \alst{þ}ęy⸗gi of sïnum &
\alst{v}eri \alst{v}ęl-glýjuð. \hld\ \alst{V}ituð ér ęnn eða hvat?\eva

\bvb A captive she saw lying beneath Wharlund: \\
a guile-eager man’s form, recognised as Lock. \\
There sits Syein not at all cheerful, \\
o’er her husband. \emph{Know ye yet, or what?}\evb\evg


\bvg\bva\mssnote{\Regius~2r/10}%
\Ballnote{Having recounted the death of Balder and the binding of Lock, the wallow moves on to describing the places of evil in the world (sts.~35--41).  The world has lost its innocence and the forces of chaos and disorder in the east and north plot its destruction.  The present stanza is somewhat obscure, likely because half of it is missing.  It is not in any way alluded to in \Gylfaginning, nor is it found in \Hauksbok.  The river Slide is also mentioned in the long list of rivers in \textlink{Grimnismal}[28]/4.}%
\alst{Ǫ́} fęllr \alst{au}stan \hld\ of \alst{Ęi}tr-dala &
\alst{s}ǫxum ok \alst{s}verðum, \hld\ \edtrans{\alst{S}líðr}{Slide}{\Bfootnote{An adjective describing a blade so sharp that the finger is cut when slid across it.  Cf.~\textlink{Atlakvida}[23]: \emph{saxi slíðr-bęitu} (dat. sg.) ‘slide-biting sax’.}} hęitir sú.\eva

\bvb A river falls from the east through the Venom-dales, \\
{[a river]} of saxes and swords; Slide it is called.\evb\evg


\bvg\bva\mssnote{\Regius~2r/11}%
\Ballnote{Sts. 36--38 are paraphrased in \Gylfaginning~52:1--10, describing the world after its destruction and rebirth:

\emph{Þá mę́lti Gangleri: „Hvat verðr þá eptir, er brenndr er himinn ok jǫrð ok heimr allr, ok dauð goð’in ǫll ok allir Ein-herjar ok alt mann-folk, ok hafið ér áðr sagt, at hverr maðr skal lifa í nǫkkvǫrum heimi um allar aldir?“  Þá svarar Þriði: „Margar eru þá vistir góðar ok margar illar; batst er þá at vera á Gim·léi á himni, ok all-gótt er til góðs drykkjar þeim, er þat þykkir gaman, í þeim sal, er Brimir heitir; hann stendr ok á himni. Sá er ok góðr salr, er stendr á Niða-fjǫllum, gǫrr af rauðu gulli; sá heitir Sindri. Í þessum sǫlum skulu byggja góðir menn ok sið-látir.  Á Ná-strǫndum er mikill salr ok illr ok horfa norðr dyrr; hann er ok ofinn allr orma-hryggjum sem vanda-hús, en orma hǫfuð ǫll vitu inn í hús’it ok blása eitri, svá at eptir sal’num renna eitr-ár, ok vaða þę́r ár eið-rofar ok morð-vargar, svá sem hér segir:“}

‘Then spoke Gangler: “What will then happen thereafter when the heaven and the earth and the whole world are burned, and the Gods all are dead and all the Oneharriers and all man-kind---and [yet] ye have said earlier that every man will live in a certain home for all ages?”  Then answers Third: “Many good dwellings are there then, and many evil ones; then it is best to be in Gemlee in the heaven.  And in the hall which is called Brimmer it is fully furnished with good drink for those who find it pleasurable; that one also stands on the heaven.  Another good hall is the one which stands on the Nithfells, made out of red gold; it is called Sinder.  In these halls shall dwell good and virtuous men.  On Neestrand there is a great hall and evil, and its doors face north.  It is all interlaced with the spines of serpents like a wicker-house, but the heads of the serpents all face into the house and blow venom so that there run through the hall rivers of venom, and in those rivers wade oath-breakers and murder-wargs, just as it says here:”’

After this passage \Gylfaginning\ cites sts.~37 and 38/1--2, followed by the prose: \emph{En í Hver-gelmi er verst} ‘But in Wharyelmer it is the worst’, and a citation of st.~38/4.

It ought to be noted that the placement and contents of sts.~36--38 most likely exclude them from being descriptions of the cosmology of the world to come.  That interpretation instead seems to be one of Snorre’s own inventions, as is strongly supported by the grammatical form \emph{męin-svara} in st.~38/2.}%
Stóð fyr \alst{n}orðan \hld\ ȧ \edtrans{\alst{N}iða-vǫllum}{Nithwolds}{\Bfootnote{‘The dark plains’.  The paraphrase of this st.~in \Gylfaginning~52:6 suggests that the version of \Voluspa\ used by Snorre had the reading \emph{*Niða-fjǫllum} ‘Nithfells’ (cf.~st.~62).  The first element is, in any case, \emph{nið} ‘the waning of the moon’ and refers to great darkness; cf.~the adjective \emph{niða-myrkr} ‘pitch darkness’.}} &
\alst{s}alr ór golli \hld\ \alst{S}indra ę́ttar; &
ęn \alst{a}nnarr stóð \hld\ ȧ \alst{Ȯ}·kólni, &
\alst{b}jór-salr jǫtuns, \hld\ \edtrans{ęn sá \alst{B}rimir hęitir}{and he/it is called Brimmer}{\Bfootnote{It is syntactically ambiguous whether Brimmer is the name of the ettin or the hall itself.  \Gylfaginning~52:5 clearly considers it the name of the hall.}}.\eva

\bvb To the north on the Nithwolds stood \\
a golden hall, of Sinder’s lineage \ken{dwarfs}, \\
but another one stood on Uncolner, \\
an ettin’s beer-hall, and he/it is called Brimmer.\evb\evg


\bvg\bva\mssnote{\Regius~2r/13, \Hauksbok~20v/19, \\
\RegiusProse~17v/1, \Trajectinus~TODO, \\
\Wormianus~TODO, \Upsaliensis~18v/19}%
\Ballnote{Cited in \Gylfaginning~52; see st.~36 above.}%
\alst{S}al \edtrans{sá hǫ̇n}{she saw}{\Afootnote{\emph{vęit’k} ‘I know’ \GylfMS.}} standa \hld\ \alst{s}ólu fjarri &
\alst{N}á-strǫndu ȧ, \hld\ \alst{n}orðr horfa dyrr; &
falla \alst{ęi}tr-dropar \hld\ \alst{i}nn umb ljóra, &
sá ’s \alst{u}ndinn salr \hld\ \alst{o}rma hryggjum.\eva

\bvb A hall she saw standing far from the sun \\
on Neestrand; north face its doors. \\
Venom-drops fall in through the smoke-vent; \\
that hall is wound with the spines of snakes.\evb\evg


\bvg\bva\mssnote{\Regius~2r/15, \Hauksbok~20v/21, \\
\RegiusProse~17v/2, \Trajectinus~TODO, \\
\Wormianus~TODO, \Upsaliensis~18v/21}%
\Ballnote{Ll.~1--2, 4 are cited in \Gylfaginning~52; see st.~36 above. --- This st.~takes a distinctly moralising tone otherwise rare in pre-Christian poetry; cf.~st.~44 below.  The crimes are what one might expect in a source from Germanic antiquity: perjury, murder and adultery.
In Scandinavian and Anglo-Saxon laws the perpetrator of such crimes held the dishonourable title of \inx[C]{nithing}, that is, one afflicted with \inx[C]{nithe} (‘evil, shame’), and was thereby stripped of his personal rights.  It is fitting, then, that such nithings should be tortured by a creature named Nithehewer (‘Nithe-striker’).

For the crimes we may cf.~the list of \emph{ú·dáða-męnn} (‘criminals’, lit.~‘men of misdeeds’) not eligible for burial in church grounds according to \Gulatingslog\ 23: \emph{dróttins-svikar} ‘traitors to their lord’, \emph{morð-vargar} ‘murder-wargs’, \emph{trygg-rofar} ‘truce-breakers’, \emph{þjófar} ‘thieves’, and those men who \emph{sjalfir spilla ǫnd sinni} ‘themselves destroy their spirit’, i.e., commit suicide.  According to that law these men were instead to be buried in a \emph{flǿðar-mál} ‘flood-mark’, where \emph{sę́r mǿti⸗sk ok grǿn torfa} ‘see and green turf meet’ \parenciteshorthand[13]{NGL1}, something which also has clear parallels in the \emph{Lex Frisionum} (TODO).

The practice of burying criminals deemed particularly shameful not just in unhallowed ground but specifically in bogs and flood-marks is well attested both in archeology and in literary sources relating to the Germanic peoples from Tacitus (\emph{Germania}, ch. 12) onwards.  Upon examining these parallels together, it is hard not to see the \emph{þungir straumar} ‘heavy streams’ of the present stanza as a reflex of such customs, so that the whole st.~in fact appears to describe the Heathen religious belief about bog burials, namely that those buried in them were punished by a watery dragon in the afterlife.

Watery punishment in the afterlife is also described in \textlink{Reginsmal}[3]--4 and possibly in \textlink{Grimnismal}[21], and a perjurious woman is drowned in a bog in \textlink{GudrunThree}[10].  The argument is pursued at at greater length in \textcite{RosenbergForthcoming}.}%
\edtrans{Sá hǫ̇n}{she saw}{\Afootnote{so \Regius \emph{;} \emph{sér hǫ̇n} ‘she sees’ \Hauksbok \emph{;} \emph{skulu} ‘shall [wade]’ \GylfMS}} \alst{þ}ar vaða \hld\ \alst{þ}unga strauma &
\alst{m}ęnn \edtext{\edtrans{\edtext{\alst{m}ęin-svara}{\Afootnote{\emph{męins-vara} \Regius \emph{;} \emph{męin-svarar} \Upsaliensis}}}{false-swearing}{\Bfootnote{Acc.~pl.~of \emph{*męin-svari} or \emph{*męin-svarr}, an otherwise unattested adj.  Related words are \emph{męin-sǿri} ‘false oath, perjury’ and \emph{męin-svari} m. ‘perjurer’.
This adjective is in the acc.~pl.~in 3/4 mss.~of \Gylfaginning, but this is grammatically irregular since the \GylfMS\ version of line 1 has the verb \emph{skulu} ‘shall [wade]’ which takes a nominative subject.
The most likely explanation for this is that Snorre, when quoting the present stanza, deliberately changed the first words of l.~1 from an archetypal \emph{*sér/sá hǫ̇n} ‘she sees/saw’ (shared with \Regius\Hauksbok) to \emph{skulu} in order to make the poetry fit better with his eschatological interpretation (cf.~st.~36 n.).  In doing so he overlooked the grammatical form of the object phrase \emph{męnn męin-svara}, probably due to the rarity of the adjective.
This has important consequences for our evaluation of Snorre’s reliability when citing Eddic verse since it shows an instance of his deliberately changing the wording of a stanza to conform to his prose.  Even so, the present instance should probably be considered exceptional, since most other textual variants in Eddic sts.~cited in \Gylfaginning\ are better explained as resulting from scribal or oral transmission.}} \hld\ ok \edtrans{\alst{m}orð-varga}{murder-wargs}{\Afootnote{\emph{morð-vargar} (nom.~pl.) \Regius\RegiusProse\Trajectinus\Wormianus \emph{;} \emph{morðingjar} ‘murderers’ \Upsaliensis.}\Bfootnote{Men outlawed for murder, which is not all instances of the taking of human life,
but specifically when done clandestinely or in ways deemed particularly reproachable.}}}{\lemma{męin-svara ok morð-varga}\Afootnote{corr.~from \emph{morð-varga męin-svara ok} by an inversion mark in same hand \Regius.}} &
\edtext{ok þann’s \alst{a}nnars glępr \hld\ \edtrans{\edtext{\alst{ęy}ra-ru̇nu}{\Afootnote{\emph{†eyrna runa†} \Hauksbok}}}{ear-whisperess}{\Bfootnote{Acc.~sg.~of \emph{ęyra-ru̇na}, with the second element derived from \emph{ru̇n} ‘whisper, secret’; cf.~the masculine \emph{(of) ru̇ni} ‘friend, confidant’, OE \emph{ge·rúna} ‘counsellor’.  This word is also used in \textlink{Havamal}[115] which likewise contains a warning not to seduce another man’s wife.}}.}{\Afootnote{om.~\GylfMS.}} &
Þar \edtrans{saug}{sucked on}{\Afootnote{so \Hauksbok \emph{;} \emph{†sv́g†} \Regius \emph{;} \emph{kvęlr} ‘torments’ \GylfMS}} \edtrans{\alst{N}íð·hǫggr}{Nithehewer}{\Bfootnote{Also appearing in st.~62 below and further in \textlink{Grimnismal}[36]/4.}} \hld\ \alst{n}ái fram-gingna; &
slęit \alst{v}argr \alst{v}era. \hld\ \alst{V}ituð ér ęnn eða hvat?\eva

\bvb There she saw wading through heavy streams \\
false-swearing men and murder-wargs, \\
and him who seduces another man’s ear-whisperess \ken{wife}. \\
There \inx[P]{Nithehewer} sucked on passed-on corpses; \\
the warg tore at men. \emph{Know ye yet, or what?}\evb\evg


\bvg\bva\mssnote{\Regius~2r/17, \Hauksbok~20v/2, \\
\RegiusProse~4r/10, \Trajectinus~TODO, \\
\Upsaliensis~5r/27, \Wormianus~TODO}%
\Ballnote{The old hag, who is not identified in any other source, raises the cubs of the wolf Fenrer.  One of these wolves, identified by Snorre as Moongarm (\emph{Mȧna-garmr}), will eventually swallow the moon.  For the myth of the wolves that chase the sun and moon cf.~\textlink{Grimnismal}[40], according to which the sun is chased by the wolf Sholl, while another wolf, Hater, runs in front of it, and \textlink{Vafthrudnismal}[46]--47, where the sun is said to be swallowed by Fenrer.  It is likely that even the Heathen poets themselves were not entirely in agreement about the exact details of these future events; for further “wolfish” confusion, cf.~st.~43/1--2 below.

These wolves are dealt with in \Gylfaginning~12, which synthesizes from \Voluspa\ 39--40 along with the two aforementioned sources.  Relevant for the present sts.~is \Gylfaginning~12:7, 9--12:

\emph{Gýgr ein býr fyrir austan Mið·garð í þeim skógi, er Járn-viðr heitir. [...] In gamla gýgr fǿðir at sonum marga jǫtna ok alla í vargs líkjum, ok þaðan af eru komnir þessir ulfar. Ok svá er sagt, at af ę́tt’inni verðr sá einna máttkastr, er kallaðr er Mána-garmr. Hann fylli⸗sk með fjǫrvi allra þeira manna, er deyja, ok hann gleypir tungl, en støkkvir blóði himin ok loft ǫll.  Þaðan týnir sól skini sínu, ok vindar eru þá ó·kyrrir ok gnýja heðan ok handan.}

‘A lonely ogress lives to the east of Middenyard in the forest called Ironwood.  The old ogress raises many ettins as her sons, all in the likenesses of wolves, and thereof these wolves [= Sholl and Hater] come.  And so it is said that from that lineage a single one becomes the mightiest, and he is called \inx[P]{Moongarm}.  He fills himself with the life of all those men who die and he swallows the moon and stains the heaven and all the air with blood.  Thereof the sun loses her shine and then the winds are violent and howl hither and thither.’

Immediately after this passage \Gylfaginning~12:13 quotes \Voluspa\ 39--40, introduced by the words \emph{svá segir í Vǫlu-spá:} ‘so it says in the Spae of the Wallow:’.}%
\edtrans{\alst{Au}str}{In the east}{\Bfootnote{The cardinal direction associated with ettins and other fiendish creatures; cf.~\textlink{Hymiskvida}[5]/1 n.}} \edtrans{býr}{dwells}{\Afootnote{so \Hauksbok\GylfMS \emph{;} \emph{sat} ’sat/stayed’ \Regius}} \edtrans{hin \edtrans{\alst{a}ldna}{old}{\Afootnote{\emph{arma} ‘wretched’ \Upsaliensis}}%
}{the old hag}{\Bfootnote{Comparatively, the ‘hag’ (an insertion in the translation to clarify the gender of ON \emph{aldna}) appears to correspond to the Finnish \emph{Louhi}, the toothless mistress of \emph{Pohjola} (the Northland).
In \emph{Kalevala} 47--49 she steals away the sun, moon, and fire (see st.~24 n.~above).
Both women are part of the evil forces that constantly struggle against the Gods by disrupting the natural cosmic order and its cycles.}} \hld\ ï \edtrans{\alst{Éa}rn-viði}{Ironwood}{\Afootnote{metr. emend.; \emph{Járn-viði} \Regius\Hauksbok\RegiusProse\Wormianus\Upsaliensis \emph{;} \emph{Járn-viðjum} ‘Ironwoods’ \Trajectinus}} &
ok \edtrans{\alst{f}ǿðir}{rears}{\Afootnote{so \Hauksbok\GylfMS \emph{;} \emph{fǿddi} ‘reared’ \Regius}} þar \hld\ \alst{F}ęnris kindir; &
verðr \edtext{af}{\Afootnote{\emph{ór} \RegiusProse\Trajectinus}} þęim \alst{ǫ}llum \hld\ \alst{ęi}nna nøkkurr &
\alst{t}ungls \edtrans{\alst{t}júgari}{snatcher}{\Afootnote{\emph{†tuigan†} \Trajectinus \emph{;} \emph{tregari} ‘troubler’ \Upsaliensis.}\Bfootnote{Hapax, generally understood as either \emph{*tjúga} ‘drag, pull’ or \emph{tjúga} f.~‘pitchfork’ + \emph{-ari} ‘er’.
The agentive suffix \emph{-ari} is a young feature found nowhere else in the poem, and it is therefore possible that this word is corrupt.  Note that \Upsaliensis’s \emph{tregari} ‘troubler’ is metrically impossible and prob.~an instance of trivialisation from \emph{*tjúgari}, as found in the other mss.~of \Gylfaginning\ \parencite[181, 194]{Males2023Textual}.
If the word is corrupt, it must have occurred early in transmission, since \emph{tjúgari} is found in all three recensions (\Regius\Hauksbok\GylfMS).}} \hld\ ï \alst{t}rolls hami.\eva

\bvb In the east dwells the old hag in \inx[P]{Ironwood} \\
and there rears the brood of \inx[P]{Fenrer} \ken{wolves}. \\
Out of them all arises one, I know not which: \\
a snatcher of the moon in a troll’s \inx[C]{hame}.\evb\evg


\bvg\bva\mssnote{\Regius~2r/19, \Hauksbok~20v/4, \\
\RegiusProse~4r/11, \Trajectinus~TODO, \\
\Upsaliensis~5r/28, \Wormianus~TODO}%
\alst{F}ylli⸗sk \alst{f}jǫrvi \hld\ \alst{f}ęigra manna, &
\alst{r}ýðr \alst{r}agna sjǫt \hld\ \alst{r}auðum dręyra, &
\edtrans{\alst{s}vǫrt verða \alst{s}ól-skin}{Black turn the sun’s rays}{\Afootnote{so \Hauksbok\GylfMS \emph{;} \emph{svart var þá sól-skin} ‘Black was then the sun’s shine’ \Regius}\Bfootnote{As part of the Fimblewinter; see \textlink{Vafthrudnismal}[44]/3--4 n. ---
The reading of \Regius, although semantically meaningful, is outvoted by the other mss.  It is also unlikely on its own merit, since the past tense of \emph{var} ‘was’ disagrees with the present tense of the other verbs in this st.  For this reason, \Regius’s \emph{‘vͬ þa’} (resolving as \emph{var þa} ‘was then’) is most likely a scribal error, where the archetypal \emph{*‘v͛þa’} (resolving as \emph{*verþa}, so \RegiusProse\Upsaliensis) was miscopied as \emph{*‘vͬþa’} (resolving as \emph{*varþa}), and then reinterpreted as \emph{var þá} by the addition of a space.}} \hld\ of \alst{s}umur ęptir, &
\alst{v}eðr ǫll \edtext{\alst{v}á-lynd}{\Afootnote{\emph{†va. ly.} \Upsaliensis}}. \hld\ \alst{V}ituð ér ęnn eða hvat?\eva

\bvb He fills himself with the lifeblood of \inx[C]{fey} men; \\
he reddens the abode of the \inx[P]{Reins} \ken{heaven} with red gore. \\
Black turn the sun’s rays in the summers thereafter, \\
the winds all woeful. \emph{Know ye yet, or what?}\evb\evg


\bvg\bva\mssnote{\Regius~2r/21, \Hauksbok~20v/16}%
\edtrans{\alst{S}at þar ȧ haugi}{There he sat on the mound}{\Bfootnote{The motif of ettins sitting on burial mounds is also found in \textlink{Thrymskvida}[6] and \textlink{Skirnismal}[11].
The precise significance of this is uncertain, but mounds are liminal spaces where the worlds of the dead and the living are particularly close, and they feature prominently in pre-Christian ritual (TODO: PCRN-HS reference).}} \hld\ ok \alst{s}ló hǫrpu &
\alst{g}ýgjar hirðir, \hld\ \alst{g}laðr Ęgg·þér; &
\edtrans{\alst{g}ól of hǫ̇num}{Over him crowed}{\Bfootnote{The crowing of three cocks---the first in Ettinham, the second in Walhall, the third in Hell---heralds the destruction of the world.}} \hld\ ï \edtrans{\alst{G}agl·viði}{Galewood}{\Bfootnote{An otherwise unknown location; the first element is \emph{gagl} ‘wild goose’.  Galewood is perhaps the same as Ironwood.}} &
\alst{f}agr-rauðr hani, \hld\ sá’s \alst{F}jalarr hęitir.\eva

\bvb There he sat on the mound and struck the harp, \\
the gow’s herdsman, glad \inx[P]{Edgethew}. \\
Over him crowed in \inx[P]{Galewood} \\
the fair-red cock who is called Feller.\evb\evg


\bvg\bva\mssnote{\Regius~2r/23, \Hauksbok~20v/18}%
\alst{G}ól of ǫ̇sum \hld\ \edtrans{\alst{G}ollin·kambi}{Goldencomb}{\Bfootnote{This cock may be depicted on two of the 11th or 12th c.~tapestries from Överhogdal, Härjedalen in Sweden.  The relevant tapestries show a great tree surrounded by a mass of animals, and on the top of the tree stands a bird.  The likely pre-Christian symbolism is confirmed by a depiction of an eight-legged horse, certainly to be identified with Weden’s horse Slapner.}}, &
sá vękr \alst{h}ǫlða \hld\ \edtrans{at \alst{H}ęrja-fǫðrs}{at the Father of Hosts’}{\Bfootnote{In Weden’s hall, Walhall.}}, &
ęn \alst{a}nnarr gęlr \hld\ fyr \alst{jǫ}rð neðan &
\alst{s}ót-rauðr hani \hld\ at \alst{s}ǫlum Hęljar.\eva

\bvb Over the Eese crowed Goldencomb; \\
he wakes the heroes at the Father of Hosts’, \\
but another one crows beneath the earth: \\
a soot-red cock in the halls of Hell.\evb\evg


\bvg\bva\mssnote{\Regius~2r/25}%
\Ballnote{The first occurrence of the third refrain (\emph{stęf}), repeated several times in the section of the poem dealing with the Rakes of the Reins (ON \emph{ragna rǫk}), culminating with st.~55.  The destruction of the world, involving the deaths of all major Gods save Balder, is treated at length in \Gylfaginning~51.  Furthermore, it is the subject of a wisdom contest in \textlink{Vafthrudnismal}[44]--53 and alluded to in \textlink{Lokasenna}[39], 41.}%
\edtext{\alst{G}ęyr \alst{G}armr mjǫk \hld\ fyr \alst{G}nipa-hęlli, &
\alst{f}ęstr mun slitna, \hld\ ęn \alst{f}reki rinna}{\lemma{Garm \dots\ rinna ‘Garm \dots\ run’}\Bfootnote{\Gylfaginning~51 distinguishes between the Fenrerswolf fought by Weden and the wolf Garm fought by Tew.  Yet this is probably the result of later systematizing.  The wolf which is said to come loose in the present stanza must surely be the Fenrerswolf, and since in \Gylfaginning~34:10--42 it is primarily Tew who deals with the Fenrerswolf, we would naturally expect him to fight it during the Rakes of the Reins as well, just like Thunder fights his rival, the Middenyardswyrm, and Lock fights Homedall (\Gylfaginning~51:48; see \textlink{Thrymskvida}[13]/3 n.).
\Voluspa\ 51, however, sees Weden fight the Wolf instead, and this requires a separate wolf for Tew to fight, hence Garm.  It is thus likely that these two wolves were originally not distinct creatures.  For the story of the Fenrerswolf see \textlink{Lokasenna}[38]/3--4; for further “wolfish” confusion, cf.~st.~39 above.}}; &
\alst{f}jǫlð vęit hǫ̇n \alst{f}rǿða, \hld\ \alst{f}ramm \edtrans{sé’k}{I see}{\Bfootnote{In \Regius\ the use of the present tense \emph{sé’k, sér} ‘I see, she sees’ marks a shift from the past tense (\emph{sá} ‘saw’) of the Balder-episode (31--34) and the description of Neestrand (37--38), a subtle but important change.  The use of the present tense continues for the rest of the poem.
Note that in \Hauksbok\ the verb is found in the present tense already in 37--38, and we lack the Balder-episode to compare with.}} lęngra &
of \alst{r}agna \alst{r}ǫk, \hld\ \alst{r}ǫmm sig-tíva.\eva

\bvb Garm barks loudly before the Gnip-caves; \\
the rope will snap and the Wolf will run. \\
She knows much lore; I see further ahead \\
about the mighty \inx[P]{Rakes of the Reins}, of the Victory-Tews \ken{gods}.\evb\evg


\bvg\bva\mssnote{\Regius~2r/28, \Hauksbok~20v/24, \\
\RegiusProse~16r/34, \Trajectinus~TODO, \\
\Wormianus~TODO, \Upsaliensis~TODO}%
\Ballnote{St.~44 is cited in \Gylfaginning\ following \Gylfaginning~51:5--7, which puts the stanza in the context of three winters of war preceding the three winters that make up the Fimblewinter: \emph{En áðr ganga svá aðrir þrír vetr, at þá er um alla verǫld orrostur miklar.  Þá drepa⸗sk brǿðr fyrir á·girni sakar, ok engi þyrmir fǫður eða syni í mann-drǫ́pum eða sifja-sliti.  Svá segir í Vǫlu-spǫ́:} ‘But before them come three other winters when throughout all the world there are great battles.  Then do brothers slay each other for reasons of covetousness, and no man spares his father or son in murder of men or destruction of kinship.  So it says in the Spae of the Wallow:’.
For the Fimblewinter, as described in \Gylfaginning~51:2--4, see \textlink{Vafthrudnismal}[44].

This stanza takes a noticeably moral tone, which may be compared with st.~38 above.  In the view of the wallow, human moral failures accompany the decline and ultimate destruction of the World, and perhaps even hasten its arrival.  This moralising may be seen as a sign of Christian influence (cf.~e.g.~\textlink{Muspilli} 59 for a similar condemnation), but the crimes listed both here and in st.~38 (perjury, murder, adultery, kinslaying, and incest) were all undoubtedly also forbidden in pre-Christian Germanic law.  In the pre-Christian worldview a threat to the established social order was thus also seen as a threat to the very existence of the World itself.}%
\alst{B}rǿðr munu \alst{b}ęrja⸗sk \hld\ ok at \alst{b}ǫnum verða⸗sk, &
munu \edtrans{\alst{s}ystrungar}{the children of sisters}{\Afootnote{\emph{†stystrungar†} \Trajectinus}} \hld\ \edtrans{\alst{s}ifjum spilla}{defile the kinship}{\Bfootnote{I.e., ‘commit incest’, apparently referring to marriages between first cousins, which are prohibited in all mediæval Scandinavian law codes.  Compare related words found in the laws, like \emph{frę́nd-semis-spell} ‘incest’ and especially \emph{sifja-spell} ‘id.’ --- The idea of incest as a sign of later ages is also found in \Rigveda\ 10.10.10a--b (norm. and tr., N. S. Dwibhashyam.): \textoi{Ā́ gʰā tā́ gaccʰān \hld\ úttarā yugā́ni, // yátra jāmáyaḥ \hld\ kr̥ṇávann ájāmi} ‘There shall come indeed those later ages when relatives shall do (acts) not (fit for) relatives.’}}; &
\alst{h}art ’s \edtrans{ï \alst{h}ęimi}{in the Home}{\Afootnote{so \Regius\Hauksbok\Upsaliensis \emph{;} \emph{með hǫlðum} ‘among men’ \RegiusProse\Trajectinus\Wormianus}}, \hld\ \alst{h}ór-dómr mikill, &
\alst{sk}ęggj-ǫld, \alst{sk}alm-ǫld, \hld\ \edtrans{\alst{sk}ildir}{shields}{\Afootnote{add. \emph{’ro} ‘are’ \Regius}} \edtrans{klofnir}{split}{\Afootnote{\emph{klofna} ‘become split’ \Upsaliensis}}, &
\edtrans{\alst{v}ind-ǫld}{wind-age}{\Bfootnote{In \Hauksbok\ the \emph{v} is capitalized, marking the beginning of a new stanza.}}, \alst{v}arg-ǫld, \hld\ \edtrans{áðr}{before}{\Afootnote{\emph{unz} (norm.) ‘until’ \Upsaliensis}} \edtrans{\alst{v}er-ǫld}{man-age}{\Bfootnote{Translated as such since it stands next to various other compounds with the second element \emph{ǫld} ‘age’.
ON \emph{ver-ǫld} is cognate with English “world” (cf.~29/3 above), but in ON the sense ‘cosmos, universe’ is usually expressed by \emph{hęimr} (e.g., in l.~3 of the present stanza).}} \edtrans{stęypi⸗sk}{tumbles down}{\Afootnote{add. after this l. \emph{grundir gjalla \hld\ gífr fljúgandi} (norm.) ‘foundations shrill, fiends flying’ \Hauksbok}} &
\edtrans{mun \edtext{\alst{ę}ngi}{\Afootnote{\emph{†enn†} \Upsaliensis}} maðr \hld\ \alst{ǫ}ðrum þyrma.}{no man will spare another}{\Afootnote{om. \RegiusProse\Trajectinus\Wormianus}}\eva

\bvb Brothers will fight and become each other’s banes; \\
the children of sisters will defile the kinship. \\
It is hard in the Home! Great whoredom, \\
axe-age, sword-age, shields split asunder, \\
wind-age, warg-age; before the man-age tumbles \\
no man will spare another.\evb\evg


\bvg\bva\mssnote{\Regius~2r/32, \Hauksbok~20v/27, \\
\RegiusProse~17r/7, \Trajectinus~TODO, \\
\Wormianus~TODO, \Upsaliensis~18v/7}%
\Ballnote{Sts. 45--54 are cited, with the omission of the refrain-stanza 47, in sequence in \Gylfaginning~51:50.  Note that \Upsaliensis\ only contains 45--46 and 53.}%
\edtext{Lęika \alst{M}íms synir, \hld\ ęn \edtext{\alst{m}jǫtuðr kyndi⸗sk}{\lemma{mjǫtuðr kyndi⸗sk ‘the Measurer advances’}\Bfootnote{The sense seems to be that the destruction swiftly comes to pass through the great predestined calamities, but both words require some elucidation.
ON \emph{mjǫtuðr} ‘the Measurer’ often means merely ‘Fate’ (e.g. \textlink{Sigurdskamma}[71]), but is a m.~agent noun and as such is not entirely impersonal, perhaps rather referring to an overarching fate-principle.
In OS poetry the cognate \emph{metod} means ‘fate’ (\textlink{Heliand}[2][128], 511) and also occurs in the cpds.~\emph{metod-gi·skaft, metod(o)-gi·skapu} ‘shapes of the Measurer’, i.e., ‘decrees of fate’ (\textlink{Heliand} 2190, 2210, 4827; cf.~\emph{ręgin(o)-gi·skapu} ‘shapes of the Reins’, \textlink{Heliand} 2593, 3347, \emph{wurd-gi·skapu} ‘shapes of Weird’, \textlink{Heliand} 127, 3692); in OE poetry \emph{meotod} comes to refer to the Christian God (e.g. \textlink{Cadman} 2).
\emph{kyndi⸗sk} ‘advances’ must be in the pres.~sg.~and is a form of \emph{kynda⸗sk} ‘be kindled, be set ablaze’, which is here understood to refer metaphorically to the swift rising of the calamity.
\CV\ suggests emending to \emph{kynni⸗sk} ‘makes itself known’.  The literal sense ‘is kindled’ could be used to argue for an emendation of \emph{mjǫtuðr} to \emph{mjǫt-viðr} ‘Measuring-Tree’ (cf.~st.~2/4 above), but this would yield an A1 verse with anacrusis, which is not metrically acceptable in \Voluspa.}} &
at \edtext{hinu}{\Afootnote{so \Hauksbok \emph{;} \emph{†en†} \Regius.}} \alst{g}alla \hld\ \edtrans{\alst{G}jallar-horni}{Yellhorn}{\Bfootnote{Or the ‘Horn of Yell’, viz. the river \emph{Gjǫll} ‘Yell’ as mentioned in \textlink{Grimnismal}[28]/6.  The root of the river’s name is the same as English “yell”.}};}{\lemma{Lęika \dots\ Gjallar-horni; ‘Mime’s \dots\ Yell.’}\Afootnote{om. \GylfMS}} &
\alst{h}ǫ́tt blę́ss \alst{H}ęim·dallr, \hld\ \alst{h}orn ’s ȧ lopti; &
\edtrans{\alst{m}ę́lir}{speaks}{\Afootnote{\emph{†mey†} \RegiusProse \emph{;} \emph{†nie†} \Trajectinus}} Óðinn \hld\ við \edtext{\alst{M}ïms}{\Afootnote{\emph{Mïmis} \Trajectinus\Upsaliensis}} hǫfuð; &
\edtext{skelfr \alst{Y}gg·drasils \hld\ \alst{a}skr standandi, &
\alst{y}mr}{\lemma{}\Afootnote{These lines are reversed in \Regius.}} it \alst{a}ldna tré, \hld\ \edtrans{ęn \alst{jǫ}tunn \edtext{losnar}{\Afootnote{add. H41 \Hauksbok.}}}{and the ettin comes loose}{\Afootnote{\emph{ę̇sir ’ro ȧ þingi} \Upsaliensis}}.\eva

\bvb Mime’s sons \ken{ettins} play and the Measurer advances \\
at [the sounding of] the shrill Yellhorn. \\
High blows Homedall; the horn is aloft; \\
Weden speaks with Mime’s head. \\
Ugdrassle’s Ash quakes, standing: \\
the old tree whines and the ettin comes loose.\evb\evg


\bvg\bva\mssnote{\Regius~2v/8, \Hauksbok~20v/30, \\
\RegiusProse~17r/9, \Trajectinus~TODO, \\
\Wormianus~TODO, \Upsaliensis~18v/9}%
\edtrans{Hvat ’s með \alst{ǫ̇}sum? \hld\ hvat ’s með \edtrans{\alst{ǫ}lfum}{Elves}{\Afootnote{\emph{ǫ̇synjum} ‘Ossens’ \Upsaliensis}}}{What is with the Eese? What is with the Elves?}{\Bfootnote{A line indicating cosmic distress.  It also occurs in \textlink{Thrymskvida}[7]/1, where it is answered by \emph{Illt ’s með ǫ̇sum, \hld\ illt ’s með ǫlfum!} ‘’Tis ill with the Eese! ’Tis ill with the Elves!’.  This is certainly the implication of the present stanza.}}? &
\edtext{gnýr \alst{a}llr \alst{Jǫ}tun-hęimr, \hld\ \alst{ę̇}sir ’ro ȧ \edtrans{þingi}{the Thing}{\Bfootnote{The \inx[P]{Thing of the Gods}, where the Gods assemble in times of distress; see note to st.~6/1--2 above.}},}{\lemma{gnýr \dots\ þingi}\Afootnote{om. \Upsaliensis}} &
\alst{st}ynja dvergar \hld\ fyr \edtext{\alst{st}ęin-durum}{\Afootnote{\emph{stęins dyrum} \Upsaliensis \emph{;} \emph{stęin-dyrum} \Hauksbok\Wormianus}} &
\edtext{\edtext{\alst{v}ęgg-bergs}{\Afootnote{\emph{veg-bergs} \Hauksbok\Trajectinus\Wormianus}} \alst{v}ísir}{\Afootnote{om. \Upsaliensis}}. \hld\ \alst{V}ituð ér ęnn eða hvat?\eva

\bvb What is with the Eese? What is with the Elves? \\
All Ettinham roars; the Eese are at the Thing. \\
Dwarfs groan before gates of stone, \\
the lords of the cliff-side. \emph{Know ye yet, or what?}\evb\evg


\bvg\bva\mssnote{\Regius~2v/4, \Hauksbok~20v/32}%
\alst{G}ęyr nú \alst{G}armr mjǫk \hld\ fyr \alst{G}nipa-hęlli, &
\alst{f}ęstr mun slitna, \hld\ ęn \alst{f}reki rinna; &
\alst{f}jǫlð vęit hǫ̇n \alst{f}rǿða, \hld\ \alst{f}ramm sé’k lęngra &
of \alst{r}agna \alst{r}ǫk, \hld\ \alst{r}ǫmm sig-tíva.\eva

\bvb Now Garm barks loudly before the Gnip-caves; \\
the rope will snap and the Wolf will run. \\
She knows much lore; I see further ahead \\
about the mighty Rakes of the Reins, of the Victory-Tews \ken{gods}.\evb\evg


\bvg\bva\mssnote{\Regius~2v/4, \Hauksbok~20v/32, \\
\RegiusProse~17r/11, \Trajectinus~TODO, \\
\Wormianus~TODO}%
\alst{H}rymr ękr austan, \hld\ \alst{h}ęf⸗sk lind fyrir, &
sný⸗sk \alst{Jǫ}rmun·gandr \hld\ ï \alst{jǫ}tun-móði, &
\alst{o}rmr knýr \alst{u}nnir, \hld\ \edtrans{ęn \alst{a}ri hlakkar}{and the eagle screams}{\Afootnote{\emph{ǫrn mun hlakka} ‘the eagle will scream’ \RegiusProse\Trajectinus}}, &
slítr \alst{n}ái \alst{n}ęf-fǫlr; \hld\ \edtrans{\alst{N}agl·far}{Nailfare}{\Bfootnote{The ship on which the fiends sail from the east.  It is described in \Gylfaginning~51:13--16: \emph{Þá geysisk haf’it á lǫnd’in, fyrir því at þá sný⸗sk Mið·garðs-ormr í jǫtun-móð ok sǿkir upp á land’it.  Þá verðr ok þat, at Nagl·far losnar, skip þat, er svá heitir.  Þat er gert af nǫglum dauðra manna, ok er þat fyrir því varnanar vert, ef maðr deyr með ó·skornum nǫglum, at sá maðr eykr mikit efni til skips’ins Nagl·fars, er goð’in ok menn vildi seint, at gert yrði.  En í þessum sę́var-gang flýtr Nagl·far.} ‘Then the sea surges onto the lands because the Middenyardswyrm writhes in ettin-wrath and attacks the land.  Then it also happens that Nailfare comes loose---the ship which is called thus; it is made from the nails of dead men, and therefore it is worth a warning if a man dies with uncut nails, that that man greatly increases the material for the ship Nailfare, which the Gods and men would rather were not done.  But in this motion of the sea Nailfare floats.’

The motif of the Devil rowing a boat made out of improperly discarded nail clippings is widespread in Scandinavian and Finnish folklore \parencite{Krohn1912}, and thus the “nail-ship” cannot be a mere Snorronean invention.  Specific rituals for disposing cut nails (and hair) are found in many Indo-European cultures, and clearly described ramifications for not observing them are also found in Zoroastrianism (\emph{Vidēvdāt} 17.2--3: “[W]hen one arranges and cuts his hair and clips his nails and then lets them fall into holes in the earth or into furrows, [...] demons come forth, and from these improprieties monsters come forth from the earth which mortals call lice and which devour the grain in the fields and the clothes in the closets”) \parencite{Lincoln1977}.  Beyond the Indo-Europeans, Judaism also has a taboo against discarding nail-clippings wherever a pregnant woman might walk, lest they supernaturally cause her to miscarry.

An interesting psychological commonality presents itself in all three traditions, namely that the improper clipping or disposal of nails is thought to threaten the fundamental stability of the cosmos by bringing about the principal fear of each culture.  In Judaism, singularly concerned with the fertility of its adherents (witness its famous ban on masturbation, \emph{Shulchan Arukh, Even HaEzer} 23:1: \texthebrew{ועון זה חמור מכל עבירות שבתורה} ‘more severe than all Torah transgressions’) it is the miscarriage; in the dualistic Zoroastrianism, the proliferation of demons; and in the Germanic religion, with its Wiking Age eschatological obsession (cf.~Introduction to the poem), it is the End Times.}} losnar.\eva

\bvb Rim drives from the east; he holds his shield before him; \\
Ermengand writhes in ettin-wrath. \\
The Wyrm drives forth the waves and the eagle screams, \\
the pale-beak tears at corpses; Nailfare comes loose.\evb\evg


\bvg\bva\mssnote{\Regius~2v/6, \Hauksbok~20v/34, \\
\RegiusProse~17r/13, \Trajectinus~TODO, \\
\Wormianus~TODO}%
\edtrans{\alst{K}jóll}{The ship}{\Bfootnote{Nailfare, as mentioned in the previous st.}} fęrr austan \hld\ \alst{k}oma munu \edtext{Mú·spells &
of \alst{l}ǫg \alst{l}ýðir}{\lemma{Mú·spells \dots\ lýðir ‘Muspell’s subjects’}\Bfootnote{According to \Gylfaginning~4:5--8 \emph{Mú·spell} ‘Muspell’ or \emph{Mú·spells-hęimr} ‘Muspellsham’ is the name of the primordial realm of fire in the South, ruled over by Surt (cf.~note to st.~3 above).

During the destruction the world will be assailed by the nebulous \emph{Mú·spells-synir} ‘Muspell’s sons’ (\textlink{Lokasenna}[42], \Gylfaginning~13:9, 37:19, 51:21, 51:28) or \emph{Mú·spells-męgir} ‘Muspell’s lads’ (\Gylfaginning~13:6, 51:25), who are to be identified with \emph{Mú·spells lýðir} and the \emph{fífl-męgir} ‘devil-lads’ in the pres.~st.

ON \emph{Mú·spell} has West Germanic cognates in OHG \emph{Mú·spilli} (\textlink{Muspilli} 56) and OS \emph{Mú·spelle} (\textlink{Heliand}[31][2591], 4358), both referring to the destruction of the world in a Christian context.  For the etymology see \textlink{Muspilli} 56 n.}}, \hld\ ęn \alst{L}oki stýrir; &
\alst{f}ara \alst{f}ífl-męgir \hld\ með \alst{f}reka allir, &
þęim es \alst{b}róðir \hld\ \alst{B}ý·lęists ï fǫr.\eva

\bvb The ship fares from the east---come will Muspell’s \\
subjects o’er the sea---and Lock steers it. \\
The devil-lads all fare with the Wolf; \\
with them comes the brother of Bilast \ken*{= Lock} along.\evb\evg


\bvg\bva\mssnote{\Regius~2v/10, \Hauksbok~20v/36, \\
\RegiusProse\textsubscript{1}~2r/35, \RegiusProse\textsubscript{2}~17r/15, \\
\Trajectinus\textsubscript{1}~3r/11, \\
\Wormianus~TODO, \Upsaliensis~TODO}%
\Ballnote{This st.~is cited twice in all mss.~of \Gylfaginning\ save \Upsaliensis.  The second citation is in the sequence of sts.~45--54 cited in \Gylfaginning~51:50 (see st.~45 n.), but the first is alone, introduced by \Gylfaginning~4:5--9:

\emph{Þá mę́lti Þriði: „Fyrst var þó sá heimr í suðr-hálfu, er Múspell heitir.  Hann er ljóss ok heitr.  Sú átt er logandi ok brennandi.  Er hann ok ó·fǿrr þeim, er þar eru út-lendir ok eigi eigu þar óðul.  Sá er Surtr nefndr, er þar sitr á lands-enda til land-varnar.  Hann hefir loganda sverð, ok í enda ver-aldar mun hann fara ok herja ok sigra ǫll goð’in ok brenna allan heim með eldi.  Svá segir í Vǫluspǫ́:}

‘Then spoke Third: “Yet, first there was the home in the southern half which is called Muspell.  It is hot and bright.  That domain is blazing and burning.  It is also unpassable to those who are foreigners and do not have their home there.  Surt is he called who sits there at the border of the land to defend the land.  He has a blazing sword, and in the end of the world he will go and harry and vanquish all the Gods and burn all the Home with fire.  So it says in the Spae of the Wallow.’}%
\edtext{\alst{S}urtr}{\Afootnote{\emph{Svartr} \Upsaliensis}} fęrr \alst{s}unnan \hld\ með \alst{s}viga lę́vi, &
\edtrans{skïnn af \alst{s}verði \hld\ \alst{s}ól val-tíva}{from the swords of the slain-Tews \name{gods} [it] shines [like] the sun}{\Bfootnote{A difficult line.  As the sun has been blackened (st.~40) the world is instead lit up by the flashing swords of the Gods as they head into the battle with Surt.  Snorre instead interprets it as a reference to Surt’s sword, \Gylfaginning~51:22--23: \emph{Surtr ríðr fyrst ok fyrir honum ok eptir bę́ði eldr brennandi.  Sverð hans er gott mjǫk; af því skínn bjartara en af sólu.} ‘Surt rides first and both ahead of and behind him is a burning fire.  His sword is very good; from it, it shines brighter than the sun.’

The main difficulty lies in the fact that \emph{val-tíva} may be read either (1) as gen.~pl.~of the attested \emph{val-tívar} ‘the slain-Tews \ken{gods}’ (so \textcite{LaFargeGlossary}; cf.~below st.~59 (n.) and \textlink{Hymiskvida}[1]/1) or %
(2) as gen.~sg.~of an unattested m.~\emph{n}-stem \emph{*val-tívi} ‘the slaughter-Tew \ken*{= Surt}’ (so \CV, Fritzner, and presumably also Snorre).
Reading 1 is adopted here as the simpler explanation because it does not presuppose an otherwise unattested \emph{n}-stem derivative of \emph{týr} ‘tew, god’---especially as Surt is not a god---and because the pl.~\emph{val-tívar} demonstrably occurs elsewhere in the poem.

A second issue arises over whether the gen.~\emph{val-tíva} modifies (a) \emph{sverði} ‘sword’ or (b) \emph{sól} ‘sun’.

In readings 1a and 2a the sense would prob.~be that the flashing sword shines like the sun, reflecting the intensity of the battle.  Although one may expect a pl.~\emph{sverðum} ‘swords’ for 1a, that is not a major problem, since sg.~forms are often used with pl.~possessors in the Old Germanic languages (e.g. \textlink{Harbardsljod}[39]/3).  Case 1b does not seem very likely since ‘sun of the \textsc{gods}’ is not a known kenning-type.  If 2b is adopted we may see a kenning “sun of Surt \ken{fire}”, but “from his sword shines the \ken{fire}” seems excessively repetitive when compared to l.~1.  In conclusion, my preferred reading is 1a.}}; &
\alst{g}rjót-bjǫrg \alst{g}nata, \hld\ ęn \edtrans{\alst{g}ífr rata}{the fiends roam}{\Afootnote{\emph{guðar hrata} ‘the gods(?) stagger’ (wo. doubt corrupt) \Upsaliensis}\Bfootnote{The reading of \Upsaliensis\ is certainly corrupt; the anachronistic m.~pl. ending \emph{-ar} is proof enough, for the word \emph{goð} \char`~\ \emph{guð} ‘gods’ was always neuter in heathen times.}}, &
\edtrans{troða \alst{h}alir \alst{h}ęl-veg}{men tread the Hellway}{\Bfootnote{The Hellway is the road between Earth and the Underworld (Hell) upon which one has to travel after death to reach one’s final resting place (cf.~the story told in \textlink{Helreid}).  The verse may thus be understood as ‘men die off’.

On the other hand, according to \Gylfaginning~51:27, Lock is followed by \emph{allir Heljar sinnar} ‘all Hell’s companions’ who are presumably thought to be some sort of revenants, so that may instead be what is being referred to here.}}, \hld\ ęn \edtrans{\alst{h}iminn klofnar}{the Heaven cracks}{\Bfootnote{As discussed above (st.~3) Heaven is thought of as a solid dome made out of stone (cf.~\textlink{Harbardsljod}[15]/2 n.); its cracking in two is therefore a horrible sight.  Cf.~the prophetic line in Sö 154 Skarpåker (edited below under Runic Poetry from Sweden and Gotland).}}.\eva

\bvb Surt fares from the south with the twig’s betrayer \ken{fire}; \\
from the swords of the slain-Tews \name{gods} [it] shines [like] the sun. \\
Rocky cliffs clash, but the fiends roam; \\
men tread the \inx[P]{Hellway}, but the Heaven cracks.\evb\evg


\bvg\bva\mssnote{\Regius~2v/13, \Hauksbok~20v/37, \\
\RegiusProse~17r/17, \Trajectinus~TODO, \\
\Wormianus~TODO}%
\Ballnote{Sts. 51--53 describe the deaths of the major gods Weden, Thunder, and Free as they fight against their respective enemies: the Fenrerswolf, the Middenyardswyrm, and Surt.  The battle takes place on the great plain Wighride as told in \textlink{Vafthrudnismal}[18].

The fight between the Wolf and Weden is also mentioned in \textlink{Vafthrudnismal}[53], which further mentions his subsequent avenging by Wider, and \textlink{Grimnismal}[24]/3--4, which mentions that he is aided by the Oneharriers.

The relevant part of the battle is described in \Gylfaginning~51:34--38: \emph{Ę́sir her-vę́ða sik ok allir Ein-herjar ok sǿkja fram á vǫllu’na.  Ríðr fyrstr Óðinn með gull-hjalm’inn ok fagra brynju ok geir sinn, er Gungnir heitir; stefnir hann móti Fenris-ulf, en Þórr fram á aðra hlið honum, ok má hann ekki duga honum, því at hann hefir fullt fang at berja⸗sk við Mið·garðs-orm.  Freyr ber⸗sk móti Surti, ok verðr harðr sam-gangr, áðr Freyr fellr.  Þat verðr hans bani, er hann missir þess ins góða sverðs, er hann gaf Skírni.} ‘The Eese and all the Oneharriers clothe themselves for war and rush forth on the plains.  Weden rides first with the golden helmet and fair byrnie and his spear which is called Gungner; he faces against the Fenrerswolf, but Thunder goes forth on one of his sides; and he cannot avail him, for he has his hands full in fighting against the Middenyardswyrm.  Free fights against Surt and it becomes a hard-fought struggle before Free falls.  It will be his bane that he is missing the good sword which he gave Shirner.’}%
Þȧ kømr \edtrans{\alst{H}lïnar \hld\ \alst{h}armr annarr}{Line’s second sorrow}{\Bfootnote{The first sorrow being the death of Balder.  Line (\emph{Hlïn}) is described in \Gylfaginning~35:19 as a minor goddess \emph{sett til gę́tslu yfir þeim mǫnnum, er Frigg vill forða við háska nǫkkurum} ‘placed to watch over those men whom Frie wants to save from any particular danger’.  In spite of this, almost all modern translators and editors have understood Line as synonymous with Frie and posited that her existence as a distinct goddess is a systematizing invention of \Gylfaginning.  \textcite{Hopkins2017} convincingly argues that this need not be the case.  As Frie’s servant put in charge of protecting those dearest to her mistress, Line’s two sorrows would consist in her failure to protect both Frie’s son and husband.}} framm, &
es \alst{Ó}ðinn fęrr \hld\ við \alst{u}lf vega, &
---\edtrans{ęn \alst{b}ani \alst{B}ęlja \hld\ \alst{b}jartr at Surti}{but the bane of Bellower \ken*{= Free}, bright, towards Surt}{\Bfootnote{In this inset line the wallow relates the death of Free; cf.~note to ALL.
Free’s obscure fight against Bellower (ON \emph{Bęli}) is mentioned in \Gylfaginning~37:16--19, which also contains the important detail that Free will be without his sword:

\emph{„Þessi sǫk er til þess, er Freyr var svá vápn-lauss, er hann barði⸗sk við Belja ok drap hann með hjartar-horni.“ Þá mę́lti Gangleri: „Undr mikit, er því⸗líkr hǫfðingi sem Freyr er vildi gefa sverð, svá at hann átti eigi annat jafn-gott.  Geysi-mikit mein var honum þat, þá er hann barði⸗sk við þann, er Beli heitir.  Þat veit trúa mín, at þeirar gjafar myndi hann þá iðra⸗sk.“ Þá svarar Hárr: „Lítit mark var þá at, er þeir Beli hittu⸗sk.  Drepa mátti Freyr hann með hendi sinni.  Verða mun þat, er Frey mun þykkja verr við koma, er hann missir sverðs’ins, þá er Múspells synir fara at herja.“}

‘“This event is the reason for why Free was so unarmed when he fought against Bellower and slew him with a hart’s antler.”  Then spake Gangler: “It is a great wonder that such a chieftain as Free is would give away his sword for which he had nothing as good.  It was a huge harm for him when he fought with the one called Bellower.  It is my true belief that he would come to regret that gift.”  Then answers High: “It was of little harm when he and Bellower met each other.  Free could have slain him with his hand.  Yet a thing will come to pass which will seem far worse to Free: that he is missing his sword when the Sons of Muspell come harrying.”’

Cf.~\textlink{Skirnismal}[8], \textlink{Lokasenna}[42], and \Gylfaginning~37:1--15 (edited as an appendix to \textlink{Skirnismal}).}}--- &
þȧ mun \alst{F}riggjar \hld\ \edtext{\alst{f}alla}{\Afootnote{om., but add.~above line in same hand \RegiusProse}} \edtext{angan}{\Afootnote{so \Hauksbok\RegiusProse\Trajectinus\Wormianus \emph{;} \emph{angantyr} \Regius}}.\eva

\bvb Then comes \inx[P]{Line}’s second sorrow to pass, \\
when Weden goes to fight against the Wolf \\
---but the bane of Bellower \ken*{= Free}, bright, towards Surt--- \\
then will Frie’s beloved \ken*{= Weden} fall.\evb\evg


\bvg\bva\mssnote{\Regius~2v/15, \RegiusProse~17r/19, \\
\Trajectinus~TODO, \Wormianus~TODO}%
\edtrans{Þȧ kømr hinn \alst{m}ikli \hld\ \alst{m}ǫgr \edtrans{Sig-fǫður}{Syefather}{\Bfootnote{‘The Father of Victory’, a name found only here and in \textlink{Lokasenna}[58]/4, likewise in the context of the battle with the Wolf.}}}{Then comes the great son of Syefather}{\Afootnote{\emph{Gęngr Óðins sonr \hld\ við ulf vega} ‘Weden’s son goes the Wolf to fight’ \RegiusProse\Trajectinus\Wormianus.}}, &
\alst{V}íð·arr \edtext{\alst{v}ega}{\Afootnote{\emph{of veg} \RegiusProse\Trajectinus\Wormianus}} \hld\ at \alst{v}al-dýri; &
lę́tr \alst{m}ęgi \edtrans{Hveðrungs}{Whethring’s}{\Bfootnote{Gen.~of \emph{Hveðrungr}, an obscure name for \inx[P]{Lock}, whose son is the Fenrerswolf.  The name otherwise appears as a name of Lock in a kenning for Hell (\Ynglingatal\ 24: \emph{Hveðrungs mę́r} ‘Whethring’s daughter’).  It is also found in a later antiquarian list (thule) of names of individual Ettins (Þul \emph{Jǫtna I} 1, \Skp\ 3), and in a thule of names for Weden (Þul \emph{Óðins} 5, \Skp\ 3), but both of these may derive from misinterpretations of the present stanza.}} \hld\ \alst{m}und of standa &
\alst{h}jǫr til \alst{h}jarta; \hld\ þȧ ’s \alst{h}efnt fǫður.\eva

\bvb Then comes the great son of \inx[P]{Syefather}, \\
Wider, to fight that slaughter-beast. \\
He lets his hand through \inx[P]{Whethring}’s lad \ken*{= the Wolf} \\
drive the sword to the heart---then is the father avenged!\evb\evg


\bvg\bva\mssnote{\Regius~2v/17, \Hauksbok~20v/41, \\
\RegiusProse~17r/20, \Trajectinus~TODO, \\
\Wormianus~TODO}%
\Ballnote{Thunder and the Middenyardswyrm kill each other, fulfilling their previous abortive encounter during the Fishing Trip (see \textlink{Hymiskvida}).  From a comparative perspective, it is notable that the Norse religion makes the ancient \emph{Chaoskampf}, the motif of a Storm God slaying a watery Serpent as part of Creation, a key part of its eschatology instead.  It is surely for this reason that the Wyrm is not slain during the Fishing Trip (as it must have been originally, and as it actually is in \Husdrapa, one of the early Scaldic retellings; see Introduction to \textlink{Hymiskvida}), since it must remain alive to fight at the Rakes of the Reins. ---

This st.~is paraphrased in \Gylfaginning~51:41--42: \emph{Þórr berr banaorð af Mið·garðs-ormi ok stígr þaðan braut níu fet. Þá fellr hann dauðr til jarðar fyrir eitri því, er ormr’inn blę́ss á hann.} ‘Thunder bears the bane-word from the Middenyardswyrm and strides nine steps away from it.  Then he falls dead to the earth from the venom (\emph{ęitri}) which the Wyrm blows on him.’  Cf.~also st.~H48 below.}%
\Aallnote{The present form of the stanza is an amalgamation of all three recensions (\Regius, \Hauksbok, and \GylfMS) based most closely on \Hauksbok\ and \GylfMS, which have the last 3 lines in the same order.  \Regius\ has the lines in the order 1, 5, 4, 2, 3 and inserts an additional line between 1 and 5.}%
\edtrans{\edtrans{Þȧ kømr}{Then comes}{\Afootnote{\emph{Gęngr} ‘Goes’ \RegiusProse\Trajectinus\Wormianus}} hinn \alst{m}ę́ri \hld\ \alst{m}ǫgr Hlǫðynjar,}{Then comes the famed lad of Lathyn}{\Afootnote{om. \Hauksbok \emph{;} add. \emph{gęngr Óðins sonr við ulf (\emph{pro} *orm) vega.} ‘Weden’s son goes the Wolf (error for ‘Wyrm’) to fight.’ \Regius.}} &
\edtrans{gęngr \alst{f}et níu \hld\ \edtext{\alst{F}jǫrgynjar}{\linenum{|2|||3}\lemma{Hlǫðynjar, Fjǫrgynjar ‘Lathyn, Firgyn’}\Bfootnote{Both names for the \inx[P]{Earth}, who is Thunder’s mother.}} burr}{nine steps takes Firgyn’s son}{\Afootnote{om. \RegiusProse\Trajectinus\Wormianus.}} &
\alst{n}ęppr frȧ \alst{n}aðri \hld\ \edtrans{\alst{n}íðs ȯ·kvíðnum}{unfearing of enmity}{\Bfootnote{Dat.~sg., modifying \emph{naðri} ‘adder, serpent’.  Numerous eds.~emend to \emph{*ȯ·kvíðinn} (nom.~sg.), referring to Thunder.}}; &
\edtrans{munu \alst{h}alir allir \hld\ \edtext{\alst{h}ęim-stǫð}{\Afootnote{\emph{hęim-†steið†}}} ryðja}{All men will clear their homesteads}{\Bfootnote{After Thunder is slain the realm of men is no longer habitable.  Cf.~\textlink{Harbardsljod}[23], \textlink{Thrymskvida}[18].}} &
\edtext{es af \alst{m}óði drepr}{\Afootnote{\emph{drepr hann af móði} \Regius}} \hld\ \edtrans{\alst{M}ið·garðs Véurr}{Middenyard’s Wighward}{\Bfootnote{Thunder, literally “the Sanctuary-Defender of Middenyard”, a fitting kenning.}}.\eva

\bvb Then comes the famed son of \inx[P]{Lathyn} \ken*{= Thunder}. \\
Nine steps takes \inx[P]{Firgyn}’s son, \\
pained, away from the Adder unfearing of enmity. \\
All men will clear their homesteads \\
when Middenyard’s Wighward strikes out of wrath.\evb\evg


\bvg\bva\mssnote{\Regius~2v/20, \Hauksbok~21r/1, \\
\RegiusProse~17r/22, \Trajectinus~TODO, \\
\Wormianus~TODO, \Upsaliensis~18v/10}%
\Ballnote{This st.~is verbally alluded to in Arn \emph{Þorfdr} 24/1--2 (\Skp\ 2, ca.~1060 \ce) in an instance of \emph{adynaton}:
\emph{Bjǫrt verðr sól at svartri; \hld\ søkkr fold ï mar døkkvan; / brestr ęrfiði Austra; \hld\ allr glymr sę́r ȧ fjǫllum} ‘Bright, the sun turns black; the land sinks into the dark sea; // Eastre’s toil \ken{heaven} shatters; all the sea roars over the fells’.
Another allusion is possibly found in SnSt \emph{Ht} 13/2, 4 (\Skp\ 3), which may, however, also be a reference to the Biblical deluge: \emph{Stóð sę́r ȧ fjǫllum; [...] skaut jǫrð ór gęima.} ‘The sea stood over the fells; the earth shot up from the ocean.’}%
\alst{S}ól \edtrans{tér}{doth}{\Afootnote{\emph{mun} ‘will’ \GylfMS}} \alst{s}ortna, \hld\ \edtrans{\edtext{\edtrans{\alst{s}ígr}{sinks}{\Afootnote{\emph{søkkr} \RegiusProse\Trajectinus\Wormianus}} fold ï mar}{\Afootnote{\emph{†sigrfolldinar†} (for \emph{*sígr fold ï mar}) \Upsaliensis,
cf.~\textcite[194, 198 n. 25]{Males2023Textual}}}}{the land sinks into the sea}%
{\Bfootnote{The stemmatically excluded reading \emph{søkkr} ‘sinks’ is probably influenced by Arn \emph{Þorfdr} 24/1b (\Skp\ 2).}}, &
\alst{h}verfa af \alst{h}imni \hld\ \alst{h}ęiðar stjǫrnur; &
gęisar \alst{ęi}mi \hld\ \edtrans{við}{from}{\Afootnote{\emph{ok} ‘and’ \Hauksbok\GylfMS}} \edtrans{\edtext{\alst{a}ldr-nara}{\Afootnote{\emph{aldr-nari} (nom.~sg.) \Hauksbok\GylfMS}}}{nourisher of life}{\Bfootnote{A rare \textsc{fire}-kenning only found here and in the \inx[C]{thule}[thules] or poetic lists of synonyms (Þul \emph{Elds} 3 in \Skp\ 3). --- It is notable that fire has two conflicting natures: it nourishes life but it also consumes it; here a kenning reflecting the former is used in the context of the latter.}}; &
lęikr \alst{h}ǫ́r \alst{h}iti \hld\ við \alst{h}imin sjalfan.\eva

\bvb The Sun doth blacken; the land sinks into the sea; \\
from heaven fade the shining stars. \\
Smoke billows from the nourisher of life \ken{fire}; \\
the high heat licks against the very heaven.\evb\evg


\bvg\bva\mssnote{\Regius~2v/22, \Hauksbok~21r/2}%
\Ballnote{With the sinking of the earth into the sea and the blackening of the heaven, the destruction has reached its climax, as signalled by the final repetition of the refrain stanza.  To indicate the dramatic pause, which was likely indicated in some way in the original oral performance, a page break has been inserted.}%
\alst{G}ęyr nú \alst{G}armr mjǫk \hld\ fyr \alst{G}nipa-hęlli, &
\alst{f}ęstr mun slitna, \hld\ ęn \alst{f}reki rinna; &
\alst{f}jǫlð vęit hǫ̇n \alst{f}rǿða, \hld\ \alst{f}ramm sé’k lęngra &
of \alst{r}agna \alst{r}ǫk, \hld\ \alst{r}ǫmm sig-tíva.\eva

\bvb Now Garm barks loudly before the Gnip-caves; \\
the rope will snap and the Wolf will run. \\
She knows much lore; I see further ahead \\
about the mighty Rakes of the Reins, of the Victory-Tews \ken{gods}.\evb\evg

\clearpage
%\pagecolor{black}
%\thispagestyle{empty}

%\null
%\vspace*{\fill}

	%\color{white}
	%\begin{center}
	%\Huge Darkness.
	%\end{center}

%\newpage
%\nopagecolor
%\color{black}

\bvg\bva\mssnote{\Regius~2v/23, \Hauksbok~21r/4}%
\Ballnote{Sts. 56--59 are paraphrased in \Gylfaginning~53:1--9:

\emph{Þá mę́lti Gangleri: „Hvárt lifa nǫkkur goð’in þá, eða er þá nǫkkur jǫrð eða himinn?“ Hárr segir: „Upp skýtr jǫrðu’nni þá ór sę́’num, ok er þá grǿn ok fǫgr. Vaxa þá akrar ó·sánir. Víð·arr ok Váli lifa, svá at eigi hefir sę́r’inn ok Surta-logi grandat þeim, ok byggja þeir á Iða-velli, þar sem fyrr var Ǫ́s·garðr, ok þar koma þá synir Þórs, Móði ok Magni, ok hafa þar Mjǫllni. Því nę́st koma þar Baldr ok Hǫðr frá Heljar, setja⸗sk þá allir samt, ok tala⸗sk við, ok minna⸗sk á rúnar sínar, ok rǿða of tíðendi þau, er fyrrum hǫfðu verit, of Mið·garðs-orm ok of Fenris-ulf.  Þá finna þeir í grasi’nu gull-tǫflur þę́r, er ę́sir’nir hǫfðu átt.“}

‘Then spoke Gangler: “Do any Gods live then, or is there then any earth or heaven?”  High says: “The earth then shoots up from the sea and it is then green and fair.  Then acres grow unsown.  Wider and Wonnel live, for the sea and Surt’s flame have not scathed them, and they settle on the Idewold where Osyard once stood.  And then come thither Thunder’s sons, Mood and Main, and they have Millner with them.  Next come Balder and Hath out of Hell.  Then they all sit down together and make speech and think back on their runes and speak about the events of antiquity, about the Middenyardswyrm and the Fenrerswolf.  Then they find in the grass those golden game-bricks which the Eese had owned.”’

After this passage \Gylfaginning\ cites \textlink{Vafthrudnismal}[51].  For more traditions of the world to come, cf.~\textlink{Vafthrudnismal}[44]--47.}%
Sér hǫ̇n \alst{u}pp koma \hld\ \edtrans{\alst{ǫ}ðru sinni}{a second time}{\Bfootnote{The first time probably being the lifting of the Earth in st.~4.}} &
\alst{jǫ}rð ór \alst{ę́}gi \hld\ \alst{i}ðja-grø̇na; &
\edtext{\alst{f}alla \alst{f}orsar, \hld\ \alst{f}lýgr ǫrn yfir, &
sá’s ȧ \alst{f}jalli \hld\ \alst{f}iska vęiðir}{\lemma{falla \dots\ vęiðir ‘Torrents \dots\ fish.’}\Bfootnote{The unusually vivid description paints a beautiful picture of the mountain-climate of Norway and Iceland.  The sense is that the serene beauty of nature is restored from the fire of the destruction.}}.\eva

\bvb She sees coming up for a second time \\
the Earth from the ocean, ever green anew. \\
Torrents fall; o’er them flies the eagle \\
which in the fells catches fish.\evb\evg


\bvg\bva\mssnote{\Regius~2v/24, \Hauksbok~21r/5}%
\edtrans{Finna⸗sk}{find each other}{\Afootnote{\emph{hitta⸗sk} \Hauksbok}\Bfootnote{The \Hauksbok\ reading provides closer parallelism with 7/1 above, but for the same reason it may also have replaced an earlier \emph{finna⸗sk}, as found in \Regius.}} \alst{ę̇}sir \hld\ ȧ \alst{I}ða-vęlli &
ok umb \edtrans{\alst{m}old-þinur}{Earth-cord}{\Bfootnote{Cf.~the kenning for the Middenyardswyrm in \Husdrapa\ 4: \emph{stirð-þinull storðar} ‘the stiff cord of the land \ken*{= Middenyardswyrm}’.}} \hld\ \alst{m}ǫ́tkan dø̇ma, &
\edtrans{ok \alst{m}inna⸗sk þar \hld\ ȧ \alst{m}ęgin-dȯma}{and there think back on mighty verdicts}{\Afootnote{so \Hauksbok \emph{;} om. \Regius}} &
ok ȧ \alst{F}imbul-týs \hld\ \alst{f}ornar ru̇nar.\eva

\bvb The Eese find each other on the Idewold \\
and of the mighty Earth-Cord \ken*{= Middenyardswyrm} speak, \\
and there think back on mighty verdicts \\
and on Fimble-Tew’s \name{= Weden’s} ancient runes.\evb\evg


\bvg\bva\mssnote{\Regius~2v/26, \Hauksbok~21r/7}%
\Ballnote{A fine literary device.  In st.~8 the golden age of the Eese, exemplified by their playing board games, was spoiled by the three ettin-women.  The rediscovery of the golden board game then betokens a new golden age.}%
Þar munu \edtrans{\alst{ę}ptir}{again}{\Afootnote{\emph{ę̇sir} ‘the Eese’ \Hauksbok}} \hld\ \alst{u}ndr-sam⸗ligar &
\alst{g}ollnar tǫflur \hld\ ï \alst{g}rasi \edtrans{finna⸗sk}{be found}{\Afootnote{\emph{finna} ‘find’ \Hauksbok}}, &
þę́r’s \edtext{\emph{\alst{ę̇}sir}}{\lemma{\emph{ę̇sir}}\Afootnote{om.~\Regius\Hauksbok}\lemma{\emph{ę̇sir} ‘the Eese’}\Bfootnote{The verb \emph{hǫfðu} ‘had’ requires a plural subject, and this is found in the paraphrase of this st.~in \Gylfaginning~53:9: \emph{Þá finna þeir í grasi’nu gull-tǫflur þę́r, er ę́sir’nir hǫfðu átt} (see st.~56 n. above).}} ï \alst{á}r-daga \hld\ \alst{á}ttar hǫfðu.\eva

\bvb There will again the wondersome \\
golden game-bricks in the grass be found, \\
which in days of yore the Eese had owned.\evb\evg


\bvg\bva\mssnote{\Regius~2v/28, \Hauksbok~21r/9}%
Munu \alst{ȯ}·sánir \hld\ \alst{a}krar vaxa, &
\edtrans{\alst{b}ǫls}{the bale}{\Bfootnote{The evil of Hath’s slaying Balder (cf.~sts.~31--34 above).  It will be forgotten as the two brothers reconcile and live alongside each other.}} mun alls \alst{b}atna, \hld\ mun \alst{B}aldr koma; &
búa \alst{H}ǫðr ok Baldr \hld\ \alst{H}ropts sig-toptir, &
\alst{v}ęl \edtrans{\alst{v}al-tívar}{the slain-Tews}{\Bfootnote{A rare epithet otherwise only found in 50/2 above and \textlink{Hymiskvida}[1]/1; it is not found in Snorre’s Edda or in Scaldic poetry.  The first element is most likely \emph{valr} ‘the battle-slain’.  This term probably originally referred exclusively to Hath and Balder, since just these are the two slain gods; it would later have been semantically widened to refer to the Gods in general (so \textlink{Voluspa}[50]/2, \textlink{Hymiskvida}[1]/1).}}. \hld\ \alst{V}ituð ér ęnn eða hvat?\eva

\bvb Unsown will acres grow; \\
the bale will all be bettered; Balder will come. \\
Hath and Balder settle Roft’s \name{= Weden’s} victory-plots \\
well, the slain-Tews. \emph{Know ye yet, or what?}\evb\evg


\bvg\bva\mssnote{\Regius~2v/30, \Hauksbok~21r/11}%
Þȧ kná \alst{H}ø̇nir \hld\ \edtrans{\alst{h}laut-við kjósa}{choose the leat-wood}{\Bfootnote{Foresee the future by means of lots, specifically twigs drenched in the blood of beasts killed during the bloot; cf.~\textlink{Hymiskvida}[1].  According to Tacitus (\emph{Germania}, ch. 10) such divination was done \emph{publice consultetur, sacerdos civitatis, sin privatim, ipse pater familiae} ‘in public questions [by] the priest of the particular state, in private [by] the father of the family’.  In later Scandinavian society the positions of high Gid (\emph{goði}) and chieftain or king were largely synonymous, and so this would mean that in the world to come, the peaceful Heener takes over the role of the chief of the Gods after the more warlike Weden.  I thank John Newman for this observation.}} &
ok \alst{b}urir \alst{b}yggva \hld\ \edtrans{\alst{b}rǿðra tvęggja}{of the two brothers}{\Bfootnote{The present translation takes \emph{tvęggja} to be the gen.~pl.~of \emph{tvęir} ‘two’.  The two brothers are presumably Hath and Balder, mentioned in the previous stanza.  Since the original mss.~do not capitalize proper nouns, we could also read \emph{brǿðra Tvęggja} ‘the brothers of Tway \name{= Weden}’, but although Weden’s brothers are attested in \Gylfaginning~6:10 as \inx[P]{Will} and \inx[P]{Wigh} (cf.~st.~4/1 above) they are never said to have children.}} &
\alst{v}ind-hęim \alst{v}íðan. \hld\ \alst{V}ituð ér ęnn eða hvat?\eva

\bvb Then does Heener choose the \inx[C]{leat}-wood, \\
and the sons of the two brothers settle \\
the wide wind-home \ken{sky/heaven}. \emph{Know ye yet, or what?}\evb\evg


\bvg\bva\mssnote{\Regius~2v/31, \Hauksbok~21r/12, \\
\RegiusProse~5v/12, \Trajectinus~TODO, \\
\Wormianus~TODO, \Upsaliensis~TODO}%
\Ballnote{This stanza is cited immediately following \Gylfaginning~17:10--12, which read: \emph{Á sunnan-verðum himins enda er sá salr er allra er fegrstr ok bjartari en sól’in, er Gim·lé heitir.  Hann skal standa þá er bę́ði himinn ok jǫrð hefir farit⸗sk, ok byggja þann stað góðir menn ok rétt-látir of allar aldir.  Svá segir í Vǫlu-spǫ́:} ‘At the southern end of heaven is the hall which is the fairest of all and brighter than the sun, which is called Gemlee.  It shall stand when both the heaven and the earth have been destroyed, and that place is settled by good men and righteous throughout all ages.  So it says in the Spae of the Wallow.’

\Gylfaginning~17:13--16 continues: \emph{Þá mę́lti Gangleri: „Hvat gę́tir þess staðar, þá er Surta-logi brennir himin ok jǫrð?“ Hárr segir: „Svá er sagt, at annarr himinn sé suðr ok upp frá þessum himni, ok heitir sá And·langr, en inn þriði himinn sé enn upp frá þeim, ok heitir sá Víð·bláinn, ok á þeim himni hyggjum vér þenna stað vera.  En Ljós-alfar einir, hyggjum vér, at nú byggvi þá staði.“} ‘Then spoke Gangler: “What protects that place when the flame of Surt burns the heaven and the earth?” High says: “So it is said, that there is another heaven southwards and upwards from this heaven and that one is called Endlong, but the third heaven is even further upwards from them and that one is called Wideblown, and on that heaven we think this place is.  But Light-Elves alone, we think, do now inhabit those places.”’}%
\alst{S}al \edtrans{sér hǫ̇n}{she sees}{\Afootnote{\emph{vęit’k} ‘I know’ \GylfMS}} standa \hld\ \alst{s}ólu fęgra, &
\edtrans{\alst{g}olli þakðan}{thatched with gold}{\Afootnote{\emph{golli bętra} ‘better than gold’ \RegiusProse\Trajectinus}}, \hld\ ȧ \edtext{\alst{G}im·lé\emph{i}}{\Afootnote{metr. emend.~by restoration of hiatus; \emph{Gim·lé} \Regius\Hauksbok\GylfMS}}; &
\edtrans{þar}{there}{\Afootnote{\emph{þann} ‘[in] that [hall]’ \Trajectinus\Wormianus}} skulu \alst{d}yggvar \hld\ \alst{d}róttir byggva &
ok umb \alst{a}ldr-daga \hld\ \alst{y}nðis njóta.\eva

\bvb A hall she sees standing fairer than the sun, \\
thatched with gold on Gemlee. \\
There faithful folk shall settle \\
and in their days of life enjoy delights.\evb\evg


\bvg\bva\mssnote{\Regius~3r/2, \Hauksbok~21r/15}%
\Ballnote{The function of the final st.~is disputed.  Earlier eds.~see the last half-line as a reference to Nithehewer sinking down into the abyss, permanently ridding the world of evil and resulting in an everlasting world-to-come---a clearly Christian conception.  Yet this interpretation requires either emending \emph{hǫ̇n} ‘she’ in l.~4 to \emph{hann} ‘he’ (so \cite{FinnurEdda}) or removing it entirely (Sijmons and Gering, III, 77--78).

I instead agree with \textcite{Schjodt1981} in seeing the last half-line as part of the frame of the wallow’s recitation and taking the rest of the stanza as a reference to the pre-Christian Indo-European cyclical conception of the world.  The presence of Nithehewer illustrates that the next world is still not entirely free from ills, for although most of the fiends have been slain in the great battle, a seed of evil still remains.  The Dragon will infest the \emph{axis mundi} (\textlink{Grimnismal}[36]) in the next world also, and the fact that he still has corpses to eat further shows that evil in mankind has not been eradicated (st.~38 above).
Since there is still evil in the next world it too must at some point fall, be destroyed, and be reborn in the endless cosmic cycle; see Introduction to the poem above.}%
Þar kømr hinn \alst{d}immi \hld\ \alst{d}ręki fljúgandi, &
\alst{n}aðr frȧnn \alst{n}eðan \hld\ frȧ \alst{N}iða-fjǫllum; &
berr sér ï \alst{f}jǫðrum \hld\ ---\alst{f}lýgr vǫll yfir--- &
\alst{N}íð·hǫggr \alst{n}ái; \hld\ \edtrans{\alst{n}ú mun hǫ̇n søkkva⸗sk}{Now will she sink!}{\Bfootnote{The spae is concluded and the wallow, referring to herself in the third person, descends back down into the grave whence Weden woke her.  Cf.~\textlink{Helreid}[14]/4b, the very last half-line of that poem: \emph{søkkst-u, gýgjar-kyn} ‘sink, thou gow’s kin!’}}.“\eva

\bvb Then comes the dusky Dragon flying, \\
the gleaming Adder up from the \inx[P]{Nithfells}. \\
He carries in his feathers---he flies o’er the field--- \\
Nithehewer, corpses.  Now will she sink!”\evb\evg

\sectionline

\newpage

\section{Appendix}

\subsection{Dwarf-tallies (sts.~11--15)}

The following sts.~(11--15) contain two or three originally distinct lists of dwarf-names.  They interrupt the flow of the whole poem so much that they have been moved to the Appendix.  Their nature as a later insertion is in any case confirmed by their faulty meter \parencite[373]{Males2025Hlod}.

That the lists themselves are originally distinct from each other is seen by the repetition of names (13/4, 15/4: \emph{Ęikin·skjaldi} ‘Oakenshield’; 11/4, 15/3: \emph{Ái} ‘Great-grandfather’) and the existence of two formulaic conclusions (12/3--4, 15/6--7).  Sts.~11--13, having no repeated names, seem to belong together.  If this is the case, st.~12, which contains the formulaic conclusion to the list, should probably switch places with 13.

Sts.~14--15 form a second group, having an introduction and a conclusion which both mention the dwarf Loffer.

NOTE: The variants from the four \Gylfaginning~manuscripts have not yet been collated.%TODO.

\bvg\bva[11]\mssnote{\Regius~1r/23, \Hauksbok~20r/17, \\
\RegiusProse~TODO, \Trajectinus~TODO, \\
\Wormianus~TODO, \Upsaliensis~TODO}%
\Ballnote{Sts.~11--12 are cited in \Gylfaginning~14:16, introduced by the words \emph{Ok þessi segir hón nǫfn þeira} ‘And in this way she says their names’.}%
\alst{N}ýi ok \alst{N}iði, \hld\ \edtext{\alst{N}orðri, Suðri, &
\alst{Au}stri, Vestri}{\lemma{Norðri, Suðri, Austri, Vestri ‘Norther and Souther, Easter and Wester’}\Bfootnote{According to \Gylfaginning~8:10, these four dwarfs were placed by the gods to hold up the Heaven, Yimer’s skull, one in each cardinal direction; cf.~st.~3/4 n.~above.}}, \hld\ \alst{A}l·þjófr, Dvalinn, &
\alst{B}ífurr, \alst{B}ǫ́furr, \hld\ \alst{B}ǫmburr, Nóri, &
\alst{Ȧ}nn ok \alst{Ȧ}narr, \hld\ \alst{Á}i, Mjǫð·vitnir.\eva

\bvb New and Nithe, Norther and Souther, \\
Easter and Wester, Allthief, Dwollen, \\
Biber, Babber, Bamber, Noor, \\
Own and Owner, Great-grandfather, Meadwitner.\evb\evg


\bvg\bva[12]\mssnote{\Regius~1r/25, \Hauksbok~20r/18, \\
\RegiusProse~TODO, \Trajectinus~TODO, \\
\Wormianus~TODO, \Upsaliensis~TODO}%
\alst{V}ęigr ok Gand·alfr, \hld\ \alst{V}ind·alfr, Þráinn, &
\alst{Þ}ękkr ok \alst{Þ}orinn, \hld\ \alst{Þ}rór, Vitr ok Litr, &
\alst{N}ár ok \alst{N}ý·ráðr, \hld\ ---\edtext{\edtext{\alst{n}ú hęf’k dverga}{\Afootnote{om. \GylfMS}}--- &
\alst{R}ęginn ok \alst{R}áð·sviðr \hld\ ---\edtext{\alst{r}étt of talða}{\Afootnote{om. \GylfMS}}.}{\lemma{nú hęf’k dverga \dots\ rétt of talða ‘now have I the dwarfs \dots\ rightly tallied’}\Bfootnote{An inset clause; since it does not list the names of any dwarfs, it was omitted by Snorre in \Gylfaginning.}}\eva

\bvb Wey and Gandelf, Windelf, Thrown, \\
Thetch and Thorn, Threw, Wit and Lit, \\
Nee and Newred---now have I the dwarfs--- \\
Rain and Redswith---rightly tallied.\evb\evg


\bvg\bva[13]\mssnote{\Regius~1r/28, \Hauksbok~20r/20, \\
\RegiusProse~TODO, \Trajectinus~TODO, \\
\Wormianus~TODO, \Upsaliensis~TODO}%
\alst{F}íli, Kíli, \hld\ \alst{F}undinn, Náli, &
\alst{H}ępti, Víli, \hld\ \alst{H}annarr, \edtext{Svíurr}{\Afootnote{\emph{†Sviðrr†, / Nár ok Náinn, \hld\ Nípingr, Dáinn / Billingr, Brúni \hld\ Bildr ok Búri} ‘\dots\ Nee and Nowen, Nipng, Dowen, Billing, Brown, Bild and Boor’ \Hauksbok}}, &
\alst{F}rár, Horn·bori, \hld\ \alst{F}rę́gr ok Lȯni, &
\alst{Au}r·vangr, \alst{Ja}ri, \hld\ \alst{Ęi}kin·skjaldi.\eva

\bvb Filer, Chiler, Found and Needler, \\
Hefter, Wiler, Hanner, Swigher, \\
Fraw, Hornborer, Fray and Looner, \\
Earwong, Earer, Oakenshield.\evb\evg


\bvg\bva[14]\mssnote{\Regius~1r/30, \Hauksbok~20r/22}%
\Ballnote{This st.~is paraphrased in \Gylfaginning~14:18 by the words: \emph{En þessir kómu frá Svarins-haugi til Aur-vanga á Jǫru-vǫllu, ok er kominn þaðan Lofarr.} ‘But these ones came from Swornshigh to the Earwongs on the Erwolds, and from them Loffer is come.’  After this passage, \Gylfaginning\ cites st.~15/3--5.}%
\edtext{Mál es \alst{d}verga \hld\ ï \alst{D}valins liði &
\alst{l}jȯna kindum \hld\ til \alst{L}ofars tęlja}{\lemma{Mál es \dots\ tęlja ‘It is time \dots\ of men’}\Bfootnote{A variant of a standard genealogical introduction; cf.~\Haleygjatal\ 1: \emph{meðan hans ę́tt \dots\ til goða tęljum} ‘while we tally his line \dots\ back to the Gods’.  The patrilineal line of dwarfs is to be counted back to their progenitor, Loffer.  This possibly disagrees with st.~10 where Moodsowner is said to be the foremost (and presumably the oldest) of the dwarfs, and Loffer is not mentioned, but consistency in such details can probably not be expected.}}, &
\edtext{þęir}{\Afootnote{\emph{þeim} \Hauksbok}} es \alst{s}óttu \hld\ frȧ \alst{s}alar stęini &
\alst{Au}r-vanga sjǫt \hld\ til \alst{Jǫ}ru-valla.\eva

\bvb It is time to tally the dwarfs in Dwollen’s troop \\
back to Loffer for the race of men; \\
they who sought, from the stone of the hall, \\
the seat of the \inx[P]{Earwongs} unto the \inx[P]{Erwolds}.\evb\evg


\bvg\bva[15]\mssnote{\Regius~1r/32, \Hauksbok~20r/24}%
\Ballnote{The dwarf-names in this st.~are split up acros several sections in \Gylfaginning~14.
In \Gylfaginning~14:17, variants of the names in ll.~1--2 are cited, followed by four names not found in \Voluspa\ and then by the names in st.~13/2.  That name-list is introduced by the words \emph{En þessir eru ok dvergar ok búa í steinum, en inir fyrri í moldu:} ‘But these ones are also dwarfs and dwell in stones, but the former ones in the earth.’
The names of ll.~3--5 are cited in \Gylfaginning~14:19 immediately following a paraphrase of st.~14, introduced by the words \emph{þessi eru nǫfn þeira} ‘these are their names’.  Although these name-lists cannot be considered citations of this stanza proper, the variants of the names have been included (TODO!) in the critical edition.}%
Þar vas \alst{D}raupnir \hld\ ok \alst{D}olg·þrasir, &
\alst{H}ár, \alst{H}aug·spori, \hld\ \alst{H}lé·vangr, Glói, &
\alst{Sk}irfir, Virfir, \hld\ \alst{Sk}á·fiðr, Ái, &
\alst{A}lfr ok \alst{Y}ngvi, \hld\ \alst{Ęi}kin·skjaldi, &
\alst{F}jalarr ok \alst{F}rosti, \hld\ \alst{F}innr ok Ginnarr; &
\edtrans{Þat mun \edtext{\alst{ę́}}{\Afootnote{so \Hauksbok \emph{;} om. \Regius}} \alst{u}ppi, \hld\ meðan \alst{ǫ}ld lifir}{That will ever be remembered while the age lives}{\Bfootnote{Both verse-halves are formulaic, though their use in a list of dwarf-names is unparalleled.
For the first cf.~\HervararSaga\ 119 (in \Skp\ 7): \emph{þat mun ę́ uppi; \hld\ illr es dómr norna!} ‘That will always be remembered; evil is the judgement of the norns!’
For the second cf.~runic inscription U 323 Sälna (ca. 980--1015 \ce), edited under Runic Poetry from Sweden and Gotland.
What is essentially a Scaldic rewrite of the whole line is found in Þstf \emph{Stuttdr} 5/1--2 (\Skp\ 2): \emph{Ęy mun uppi \hld\ Ęndils, meðan stęndr // sól-borgar salr, \hld\ svǫr-gǿðis fǫr} ‘Ever will be remembered---while the hall of the sun’s stronghold \ken{sky/heaven > earth} stands---the journey of the fattener of Andle’s bird \ken{raven/eagle > warrior}.’}}, &
\alst{l}ang-niðja-tal \hld\ \edtext{\alst{L}ofars hafat}{\Bfootnote{This line would be unmetrical in the rest of the poem since the first lift is on a short syllable \parencite[373]{Males2025Hlod}.  This strongly suggests that the Dwarf-tally section is a later insert.}}.\eva

\bvb There was Dreepner and Dollowthrasher, \\
High, Highspurer, Leewong, Glower, \\
Sherver, Werver, Showfind, Great-grandfather, \\
Elf and Ing, Oakenshield, \\
Feller and Frost, Finn and Ginner. \\
That will ever be remembered while the age lives, \\
the recited tally of the ancestors of Loffer.\evb\evg

\subsection{Stanzas from \emph{Hauksbók}}

\Hauksbok\ has a few substantial insertions and variants not found in \Regius.  The following four stanzas, H31, H41, H48, H58, are the only instances longer than a single line (but cf.~st.~7/2 n.~above).

The style and content of H41, H48, and especially H58, comes off as rather different from the rest of the poem; they are therefore unlikely to be authorial.  H31 may, however, be genuine.

\bvg\bva[H31]\mssnote{\Hauksbok~20v/12}%
\Ballnote{These two lines replace the present 34/1--2 in \Hauksbok. ---

Lock was bound by Weden’s son Wonnel, who thus avenged Balder by punishing both of those responsible for his death: he killed Hath (who held the weapon) and bound Lock (who guided Hath’s hand).
Cf.~\Gylfaginning~50:20--22: \emph{Þá vǫ́ru teknir synir Loka, Váli ok Nari eða Narfi.  Brugðu Ę́sir Vála í vargs líki ok reif hann í sundr Narfa, bróður sinn.  Þá tóku ę́sir þarma hans ok bundu Loka með yfir þá þrjá egg-steina.  Stendr einn undir herðum, annarr undir lendum, þriði undir knés-fótum, ok urðu þau bǫnd at járni.} ‘Then were taken Lock’s sons, Wonnel and Nare or Narve.  The Eese turned Wonnel into the likeness of a wolf and he tore apart Narve, his brother.  Then the Eese took his intestines and with them they bound Lock over three sharp stones.  One stands under his shoulder-blades, the other under his lower back, the third under his houghs; and those bonds turned into iron.’

It is likely that Snorre’s version of \Voluspa\ contained this st.~and also had the oblique sg.~form \emph{*Vála} in l.~1, which led him to misinterpret Wonnel as the name of one of Lock’s sons, rather than the name of Weden’s son Wonnel, the avenger of Hath.  No other source attests Wonnel as the name of a son of Lock.  Cf.~the version of the binding of Lock found in \textlink{Lokasenna}[P8].}%
Þȧ kná \edtrans{\alst{V}ȧl\emph{i}}{Wonnel}{\Afootnote{\emph{Vála} \Hauksbok}} \hld\ \alst{v}íg-bǫnd snúa; &
\alst{h}ęldr vǫ́ru \alst{h}arð-gǫr \hld\ \alst{h}ǫpt ór þǫrmum.\eva

\bvb Then did \inx[P]{Wonnel} the war-bonds twist: \\
the most sturdy fetters were made from entrails.\evb\evg


\bvg\bva[H41]\mssnote{\Hauksbok~20v/28}%
\Ballnote{St.~45 is followed by these lines in \Hauksbok, in effect forming another four-line stanza with 45/5--6.  In context, the lines seem to be referring to an ettin swallowing the World Tree.}%
\alst{H}rę́ða⸗sk allir \hld\ ȧ \alst{h}ęl-vegum &
áðr \alst{S}urtar þann \hld\ \alst{s}efi of glęypir.\eva

\bvb All are frightened on the Hell-ways \\
before Surt’s kinsman devours it.\evb\evg


\bvg\bva[H48]\mssnote{\Hauksbok~20v/39}%
\Aallnote{The last part of the stanza is almost completely illegible in the ms.  I have relied on the transliteration of \textcite[13,44\psqq]{JonHelgason1971}.}%
\Ballnote{This stanza appears between 52 and 53 in \Hauksbok.  It is not found in the sequence of stanzas cited in \Gylfaginning~51:50, but some details in the preceding paraphrase are not entirely dissimilar.  \Gylfaginning~51:18--20: \emph{En Fenris-ulfr ferr með gapandi munn, ok er inn neðri kjǫptr við jǫrðu, en in øfri við himin.  Gapa myndi hann meira, ef rúm vę́ri til.  Eldar brenna ór augum hans ok nǫsum.  Mið·garðs-ormr blę́ss svá eitri’nu, at hann dreifir lopt ǫll ok lǫg, ok er hann all-ógur⸗ligr, ok er hann á aðra hlið ulfi’num.} ‘But the Fenrerswolf fares with gaping mouth; and the lower jaw is on the earth, but the upper is in the heaven; he would gape yet broader if there were room for it.  Fires burn from his eyes and nostrils.  The Middenyardswyrm blows the venom so far that it covers the whole air and sea, and he is most terrible, and he is by the side of the Wolf.’}%
Gïnn \alst{l}opt yfir \hld\ \alst{l}indi jarðar, &
gapa \alst{ý}gs kjaptar \hld\ \alst{o}rms ï hę́ðum; &
mun \alst{Ó}ðins son \hld\ \edtrans{\emph{\alst{ęi}tr}i}{venom}{\Afootnote{emend.; \emph{ormi} ‘Wyrm’ \Hauksbok.}\Bfootnote{Cf.~\Gylfaginning~51:41--42, cited in the notes to st.~53 above.}} mǿta &
\alst{v}args at \edtext{da\emph{uða}}{\Afootnote{\emph{‘da...’} \Hauksbok}} \hld\ \alst{V}íð·ars niðja.\eva

\bvb Over the air yawns the Girdle of the Earth \ken*{= Middenyardswyrm}; \\
the jaws of the fierce Wyrm gape in the heights. \\
Weden’s son \ken*{= Thunder} will meet the venom \\
of the Warg, after the deaths of Wider’s kinsmen \ken*{= the Eese}.\evb\evg


\bvg\bva[H58]\mssnote{\Hauksbok~21r/14}%
\Ballnote{This half-stanza appears between sts.~61 and 62 in \Hauksbok.  It is generally held to be a late Christian insert, since it betrays Christian theological sentiments.  An omnipotent God coming down from above to the Great Judgment (\emph{ręgin-dȯmr})---none of these elements are found in the pre-Christian religion.  Cf.~\textlink{Hyndluljod}[43], which is likely drawing on the present stanza.}%
Þȧ kømr hinn \alst{r}íki \hld\ at \alst{r}ęgin-dȯmi &
\alst{ǫ}flugr \alst{o}fan \hld\ sá’s \alst{ǫ}llu rę́ðr.\eva

\bvb Then comes the powerful one to the Great Judgment, \\
strong down from above, he who rules everything.\evb\evg

\sectionline
